% Fonction logarithme népérien — Mathématiques Tle TI — Cours signé OnBuch+
% Programme officiel MINESEC. Compiler avec : tectonic M06-fonction-logarithme-neperien.tex
\newcommand{\DOCMATIERE}{Mathématiques}
\newcommand{\DOCNIVEAU}{Tle TI}
\newcommand{\DOCMODULE}{Relations et opérations fondamentales dans l'ensemble des nombres réels}
\newcommand{\DOCLECON}{06}
\newcommand{\DOCTITRE}{Fonction logarithme népérien}
% OnBuch+ — préambule « cours signés » ENRICHI (compilé avec tectonic / XeLaTeX)
% Système visuel « Pop » d'OnBuch+ : fond crème (jamais blanc), bordures encre
% épaisses, ombres « dures » décalées, titres Archivo Black en capitales.
% Version MATHÉMATIQUES : graphes (pgfplots/tikz), boîtes pédagogiques
% (définition, propriété, méthode, attention, le savais-tu, exercice résolu,
% exercice, TP, bilan), page de garde et signature OnBuch+ ancrée en bas.
%
% Macros attendues dans main.tex AVANT \input{preamble} :
%   \DOCMATIERE  \DOCNIVEAU  \DOCMODULE  \DOCLECON (numéro)  \DOCTITRE
\documentclass[11pt]{article}

\usepackage{fontspec}
\usepackage[french]{babel}
\usepackage[a4paper,margin=2cm,top=3.4cm,bottom=2.8cm,headheight=1.4cm,footskip=1.3cm]{geometry}
\usepackage{amsmath,amssymb,amsfonts}
\usepackage{mathtools} % \xmapsto, \coloneqq... souvent utilisés par les modèles
\usepackage{cancel} % simplifications barrées
% \abs{z} (module, valeur absolue) et \norm{u} (norme), utilisés par les modèles
\providecommand{\abs}{}\let\abs\relax\DeclarePairedDelimiter{\abs}{\lvert}{\rvert}
\providecommand{\norm}{}\let\norm\relax\DeclarePairedDelimiter{\norm}{\lVert}{\rVert}
\usepackage[table]{xcolor} % [table] : fournit aussi \rowcolors (alternance de lignes)
\usepackage{pagecolor}
\usepackage{tikz}
\usetikzlibrary{arrows.meta,positioning,calc,shapes.geometric,shapes.misc,shapes.arrows,decorations.pathmorphing,decorations.markings,patterns,shadows,mindmap,trees,fit,backgrounds,matrix,intersections}
\usepackage{pgfplots}
\usepackage{pgf-pie} % diagrammes circulaires \pie (statistiques)
\pgfplotsset{compat=1.18}
\usepgfplotslibrary{fillbetween,groupplots} % aires entre courbes (calcul intégral)
\usepackage{enumitem}
\usepackage{fancyhdr}
% [most] charge toutes les bibliothèques tcolorbox (listings, minted, raster,
% documentation...) et gonfle inutilement la mémoire TeX sur un document long
% (~300 boîtes) : on ne charge que ce qui est réellement utilisé.
\usepackage[skins,breakable]{tcolorbox}
\usepackage{graphicx}
\usepackage{colortbl}
\usepackage{booktabs}
\usepackage{tabularx}
\usepackage{diagbox} % case d'en-tête diagonale des tableaux à double entrée
% Colonnes extensibles centrée / gauche / droite (C, L, R), souvent utilisées par les modèles
\newcolumntype{C}{>{\centering\arraybackslash}X}
\newcolumntype{L}{>{\raggedright\arraybackslash}X}
\newcolumntype{R}{>{\raggedleft\arraybackslash}X}
\usepackage{array}
\usepackage{multicol}
\usepackage{siunitx}
% parse-numbers=false : accepte aussi \SI{2,5 \times 10^{-3}}{...}
\sisetup{locale=FR,output-decimal-marker={,},group-minimum-digits=4,parse-numbers=false}
\DeclareSIUnit\ohm{\ensuremath{\symup{\Omega}}} % U+2126 absent de la police
\usepackage{qrcode}
% \degree : commande fréquemment utilisée par les modèles pour un angle ou une
% température en dehors de \SI{}{} (non fournie par les paquets chargés).
\providecommand{\degree}{^{\circ}}
% \qed (fin de démonstration), utilisé par les modèles sans amsthm
\providecommand{\qed}{\hfill$\square$}
% \barre : double barre des valeurs interdites dans un tableau de variations (tkz-tab)
\providecommand{\barre}{\ensuremath{\|}}
% environnement proof (amsthm non chargé) utilisé par les modèles
\ifdefined\proof\else\newenvironment{proof}[1][Démonstration]{\par\noindent\textit{#1.}\ }{\qed\par}\fi
\usepackage{lastpage}
\usepackage{hyperref}
\hypersetup{hidelinks}

% ============================================================
% Palette « Pop » OnBuch+ — mêmes hex que la classe P (plus_ui.dart)
% ============================================================
\definecolor{poporange}{HTML}{F59321}
\definecolor{popdark}{HTML}{E07A0C}
\definecolor{popcream}{HTML}{FFF6E8}
\definecolor{popink}{HTML}{1C1714}
\definecolor{poppurple}{HTML}{7A5AE0}
\definecolor{popgreen}{HTML}{2E9E6A}
\definecolor{poppink}{HTML}{E0457B}
\definecolor{popblue}{HTML}{2F7DE1}
\definecolor{popgold}{HTML}{C98A12}
\definecolor{popmuted}{HTML}{8A7F76}
\definecolor{popline}{HTML}{EADFD2}
\definecolor{popgray}{HTML}{8A7F76}
\colorlet{popgrey}{popgray}
% Alias tolérants (le modèle nomme parfois les couleurs différemment)
\colorlet{popwhite}{white}
\colorlet{popblack}{popink}
\colorlet{popred}{poppink}
\colorlet{popyellow}{popgold}
% Teintes claires pour fonds de tableaux / zones
\colorlet{popblueL}{popblue!10!white}
\colorlet{poporangeL}{poporange!14!white}
\colorlet{popgreenL}{popgreen!12!white}
\colorlet{poppurpleL}{poppurple!12!white}
\colorlet{poppinkL}{poppink!12!white}
\colorlet{popgoldL}{popgold!14!white}

\pagecolor{popcream}

% ============================================================
% Typographie
% ============================================================
\defaultfontfeatures{Ligatures=TeX}
\setmainfont{PlusJakartaSans-Medium.ttf}[Path=../../../fonts/,BoldFont=PlusJakartaSans-Bold.ttf]
% Maths en sans-serif (Fira Math) avec chiffres et lettres droites de Plus
% Jakarta, pour que les formules se fondent dans le texte.
\usepackage[math-style=ISO,bold-style=ISO]{unicode-math}
\setmathfont{FiraMath-Regular.otf}[Path=../../../fonts/]
\setmathfont{PlusJakartaSans-Medium.ttf}[Path=../../../fonts/,range={up/num,up/{Latin,latin},\mathup}]
\usepackage{icomma} % virgule décimale sans espace en mode math
% Caractères mathématiques Unicode absents des polices de texte (Plus Jakarta,
% Archivo Black) : rendus par leur équivalent mathématique, en texte comme en
% formule — sinon ils s'affichent en carré vide (ex. « Arithmétique dans ℤ »).
\usepackage{newunicodechar}
\newunicodechar{ℝ}{\ensuremath{\mathbb{R}}}
\newunicodechar{ℤ}{\ensuremath{\mathbb{Z}}}
\newunicodechar{ℂ}{\ensuremath{\mathbb{C}}}
\newunicodechar{ℕ}{\ensuremath{\mathbb{N}}}
\newunicodechar{ℚ}{\ensuremath{\mathbb{Q}}}
\newunicodechar{𝔻}{\ensuremath{\mathbb{D}}}
\newunicodechar{α}{\ensuremath{\alpha}}
\newunicodechar{β}{\ensuremath{\beta}}
\newunicodechar{γ}{\ensuremath{\gamma}}
\newunicodechar{δ}{\ensuremath{\delta}}
\newunicodechar{ε}{\ensuremath{\varepsilon}}
\newunicodechar{θ}{\ensuremath{\theta}}
\newunicodechar{λ}{\ensuremath{\lambda}}
\newunicodechar{μ}{\ensuremath{\mu}}
\newunicodechar{σ}{\ensuremath{\sigma}}
\newunicodechar{τ}{\ensuremath{\tau}}
\newunicodechar{φ}{\ensuremath{\varphi}}
\newunicodechar{ψ}{\ensuremath{\psi}}
\newunicodechar{ω}{\ensuremath{\omega}}
\newunicodechar{Σ}{\ensuremath{\Sigma}}
% majuscules produites par \MakeUppercase dans les titres (α-aminés -> Α)
\newunicodechar{Α}{\ensuremath{\alpha}}
\newunicodechar{Β}{\ensuremath{\beta}}
\newunicodechar{Γ}{\ensuremath{\Gamma}}
\newunicodechar{Δ}{\ensuremath{\Delta}}
\newunicodechar{Ε}{\ensuremath{\varepsilon}}
\newunicodechar{Θ}{\ensuremath{\Theta}}
\newunicodechar{Λ}{\ensuremath{\Lambda}}
\newunicodechar{Μ}{\ensuremath{\mu}}
\newunicodechar{Τ}{\ensuremath{\tau}}
\newunicodechar{Φ}{\ensuremath{\Phi}}
\newunicodechar{Ψ}{\ensuremath{\Psi}}
\newunicodechar{Ω}{\ensuremath{\Omega}}
\newunicodechar{↦}{\ensuremath{\mapsto}}
\newunicodechar{∈}{\ensuremath{\in}}
\newunicodechar{∉}{\ensuremath{\notin}}
\newunicodechar{∧}{\ensuremath{\wedge}}
\newunicodechar{⇔}{\ensuremath{\Leftrightarrow}}
\newunicodechar{⟺}{\ensuremath{\Longleftrightarrow}}
\newunicodechar{⇒}{\ensuremath{\Rightarrow}}
\newunicodechar{⟹}{\ensuremath{\Longrightarrow}}
\newunicodechar{∘}{\ensuremath{\circ}}
\newunicodechar{∪}{\ensuremath{\cup}}
\newunicodechar{∩}{\ensuremath{\cap}}
\newunicodechar{≡}{\ensuremath{\equiv}}
\newunicodechar{⊂}{\ensuremath{\subset}}
\newunicodechar{⊥}{\ensuremath{\perp}}
\newunicodechar{∃}{\ensuremath{\exists}}
\newunicodechar{∀}{\ensuremath{\forall}}
\newunicodechar{∛}{\ensuremath{\sqrt[3]{\,}}}
\newunicodechar{∜}{\ensuremath{\sqrt[4]{\,}}}
\newunicodechar{⃗}{\ensuremath{{}^{\to}}}
\newunicodechar{ⁿ}{\ifmmode^{n}\else\textsuperscript{\ensuremath{n}}\fi}
\newunicodechar{ˣ}{\ifmmode^{x}\else\textsuperscript{\ensuremath{x}}\fi}
\newunicodechar{ᵇ}{\ifmmode^{b}\else\textsuperscript{\ensuremath{b}}\fi}
\newunicodechar{ʳ}{\ifmmode^{r}\else\textsuperscript{\ensuremath{r}}\fi}
\newunicodechar{ᵘ}{\ifmmode^{u}\else\textsuperscript{\ensuremath{u}}\fi}
\newunicodechar{ʸ}{\ifmmode^{y}\else\textsuperscript{\ensuremath{y}}\fi}
\newunicodechar{ᵃ}{\ifmmode^{a}\else\textsuperscript{\ensuremath{a}}\fi}
\newunicodechar{ᵉ}{\ifmmode^{e}\else\textsuperscript{\ensuremath{e}}\fi}
\newunicodechar{ᵏ}{\ifmmode^{k}\else\textsuperscript{\ensuremath{k}}\fi}
\newunicodechar{ᵖ}{\ifmmode^{p}\else\textsuperscript{\ensuremath{p}}\fi}
\newunicodechar{ᴺ}{\ifmmode^{N}\else\textsuperscript{\ensuremath{N}}\fi}
\newunicodechar{ᶜ}{\ifmmode^{c}\else\textsuperscript{\ensuremath{c}}\fi}
\newunicodechar{ʰ}{\ifmmode^{h}\else\textsuperscript{\ensuremath{h}}\fi}
\newunicodechar{ᵠ}{\ifmmode^{q}\else\textsuperscript{\ensuremath{q}}\fi}
\newunicodechar{ₙ}{\ifmmode_{n}\else\textsubscript{\ensuremath{n}}\fi}
\newunicodechar{ₐ}{\ifmmode_{a}\else\textsubscript{\ensuremath{a}}\fi}
\newunicodechar{ᵢ}{\ifmmode_{i}\else\textsubscript{\ensuremath{i}}\fi}
\newunicodechar{ᵣ}{\ifmmode_{r}\else\textsubscript{\ensuremath{r}}\fi}
\newunicodechar{ₖ}{\ifmmode_{k}\else\textsubscript{\ensuremath{k}}\fi}
\newunicodechar{ᵦ}{\ifmmode_{\beta}\else\textsubscript{\ensuremath{\beta}}\fi}
\newunicodechar{ₓ}{\ifmmode_{x}\else\textsubscript{\ensuremath{x}}\fi}
\newfontfamily\popdisplay{ArchivoBlack-Regular.ttf}[Path=../../../fonts/]
\newfontfamily\poplogo{SpaceGrotesk-Bold.ttf}[Path=../../../fonts/]
\newfontfamily\popsemi{PlusJakartaSans-SemiBold.ttf}[Path=../../../fonts/]
\newfontfamily\popxbold{PlusJakartaSans-ExtraBold.ttf}[Path=../../../fonts/]
\newfontfamily\popxboldls{PlusJakartaSans-ExtraBold.ttf}[Path=../../../fonts/,LetterSpace=3.0]

\setlist[enumerate]{leftmargin=1.7em,itemsep=5pt,topsep=4pt}
\setlist[itemize]{leftmargin=1.7em,itemsep=4pt,topsep=4pt,label={\color{poporange}\textbullet}}
\setlength{\parindent}{0pt}
\setlength{\parskip}{5pt}
\renewcommand{\arraystretch}{1.25}

% pgfplots : style Pop par défaut
\pgfplotsset{
  every axis/.append style={axis line style={popink,line width=1pt},tick style={popink},
    grid style={popline,line width=0.5pt},label style={font=\small},tick label style={font=\footnotesize}},
  popcurve/.style={poporange,line width=1.8pt,no markers,smooth},
}
% Même style disponible dans un \draw TikZ simple (hors pgfplots)
\tikzset{popcurve/.style={poporange,line width=1.8pt}}

% ============================================================
% Pill / squiggle
% ============================================================
\newcommand{\poppill}[3]{%
  \tikz[baseline=-0.55ex]\node[draw=popink,line width=1.2pt,rounded corners=2.6pt,%
    fill=#2,inner xsep=6.5pt,inner ysep=3pt]{%
    \popxboldls\fontsize{7}{7}\selectfont\color{#3}\MakeUppercase{#1}};%
}
\newcommand{\popsquiggle}[1]{%
  \tikz[baseline=-0.4ex]{
    \draw[#1,line width=2.6pt,line cap=round]
      plot [smooth] coordinates {(0,0.09) (0.34,0.01) (0.68,0.21) (1.02,0.01) (1.36,0.10)};
  }%
}
% Mot-clé surligné (marqueur orange sous le texte)
\newcommand{\cle}[1]{{\popxbold\color{popdark}#1}}

% ============================================================
% En-tête / pied
% ============================================================
\pagestyle{fancy}
\fancyhf{}
\renewcommand{\headrulewidth}{0pt}
\renewcommand{\footrulewidth}{0pt}
\lhead{\raisebox{-0.15ex}{\poplogo\fontsize{13}{13}\selectfont\color{popink}OnBuch\color{poporange}+}}
\rhead{\poppill{\DOCMATIERE\ \textbullet\ \DOCNIVEAU\ \textbullet\ Leçon \DOCLECON}{white}{popink}}
\lfoot{\popxbold\fontsize{7.6}{7.6}\selectfont\color{popmuted}\MakeUppercase{Cours signé OnBuch+}\ \textbullet\ {\popsemi\DOCTITRE}}
\rfoot{\tikz[baseline=-0.6ex]\node[rounded corners=8.5pt,fill=popink,minimum height=17pt,minimum width=17pt,inner xsep=5pt,inner sep=0]{\popxbold\fontsize{8}{8}\selectfont\color{popcream}\thepage\,/\,\pageref*{LastPage}};}

% ============================================================
% Titres
% ============================================================
\newcommand{\chaptitre}[2]{%
  \begin{tcolorbox}[enhanced,colback=poporange,colframe=popink,boxrule=2.6pt,arc=11pt,
      drop shadow=popink,left=15pt,right=15pt,top=12pt,bottom=11pt,before skip=3pt,after skip=18pt]
    \poppill{#2}{popink}{popcream}\par\vspace{9pt}
    {\raggedright\hyphenpenalty=10000\exhyphenpenalty=10000\popdisplay\fontsize{22}{24}\selectfont\color{popcream}\MakeUppercase{#1}\par}
  \end{tcolorbox}
}
% Section numérotée : pastille orange avec le numéro + titre Archivo
\newcounter{popsec}
\newcommand{\coursec}[1]{%
  \stepcounter{popsec}\setcounter{popsub}{0}%
  \par\vspace{16pt}\noindent
  \begin{minipage}{\linewidth}
  \tikz[baseline=-0.7ex]\node[draw=popink,line width=1.6pt,fill=poporange,rounded corners=4pt,minimum size=22pt,inner sep=0]{\popdisplay\fontsize{12}{12}\selectfont\color{popcream}\thepopsec};%
  \hspace{8pt}{\popdisplay\fontsize{15}{17}\selectfont\color{popink}\MakeUppercase{#1}}\par
  \vspace{4pt}\noindent{\color{poporange}\rule{36pt}{3.6pt}}
  \end{minipage}\par\nopagebreak\vspace{8pt}
}
\newcounter{popsub}
\newcommand{\courssub}[1]{%
  \stepcounter{popsub}%
  \par\vspace{9pt}\noindent{\popxbold\fontsize{11.5}{13}\selectfont\color{popink}{\color{poporange}\thepopsec.\thepopsub}\hspace{6pt}#1}\par\nopagebreak\vspace{4pt}
}

% ============================================================
% Boîtes pédagogiques — même recette : fond blanc, bordure épaisse colorée,
% ombre dure encre, badge plein. Titre optionnel [#1].
% ============================================================
\tcbset{popbase/.style={breakable,enhanced,colback=white,boxrule=2.3pt,arc=9pt,
  shadow={3pt}{-3pt}{0pt}{popink},left=14pt,right=14pt,top=11pt,bottom=11pt,before skip=11pt,after skip=15pt,
  fontupper=\popsemi\fontsize{10.6}{14.5}\selectfont\color{popink},
  fontlower=\popsemi\fontsize{10.6}{14.5}\selectfont\color{popink}}}
% Coche dessinée (le glyphe ✓ n'existe pas dans les polices du document)
\newcommand{\popcheck}[1]{\tikz[baseline=-0.5ex]\draw[#1,line width=1.6pt,line cap=round,line join=round](-0.13,0.0)--(-0.04,-0.09)--(0.13,0.1);}
\newcommand{\popboxhead}[3]{\poppill{#1}{#2}{white}\ifx\relax#3\relax\else\hspace{6pt}{\popxbold\fontsize{10.6}{12}\selectfont\color{popink}#3}\fi\par\vspace{7pt}}

\newtcolorbox{aretenir}[1][]{popbase,colframe=popgreen,before upper={\popboxhead{À retenir}{popgreen}{#1}}}
\newtcolorbox{exemplebox}[1][]{popbase,colframe=poporange,before upper={\popboxhead{Exemple}{poporange}{#1}}}
\newtcolorbox{definition}[1][]{popbase,colframe=popblue,before upper={\popboxhead{Définition}{popblue}{#1}}}
\newtcolorbox{propriete}[1][]{popbase,colframe=poppurple,before upper={\popboxhead{Propriété}{poppurple}{#1}}}
\newtcolorbox{methode}[1][]{popbase,colframe=popgold,before upper={\popboxhead{Méthode}{popgold}{#1}}}
\newtcolorbox{attention}[1][]{popbase,colframe=poppink,before upper={\popboxhead{Attention}{poppink}{#1}}}
\newtcolorbox{experience}[1][]{popbase,colframe=popblue,colback=popblueL,before upper={\popboxhead{Expérience}{popblue}{#1}}}
% « Le savais-tu ? » : carte violette pleine (contexte, culture, Cameroun)
\newtcolorbox{savaistu}[1][]{popbase,colback=poppurple,colframe=popink,
  fontupper=\popsemi\fontsize{10.4}{14.2}\selectfont\color{white},
  before upper={\poppill{Le savais-tu\,?}{white}{poppurple}\ifx\relax#1\relax\else\hspace{6pt}{\popxbold\color{white}#1}\fi\par\vspace{7pt}}}
% Exercice résolu : énoncé en haut, \tcblower puis la solution sur fond crème
\newcounter{popexo}\newcounter{popexr}
\newtcolorbox{exoresolu}[1][]{popbase,colframe=poporange,colbacklower=popcream,
  segmentation style={popink,line width=1.2pt,dashed},
  before upper={\stepcounter{popexr}\popboxhead{Exercice résolu \thepopexr}{poporange}{#1}},
  before lower={\poppill{Solution}{popink}{popcream}\par\vspace{6pt}}}
% Exercice (non résolu en place) — la correction est dans la partie Corrigés
\newtcolorbox{exercice}[1][]{popbase,colframe=popink,
  before upper={\stepcounter{popexo}\popboxhead{Exercice \thepopexo}{popink}{#1}}}
% Corrigé
\newtcolorbox{corrige}[1][]{popbase,colframe=popgreen,colback=popgreenL,
  before upper={\popboxhead{Corrigé}{popgreen}{#1}}}
% Objectifs / prérequis / bilan
\newtcolorbox{objectifs}{popbase,colframe=popink,colback=poporangeL,
  before upper={\popboxhead{Objectifs de la leçon}{popink}{}}}
\newtcolorbox{prerequis}{popbase,colframe=popink,colback=popblueL,
  before upper={\popboxhead{Prérequis}{popblue}{}}}
\newtcolorbox{bilan}{popbase,colframe=popink,colback=white,boxrule=2.8pt,
  before upper={\popboxhead{Fiche bilan}{poporange}{L'essentiel à connaître pour l'examen}}}
% Figure encadrée avec légende
\newtcolorbox{popfigure}[1][]{enhanced,breakable=false,colback=white,colframe=popline,boxrule=1.4pt,arc=8pt,
  left=10pt,right=10pt,top=10pt,bottom=8pt,before skip=10pt,after skip=12pt,halign=center,
  lower separated=false,halign lower=center,
  fontlower=\popsemi\fontsize{9}{11.5}\selectfont\color{popmuted},
  #1}
\newcommand{\legende}[1]{\par\vspace{6pt}{\popsemi\fontsize{9}{11.5}\selectfont\color{popmuted}#1}}

% ============================================================
% Signature OnBuch+
% ============================================================
\newcommand{\poplogoinline}[1]{{\poplogo\fontsize{#1}{#1}\selectfont\color{popink}OnBuch\color{poporange}+}}
\newcommand{\signatureonbuch}{%
  \begin{tcolorbox}[enhanced,colback=popink,colframe=popink,boxrule=2.2pt,arc=10pt,
      drop shadow=poporange,left=14pt,right=14pt,top=10pt,bottom=10pt,before skip=0pt,after skip=0pt]
    \tikz[baseline=-0.7ex]\node[circle,fill=popgreen,draw=popcream,line width=1.3pt,minimum size=22pt,inner sep=0]{\popcheck{white}};%
    \hspace{9pt}\begin{minipage}[c]{0.62\linewidth}
      {\popxbold\fontsize{12}{13}\selectfont\color{popcream}Signé\ \textbullet\ {\poplogo OnBuch\color{poporange}+}}\par\vspace{2pt}
      {\raggedright\popsemi\fontsize{8}{10.5}\selectfont\color{popline}\MakeUppercase{Équipe pédagogique OnBuch+}\\\MakeUppercase{Conforme au programme officiel MINESEC}\par}
    \end{minipage}\hfill
    {\poplogo\fontsize{22}{22}\selectfont\color{popcream}OnBuch\color{poporange}+}
  \end{tcolorbox}%
}

% ============================================================
% Page de garde
% #1 titre · #2 matière · #3 niveau · #4 module · #5 n° leçon · #6 durée ·
% #7 description / objectifs (texte ou itemize)
% ============================================================
\newcommand{\pagegarde}[7]{%
  \thispagestyle{empty}
  \newgeometry{a4paper,margin=1.7cm,top=1.6cm,bottom=1.4cm}%
  \begin{tikzpicture}[remember picture,overlay]
    \fill[poporange!18] ([xshift=-3.2cm,yshift=-3.6cm]current page.north east) circle (4.6cm);
    \fill[poppurple!14] ([xshift=2.4cm,yshift=7.6cm]current page.south west) circle (3.8cm);
  \end{tikzpicture}%
  {\poplogo\fontsize{30}{30}\selectfont\color{popink}OnBuch\color{poporange}+}\hfill
  \raisebox{0.6ex}{\poppill{Cours signé \textbullet\ Premium}{popink}{popcream}}\par
  \vspace{1pt}\popsquiggle{poppurple}\par
  \vspace{8pt}
  \begin{tcolorbox}[enhanced,colback=poporange,colframe=popink,boxrule=2.8pt,arc=13pt,
      drop shadow=popink,left=18pt,right=18pt,top=12pt,bottom=13pt,after skip=10pt]
    \poppill{#2 \textbullet\ #3}{popink}{popcream}\hspace{5pt}\poppill{#4}{white}{popink}\par\vspace{8pt}
    {\popdisplay\fontsize{14}{14}\selectfont\color{popink}LEÇON #5}\par\vspace{4pt}
    {\raggedright\hyphenpenalty=10000\exhyphenpenalty=10000\popdisplay\fontsize{25}{27}\selectfont\color{popcream}\MakeUppercase{#1}\par}
    \vspace{8pt}
    {\popxbold\fontsize{9.5}{9.5}\selectfont\color{popink}\MakeUppercase{Durée conseillée : #6}}
  \end{tcolorbox}
  \begin{tcolorbox}[enhanced,colback=white,colframe=popink,boxrule=2.2pt,arc=10pt,
      drop shadow=popink,left=16pt,right=16pt,top=10pt,bottom=10pt,after skip=8pt]
    \poppill{Dans cette leçon}{poppurple}{white}\par\vspace{5pt}
    \popsemi\fontsize{9.3}{12.6}\selectfont\color{popink}#7
  \end{tcolorbox}
  \noindent\begin{minipage}[t]{0.487\linewidth}
    \begin{tcolorbox}[enhanced,colback=popgreen,colframe=popink,boxrule=2.2pt,arc=10pt,
        drop shadow=popink,left=11pt,right=11pt,top=10pt,bottom=10pt,height=3.1cm,valign=center]
      \tcbox[colback=white,colframe=popink,boxrule=1.3pt,arc=5pt,boxsep=0pt,nobeforeafter,
        left=3pt,right=3pt,top=3pt,bottom=3pt,valign=center]{\qrcode[height=1.9cm]{https://play.google.com/store/apps/details?id=com.luvvix.onbuch&hl=fr}}%
      \hspace{8pt}\begin{minipage}[c]{0.5\linewidth}
        {\popdisplay\fontsize{11}{12.5}\selectfont\color{white}\MakeUppercase{Télécharge OnBuch}}\par\vspace{4pt}
        {\raggedright\popsemi\fontsize{8.4}{11}\selectfont\color{white}Gratuit \textbullet\ Google Play\\Tuteur IA, annales, concours, résultats\par}
      \end{minipage}
    \end{tcolorbox}
  \end{minipage}\hfill
  \begin{minipage}[t]{0.487\linewidth}
    \begin{tcolorbox}[enhanced,colback=poppurple,colframe=popink,boxrule=2.2pt,arc=10pt,
        drop shadow=popink,left=11pt,right=11pt,top=10pt,bottom=10pt,height=3.1cm,valign=center]
      \tikz[baseline=-0.7ex]\node[circle,fill=white,draw=popink,line width=1.3pt,minimum size=1.6cm,inner sep=0]{\popdisplay\fontsize{24}{24}\selectfont\color{poppurple}+};%
      \hspace{8pt}\begin{minipage}[c]{0.6\linewidth}
        {\popdisplay\fontsize{11}{12.5}\selectfont\color{white}\MakeUppercase{Passe à OnBuch+}}\par\vspace{4pt}
        {\raggedright\popsemi\fontsize{8.4}{11}\selectfont\color{white}Tous les cours signés, Tuteur IA illimité, fiches PDF — onglet OnBuch+ de l'app\par}
      \end{minipage}
    \end{tcolorbox}
  \end{minipage}
  \vspace{8pt}
  \popcontenu
  \vfill
  \signatureonbuch
  \restoregeometry
  \newpage
}

% Rangée « Ce cours contient » (page de garde)
\newcommand{\poptile}[3]{% #1 couleur #2 gros texte #3 légende
  \begin{tcolorbox}[enhanced,colback=#1,colframe=popink,boxrule=2pt,arc=9pt,shadow={2.5pt}{-2.5pt}{0pt}{popink},
      left=6pt,right=6pt,top=9pt,bottom=9pt,halign=center,valign=center,height=1.75cm,width=\linewidth,nobeforeafter]
    {\popdisplay\fontsize{12.5}{13}\selectfont\color{white}\MakeUppercase{#2}}\par\vspace{3pt}
    {\centering\popsemi\fontsize{7.8}{9.5}\selectfont\color{white}#3\par}
  \end{tcolorbox}}
\newcommand{\popcontenu}{%
  \par\noindent{\popxboldls\fontsize{7.5}{7.5}\selectfont\color{popmuted}CE COURS SIGNÉ CONTIENT}\par\vspace{7pt}
  \noindent\begin{minipage}[t]{0.235\linewidth}\poptile{popblue}{Cours}{complet, illustré, pas à pas}\end{minipage}\hfill
  \begin{minipage}[t]{0.235\linewidth}\poptile{poporange}{Méthodes}{et exercices résolus}\end{minipage}\hfill
  \begin{minipage}[t]{0.235\linewidth}\poptile{poppink}{Exercices}{type examen + corrigés}\end{minipage}\hfill
  \begin{minipage}[t]{0.235\linewidth}\poptile{popgreen}{Fiche bilan}{l'essentiel à retenir}\end{minipage}\par}

% Sommaire
\newtcolorbox{sommaire}{enhanced,colback=white,colframe=popink,boxrule=2.2pt,arc=10pt,
  drop shadow=popink,left=18pt,right=18pt,top=13pt,bottom=13pt,before skip=6pt,after skip=20pt,
  before upper={\poppill{Sommaire}{poppurple}{white}\par\vspace{9pt}\popsemi\fontsize{10.6}{14.5}\selectfont\color{popink}}}

% Signature de fin de cours — poussée tout en bas de la dernière page
\newcommand{\findecours}{%
  \par\vfill\nopagebreak
  \begin{center}\popsquiggle{poporange}\end{center}\vspace{4pt}
  \signatureonbuch
}
\AtBeginDocument{\def\llbracket{\mathopen{[\mkern-3.5mu[}}\def\rrbracket{\mathclose{]\mkern-3.5mu]}}} % intervalles d'entiers (glyphe ⟦ absent de la police)
% Glyphes absents de Fira Math : redéfinis à partir de glyphes disponibles
\AtBeginDocument{%
  \def\setminus{\mathbin{\backslash}}\let\smallsetminus\setminus
  \def\vdots{\mathord{\vbox{\baselineskip3.2pt\lineskiplimit0pt\kern4pt\hbox{.}\hbox{.}\hbox{.}}}}%
  \def\ddots{\mathinner{\mkern1mu\raise6.5pt\hbox{.}\mkern2mu\raise3.6pt\hbox{.}\mkern2mu\raise0.7pt\hbox{.}\mkern1mu}}%
  \def\triangle{\mathord{\scalebox{0.9}{$\symup{\Delta}$}}}%
  \def\nabla{\mathord{\raisebox{\depth}{\scalebox{1}[-1]{$\symup{\Delta}$}}}}%
  \def\checkmark{\popcheck{popdark}}%
  \def\star{\mathbin{\ast}}\def\bigstar{\mathord{\ast}}%
  \def\diamond{\mathbin{\scalebox{0.6}{$\Diamond$}}}%
}

%% FIGFIX-BEGIN v5 — correctifs globaux des figures (docs/onbuch-generation/tools/figfix)
\usepackage{adjustbox}\usepackage{etoolbox}\tcbuselibrary{raster}
% toute figure TikZ/pgfplots plus large que la ligne est réduite à la largeur disponible
\BeforeBeginEnvironment{tikzpicture}{\begin{adjustbox}{max width=\linewidth}}
\AfterEndEnvironment{tikzpicture}{\end{adjustbox}}
% légendes pgfplots lisibles par-dessus les courbes
\pgfplotsset{every axis/.append style={legend style={fill=white,fill opacity=0.92,text opacity=1,draw=black!15,inner sep=3pt}}}
% étiquettes d'arêtes (syntaxe edge["..."]) avec fond blanc
\tikzset{every edge quotes/.append style={fill=white,inner sep=1.2pt}}
%% FIGFIX-END
\begin{document}

\pagegarde{Fonction logarithme népérien}{Mathématiques}{Tle TI}{Relations et opérations fondamentales dans l'ensemble des nombres réels}{06}{9 h}{Cette leçon introduit la fonction logarithme népérien comme primitive de x ↦ 1/x s'annulant en 1, établit ses propriétés algébriques fondamentales, ses limites, sa dérivée et sa courbe représentative. Elle fournit les outils pour résoudre équations, inéquations et systèmes, calculer des primitives de la forme u′/u et modéliser des situations concrètes de croissance, de désintégration ou d'échelles logarithmiques.
\vspace{4pt}
\begin{multicols}{2}\small
\begin{itemize}
\item À la fin de cette leçon, je sais définir la fonction ln et préciser son ensemble de définition
\item À la fin de cette leçon, je sais utiliser les propriétés algébriques de ln pour simplifier ou développer une expression
\item À la fin de cette leçon, je sais déterminer les limites de ln aux bornes de son domaine et celles de fonctions composées avec ln
\item À la fin de cette leçon, je sais dériver une fonction de la forme ln u et étudier ses variations
\item À la fin de cette leçon, je sais dresser le tableau de variation complet de ln et tracer sa courbe représentative avec ses éléments remarquables
\item À la fin de cette leçon, je sais déterminer une primitive de la forme u′/u
\end{itemize}\end{multicols}}

\begin{sommaire}
\begin{multicols}{2}
\begin{enumerate}
\item Définition et premières propriétés de la fonction ln
\item Propriétés algébriques de ln
\item Limites de la fonction ln
\item Dérivation, variations et courbe représentative
\item Primitives de la forme u′/u
\item Équations, inéquations, systèmes et applications
\item Activité d'intégration
\item Méthodes et exercices résolus
\item Exercices
\item Corrigés des exercices
\item Fiche bilan
\end{enumerate}
\end{multicols}
\end{sommaire}

\begin{objectifs}
\begin{itemize}
\item À la fin de cette leçon, je sais définir la fonction ln et préciser son ensemble de définition
\item À la fin de cette leçon, je sais utiliser les propriétés algébriques de ln pour simplifier ou développer une expression
\item À la fin de cette leçon, je sais déterminer les limites de ln aux bornes de son domaine et celles de fonctions composées avec ln
\item À la fin de cette leçon, je sais dériver une fonction de la forme ln u et étudier ses variations
\item À la fin de cette leçon, je sais dresser le tableau de variation complet de ln et tracer sa courbe représentative avec ses éléments remarquables
\item À la fin de cette leçon, je sais déterminer une primitive de la forme u′/u
\item À la fin de cette leçon, je sais résoudre des équations, inéquations et systèmes faisant intervenir ln x
\item À la fin de cette leçon, je sais utiliser ln pour modéliser et résoudre des problèmes concrets (croissance, pH, échelle logarithmique)
\end{itemize}
\end{objectifs}

\begin{prerequis}
\begin{itemize}
\item Notion de primitive d'une fonction continue sur un intervalle (début de Terminale)
\item Dérivation des fonctions usuelles et dérivée des fonctions composées (Première)
\item Étude de variations et limites de fonctions (Première et début de Terminale)
\item Résolution d'équations et d'inéquations du premier et du second degré
\item Fonction inverse x ↦ 1/x : ensemble de définition, continuité, représentation graphique
\end{itemize}
\end{prerequis}

\newpage

\coursec{Définition et premières propriétés de la fonction ln}

Au Cameroun, le pH de l'eau de forage contrôlée par les services d'hygiène, l'échelle de Richter utilisée pour évaluer les séismes ressentis dans la région de l'Ouest, et le temps nécessaire pour doubler un capital placé à la microfinance reposent tous sur un même outil mathématique : le \cle{logarithme népérien}. Comment une seule fonction peut-elle transformer des multiplications en additions et rendre calculables des phénomènes qui croissent très lentement ou très vite ? Cette leçon construit cette fonction à partir d'un problème d'\cle{aire} sous une hyperbole et montre comment elle résout des équations que les outils de Première ne permettaient pas de traiter.

\courssub{Un problème d'aire sous l'hyperbole}

\textbf{Le point de départ : une fonction sans primitive « simple ».}\par

Depuis la classe de Première, tu sais dériver les fonctions polynômes, les puissances, les racines carrées. Mais voici une question embarrassante : quelle fonction a pour dérivée $x \mapsto \dfrac{1}{x}$ ? Essayons. La dérivée de $x^2$ est $2x$ ; celle de $x$ est $1$ ; celle d'une constante est $0$. En dérivant, les puissances \emph{descendent} d'un cran : on ne peut jamais obtenir $\dfrac{1}{x} = x^{-1}$ en dérivant une puissance de $x$, car la dérivée de $x^0 = 1$ est $0$, pas $x^{-1}$. Aucune des fonctions « usuelles » de Première ne convient.

Pourtant, la fonction $f : x \mapsto \dfrac{1}{x}$ est \cle{continue} sur $]0\,;\,+\infty[$ (c'est une fonction rationnelle dont le dénominateur ne s'annule pas sur cet intervalle). Or un théorème fondamental de l'analyse affirme que \textbf{toute fonction continue sur un intervalle y admet des primitives}. Il existe donc bel et bien des fonctions dont la dérivée est $\dfrac{1}{x}$ : simplement, ce sont des fonctions \emph{nouvelles}, que nous allons construire.

\begin{experience}[Découverte : mesurer une aire sous l'hyperbole]
Considérons la courbe de $t \mapsto \dfrac{1}{t}$ (une hyperbole) et intéressons-nous à l'aire du domaine compris entre la courbe, l'axe des abscisses, et les droites verticales d'abscisses $1$ et $x$, pour un réel $x > 1$ fixé.
\begin{enumerate}
\item Pour $x = 1$, que vaut cette aire ? (Le domaine est réduit à un segment.)
\item Pour $x = 2$, encadre l'aire entre celle du rectangle de largeur $1$ et de hauteur $\dfrac{1}{2}$ (sous la courbe) et celle du rectangle de largeur $1$ et de hauteur $1$ (au-dessus de la courbe). Qu'en déduis-tu ?
\item Si l'on note $F(x)$ cette aire, que vaut $F(1)$ ? Le nombre $F(x)$ augmente-t-il quand $x$ augmente ?
\end{enumerate}
Tu viens de définir, sans le savoir, la fonction logarithme népérien : $F(x) = \ln x$.
\end{experience}

L'idée est donc la suivante : puisque $\dfrac{1}{x}$ admet des primitives sur $]0\,;\,+\infty[$, et que ces primitives différent toutes d'une constante, on en choisit \textbf{une seule}, de manière canonique : celle qui vaut $0$ en $x = 1$. C'est exactement l'aire mesurée à partir de l'abscisse $1$.

\courssub{Définition de la fonction logarithme népérien}

\begin{definition}[Fonction logarithme népérien]
On appelle \cle{fonction logarithme népérien}, notée $\ln$, l'\textbf{unique} primitive de la fonction $x \mapsto \dfrac{1}{x}$ sur l'intervalle $]0\,;\,+\infty[$ qui \textbf{s'annule en} $1$. Autrement dit, pour tout réel $x > 0$ :
\[
\ln x = \int_{1}^{x} \frac{1}{t}\,\mathrm{d}t .
\]
L'\cle{ensemble de définition} de $\ln$ est donc $\mathcal{D}_{\ln} = ]0\,;\,+\infty[$.
\end{definition}

\begin{attention}[Le domaine avant tout !]
L'erreur la plus fréquente — et la plus sanctionnée au Baccalauréat — est d'écrire $\ln x$ pour un nombre $x \leq 0$. \textbf{Le logarithme népérien n'existe que pour les réels strictement positifs.} Avant tout calcul, toute équation ou toute étude de fonction contenant $\ln$, commence par déterminer l'ensemble de validité. Par exemple, $\ln(x-3)$ n'a de sens que si $x - 3 > 0$, c'est-à-dire $x > 3$. Écrire $\ln(-2)$ ou $\ln 0$ est une absurdité mathématique.
\end{attention}

\begin{propriete}[Dérivabilité et dérivée de ln]
La fonction $\ln$ est \cle{dérivable} sur $]0\,;\,+\infty[$ (donc \cle{continue} sur cet intervalle), et pour tout $x > 0$ :
\[
(\ln)'(x) = \frac{1}{x}.
\]
\end{propriete}

\begin{experience}[Démonstration (immédiate, mais à savoir rédiger)]
C'est une conséquence directe de la définition. Dire que $\ln$ est \emph{une primitive} de $x \mapsto \dfrac{1}{x}$ sur $]0\,;\,+\infty[$ signifie exactement deux choses :
\begin{enumerate}
\item $\ln$ est dérivable sur $]0\,;\,+\infty[$ ;
\item sa dérivée est la fonction $x \mapsto \dfrac{1}{x}$, c'est-à-dire $(\ln)'(x) = \dfrac{1}{x}$ pour tout $x > 0$.
\end{enumerate}
De plus, toute fonction dérivable sur un intervalle est continue sur cet intervalle, donc $\ln$ est continue sur $]0\,;\,+\infty[$. \hfill$\blacksquare$
\end{experience}

\courssub{Interprétation géométrique : le logarithme comme une aire}

La définition intégrale donne une image mentale précieuse : pour $x > 1$, le nombre $\ln x$ est l'\cle{aire} (en unités d'aire) du domaine délimité par l'hyperbole d'équation $y = \dfrac{1}{t}$, l'axe des abscisses et les droites $t = 1$ et $t = x$.

\begin{popfigure}
\begin{center}
\begin{tikzpicture}[scale=1.6]
% Aire hachurée sous y=1/x entre 1 et 2.5
\fill[pattern=north east lines, pattern color=popblue, opacity=0.55]
  plot[domain=1:2.5, samples=100] (\x, {1/\x}) -- (2.5,0) -- (1,0) -- cycle;
% Axes
\draw[->, popdark, thick] (-0.4,0) -- (4.4,0) node[below right] {$t$};
\draw[->, popdark, thick] (0,-0.4) -- (0,2.6) node[above left] {$y$};
% Hyperbole
\draw[popink, very thick, domain=0.42:4.2, samples=150] plot (\x, {1/\x}) node[right] {$y=\dfrac{1}{t}$};
% Lignes verticales
\draw[popblue, thick] (1,0) -- (1,1);
\draw[popblue, thick] (2.5,0) -- (2.5,0.4);
% Graduations
\draw[popdark] (1,0.05) -- (1,-0.05) node[below] {$1$};
\draw[popdark] (2.5,0.05) -- (2.5,-0.05) node[below] {$x$};
\draw[popdark] (0.05,1) -- (-0.05,1) node[left] {$1$};
\draw[popdark] (0,0) node[below left] {$O$};
% Légende de l'aire
\node[popblue] at (1.75,0.55) {$\mathcal{A} = \ln x$};
\end{tikzpicture}
\end{center}
\legende{Pour $x > 1$, $\ln x$ est l'aire du domaine hachuré situé sous l'hyperbole $y = \dfrac{1}{t}$, entre les abscisses $1$ et $x$.}
\end{popfigure}

\begin{aretenir}[Lecture géométrique du signe]
\begin{itemize}
\item Si $x > 1$, l'intégrale $\displaystyle\int_1^x \frac{\mathrm{d}t}{t}$ est l'aire d'un domaine situé \emph{au-dessus} de l'axe des abscisses : elle est \textbf{strictement positive}.
\item Si $0 < x < 1$, on intègre « à rebours » (de $1$ vers une borne plus petite) : $\ln x = -\displaystyle\int_x^1 \frac{\mathrm{d}t}{t}$, qui est \textbf{strictement négatif}.
\item Si $x = 1$, le domaine est réduit à un segment, d'aire nulle : $\ln 1 = 0$.
\end{itemize}
\end{aretenir}

\courssub{Premières valeurs remarquables et le nombre e}

\begin{propriete}[Valeurs remarquables]
\begin{itemize}
\item $\ln 1 = 0$ (par définition même : la primitive choisie s'annule en $1$).
\item Il existe un \textbf{unique} réel, noté $\mathrm{e}$, tel que $\ln \mathrm{e} = 1$. On a $\mathrm{e} \approx \num{2.718}$.
\end{itemize}
\end{propriete}

L'existence et l'unicité de $\mathrm{e}$ seront justifiées rigoureusement dans la section 4 grâce au théorème des valeurs intermédiaires (la fonction $\ln$ est continue, strictement croissante, et prend des valeurs arbitrairement grandes). Retenons dès maintenant l'idée géométrique : $\mathrm{e}$ est l'abscisse telle que l'aire sous l'hyperbole entre $1$ et $\mathrm{e}$ vaut exactement $1$ unité d'aire.

\begin{exemplebox}[Premières valeurs à la calculatrice]
À l'aide de la touche \textbf{ln} de la calculatrice, on obtient :
\begin{center}
\begin{tabularx}{\linewidth}{|l|X|X|X|X|}
\hline
\rowcolor{popblueL}
$x$ & \centering $1$ & \centering $2$ & \centering $\mathrm{e}$ & \centering\arraybackslash $10$ \\
\hline
$\ln x$ (arrondi à $10^{-3}$) & \centering $0$ & \centering $\num{0.693}$ & \centering $1$ & \centering\arraybackslash $\num{2.303}$ \\
\hline
\end{tabularx}
\end{center}
On note aussi $\ln \num{0.5} \approx \num{-0.693}$ : cohérent avec le fait que $\num{0.5} < 1$, donc son logarithme est négatif. Remarque au passage : $\ln \num{0.5} = -\ln 2$ — ce n'est pas un hasard, la section 2 l'expliquera.
\end{exemplebox}

\begin{savaistu}[Le nombre e, star de la finance]
Le nombre $\mathrm{e} \approx \num{2.71828}$ est appelé \cle{nombre d'Euler}, en hommage au mathématicien suisse Leonhard Euler (XVIII\textsuperscript{e} siècle), même si c'est John Napier (Neper) qui donna son nom au logarithme « népérien » au début du XVII\textsuperscript{e} siècle. Le nombre $\mathrm{e}$ apparaît naturellement dans le calcul des \textbf{intérêts composés} : si une microfinance de Yaoundé plaçait un capital à un taux de $100\,\%$ par an en composant les intérêts de plus en plus souvent (chaque mois, chaque jour, chaque seconde...), le capital au bout d'un an tendrait vers $\mathrm{e}$ fois le capital initial, et non vers l'infini ! C'est pourquoi $\mathrm{e}$ gouverne tous les phénomènes de croissance continue : populations, radioactivité, épidémies.
\end{savaistu}

\courssub{Sens de variation et signe de la fonction ln}

\begin{propriete}[Stricte croissance de ln]
La fonction $\ln$ est \textbf{strictement croissante} sur $]0\,;\,+\infty[$.
\end{propriete}

\begin{experience}[Démonstration (exigible)]
Pour tout $x > 0$, on a $(\ln)'(x) = \dfrac{1}{x}$. Or $x > 0$ implique $\dfrac{1}{x} > 0$ : la dérivée est \textbf{strictement positive} sur $]0\,;\,+\infty[$. La fonction $\ln$ est donc strictement croissante sur cet intervalle. \hfill$\blacksquare$
\end{experience}

Cette stricte croissance, combinée à la valeur $\ln 1 = 0$, donne immédiatement le signe de $\ln x$.

\begin{propriete}[Signe de ln x]
\begin{itemize}
\item Si $0 < x < 1$, alors $\ln x < 0$ ;
\item si $x = 1$, alors $\ln x = 0$ ;
\item si $x > 1$, alors $\ln x > 0$.
\end{itemize}
\end{propriete}

\begin{experience}[Démonstration (exigible)]
La fonction $\ln$ est strictement croissante sur $]0\,;\,+\infty[$ et $\ln 1 = 0$.
\begin{itemize}
\item Si $x > 1$, la stricte croissance donne $\ln x > \ln 1 = 0$.
\item Si $0 < x < 1$, la stricte croissance donne $\ln x < \ln 1 = 0$.
\end{itemize}
On retrouve exactement ce que l'interprétation géométrique par les aires nous avait fait deviner. \hfill$\blacksquare$
\end{experience}

\begin{attention}[Deux conséquences à ne jamais oublier]
\begin{enumerate}
\item \textbf{Conservation de l'ordre} : pour tous réels $a > 0$ et $b > 0$,
\[
\ln a < \ln b \iff a < b .
\]
C'est la clé de toutes les inéquations avec $\ln$ (section 6) : on peut « enlever » les $\ln$ \emph{sans changer le sens de l'inégalité}, car $\ln$ est croissante.
\item \textbf{Unicité} : $\ln a = \ln b \iff a = b$ (pour $a, b > 0$). La fonction $\ln$ est \cle{injective} : deux nombres distincts ont des logarithmes distincts.
\end{enumerate}
\end{attention}

\courssub{Exercices résolus de mise en route}

\begin{exoresolu}[Domaine de définition d'une expression avec ln]
Déterminer l'ensemble de définition de la fonction $f$ définie par $f(x) = \ln(2x - 6) + \ln(x + 1)$.
\tcblower
\textbf{Solution détaillée.}

Une expression $\ln(\text{quelque chose})$ n'existe que si ce « quelque chose » est \textbf{strictement positif}. Il faut donc que \emph{les deux} conditions soient réalisées simultanément :
\[
\begin{cases}
2x - 6 > 0 \\
x + 1 > 0
\end{cases}
\iff
\begin{cases}
x > 3 \\
x > -1
\end{cases}
\]
Les deux conditions doivent être vraies en même temps : on prend la plus restrictive, c'est-à-dire $x > 3$.
\[
\mathcal{D}_f = ]3\,;\,+\infty[.
\]
Vérification mentale : pour $x = 4$, $2x - 6 = 2 > 0$ et $x + 1 = 5 > 0$, tout va bien ; pour $x = 2$, $2x - 6 = -2 < 0$, l'expression n'existe pas.
\end{exoresolu}

\begin{exoresolu}[Comparaison et signe sans calculatrice]
\begin{enumerate}
\item Comparer $\ln 5$ et $\ln 7$.
\item Déterminer le signe de $\ln \num{0.8}$ et celui de $\ln \dfrac{3}{2}$.
\item Résoudre dans $]0\,;\,+\infty[$ l'équation $\ln x = 0$.
\end{enumerate}
\tcblower
\textbf{Solution détaillée.}
\begin{enumerate}
\item La fonction $\ln$ est strictement croissante et $5 < 7$, donc $\ln 5 < \ln 7$.
\item $\num{0.8} < 1$, donc $\ln \num{0.8} < 0$. De même $\dfrac{3}{2} = \num{1.5} > 1$, donc $\ln \dfrac{3}{2} > 0$.
\item On sait que $\ln 1 = 0$. Par injectivité de $\ln$ (stricte croissance), l'équation $\ln x = 0 = \ln 1$ équivaut à $x = 1$. L'unique solution est $x = 1$.
\end{enumerate}
\end{exoresolu}

\begin{exoresolu}[Estimation d'une aire]
Un technicien de CAMTEL modélise l'atténuation d'un signal par une formule faisant intervenir $\ln 3$. En utilisant l'interprétation géométrique, encadrer $\ln 3$ entre deux aires de rectangles, puis donner la valeur approchée à la calculatrice.
\tcblower
\textbf{Solution détaillée.}

$\ln 3$ est l'aire sous l'hyperbole $y = \dfrac{1}{t}$ entre $t = 1$ et $t = 3$, sur un intervalle de largeur $2$. Sur $[1\,;\,3]$, la fonction $t \mapsto \dfrac{1}{t}$ est décroissante, donc :
\[
\frac{1}{3} \leq \frac{1}{t} \leq 1 \quad \text{pour tout } t \in [1\,;\,3].
\]
En multipliant par la largeur $2$ de l'intervalle :
\[
2 \times \frac{1}{3} \leq \ln 3 \leq 2 \times 1,
\quad \text{c'est-à-dire} \quad
\frac{2}{3} \leq \ln 3 \leq 2.
\]
La calculatrice donne $\ln 3 \approx \num{1.099}$, valeur bien comprise entre $\num{0.667}$ et $2$. L'encadrement est grossier (les rectangles épousent mal la courbe) mais il permet de \emph{vérifier} un résultat et de détecter une erreur de saisie à la calculatrice.
\end{exoresolu}

\begin{aretenir}[L'essentiel de cette section]
\begin{itemize}
\item $\ln$ est \textbf{l'unique primitive} de $x \mapsto \dfrac{1}{x}$ sur $]0\,;\,+\infty[$ qui s'annule en $1$ : $\ln x = \displaystyle\int_1^x \frac{\mathrm{d}t}{t}$.
\item $\ln x$ n'existe que pour $x > 0$ : toujours chercher le domaine en premier.
\item $\ln$ est dérivable et continue sur $]0\,;\,+\infty[$, avec $(\ln)'(x) = \dfrac{1}{x} > 0$ : $\ln$ est \textbf{strictement croissante}.
\item $\ln 1 = 0$ ; $\ln \mathrm{e} = 1$ avec $\mathrm{e} \approx \num{2.718}$.
\item $\ln x < 0$ si $0 < x < 1$ ; $\ln x > 0$ si $x > 1$.
\item Géométriquement, pour $x > 1$, $\ln x$ est l'aire sous l'hyperbole $y = \dfrac{1}{t}$ entre $1$ et $x$.
\end{itemize}
\end{aretenir}

\begin{exercice}[Pour s'entraîner]
\begin{enumerate}
\item Déterminer l'ensemble de définition de chaque fonction : $g(x) = \ln(5 - x)$ ; $h(x) = \ln(x^2 - 4)$ ; $k(x) = \dfrac{\ln x}{x - 1}$.
\item Sans calculatrice, ranger dans l'ordre croissant : $\ln \dfrac{1}{2}$, $\ln 1$, $\ln 3$, $\ln \num{0.9}$.
\item Résoudre dans $]0\,;\,+\infty[$ : a) $\ln x = 1$ ; b) $\ln x < 0$ ; c) $\ln x \geq \ln 4$.
\item En utilisant la méthode des rectangles, encadrer $\ln 2$ entre $\dfrac{1}{2}$ et $1$, puis vérifier à la calculatrice.
\end{enumerate}
\end{exercice}

\begin{corrige}[Exercice : Pour s'entraîner]
\begin{enumerate}
\item $g$ : il faut $5 - x > 0$, soit $x < 5$ : $\mathcal{D}_g = ]-\infty\,;\,5[$. $h$ : il faut $x^2 - 4 > 0$, soit $x < -2$ ou $x > 2$ : $\mathcal{D}_h = ]-\infty\,;\,-2[ \cup ]2\,;\,+\infty[$. $k$ : il faut $x > 0$ et $x \neq 1$ : $\mathcal{D}_k = ]0\,;\,1[ \cup ]1\,;\,+\infty[$.
\item La fonction $\ln$ est strictement croissante et $\dfrac{1}{2} < \num{0.9} < 1 < 3$, donc : $\ln \dfrac{1}{2} < \ln \num{0.9} < \ln 1 < \ln 3$.
\item a) $\ln x = 1 = \ln \mathrm{e}$ donne $x = \mathrm{e}$. b) $\ln x < 0 = \ln 1$ donne, par stricte croissance, $x < 1$ : $\mathcal{S} = ]0\,;\,1[$. c) $\ln x \geq \ln 4$ donne $x \geq 4$ : $\mathcal{S} = [4\,;\,+\infty[$.
\item Sur $[1\,;\,2]$, $\dfrac{1}{2} \leq \dfrac{1}{t} \leq 1$ ; en multipliant par la largeur $1$ : $\dfrac{1}{2} \leq \ln 2 \leq 1$. La calculatrice donne $\ln 2 \approx \num{0.693}$, conforme à l'encadrement.
\end{enumerate}
\end{corrige}

\coursec{Propriétés algébriques de ln}

Dans la section précédente, nous avons défini la fonction logarithme népérien comme l'unique primitive de $x \mapsto \dfrac{1}{x}$ sur $]0\,;\,+\infty[$ qui s'annule en $1$ :
\[
\ln(x) = \int_1^x \frac{1}{t}\,\mathrm{d}t, \qquad \ln(1) = 0, \qquad \ln'(x) = \frac{1}{x}.
\]
Cette définition par une intégrale est précise, mais elle cache un trésor : la fonction $\ln$ possède des \cle{propriétés algébriques} extraordinaires qui la rendent irremplaçable en calcul. L'idée centrale est la suivante : \textbf{le logarithme transforme les multiplications en additions}. Or additionner est bien plus facile que multiplier ! C'est exactement ce « miracle » que nous allons démontrer et exploiter dans toute cette section. Tu verras qu'une seule formule — la relation fonctionnelle $\ln(ab) = \ln a + \ln b$ — engendre toutes les autres, comme un tronc d'arbre engendre ses branches.

\courssub{La relation fonctionnelle fondamentale}

Commençons par une intuition. Prenons $a = 2$ et $b = 3$. À la calculatrice : $\ln 2 \approx 0,693$, $\ln 3 \approx 1,099$ et $\ln 6 \approx 1,792$. Or $0,693 + 1,099 = 1,792$. Coïncidence ? Essayons avec $a = 5$ et $b = 7$ : $\ln 5 + \ln 7 \approx 1,609 + 1,946 = 3,555$ et $\ln 35 \approx 3,555$. Encore exact ! Ce phénomène n'est pas un hasard : c'est un théorème, que nous allons démontrer proprement.

\begin{propriete}[Relation fonctionnelle fondamentale (théorème)]
Pour tous réels $a > 0$ et $b > 0$ :
\[
\ln(ab) = \ln a + \ln b.
\]
Autrement dit, \textbf{le logarithme d'un produit est la somme des logarithmes}. Cette démonstration est classique au Baccalauréat : il faut savoir la refaire.
\end{propriete}

\begin{experience}[Démonstration guidée du théorème fondamental]
Fixons un réel $a > 0$ quelconque. L'astuce consiste à faire varier $b$ : on étudie la fonction
\[
\varphi : x \longmapsto \ln(ax) - \ln x \quad \text{définie sur } ]0\,;\,+\infty[.
\]
\begin{enumerate}
\item \textbf{Dérivabilité.} La fonction $x \mapsto ax$ est dérivable et strictement positive sur $]0\,;\,+\infty[$, donc par composition $\varphi$ est dérivable sur $]0\,;\,+\infty[$.
\item \textbf{Calcul de la dérivée.} Par la formule de dérivation d'une composée $(\ln \circ u)' = \dfrac{u'}{u}$ avec $u(x) = ax$, donc $u'(x) = a$ :
\[
\varphi'(x) = \frac{a}{ax} - \frac{1}{x} = \frac{1}{x} - \frac{1}{x} = 0.
\]
\item \textbf{Conclusion sur $\varphi$.} La dérivée de $\varphi$ est nulle sur l'\emph{intervalle} $]0\,;\,+\infty[$, donc $\varphi$ est \textbf{constante} sur cet intervalle.
\item \textbf{Détermination de la constante.} Évaluons en $x = 1$ :
\[
\varphi(1) = \ln(a \times 1) - \ln 1 = \ln a - 0 = \ln a.
\]
Donc pour tout $x > 0$ : $\ln(ax) - \ln x = \ln a$, c'est-à-dire $\ln(ax) = \ln a + \ln x$.
\item \textbf{Conclusion.} En prenant $x = b$ (réel strictement positif quelconque) : $\ln(ab) = \ln a + \ln b$. \hfill$\blacksquare$
\end{enumerate}
\end{experience}

\begin{attention}[La dérivée nulle implique la constance... sur un intervalle !]
L'étape 3 de la démonstration utilise le théorème : « si $f' = 0$ sur un \textbf{intervalle}, alors $f$ est constante sur cet intervalle ». Ici, $]0\,;\,+\infty[$ est bien un intervalle, donc tout va bien. Retiens aussi l'idée générale de la méthode : pour prouver qu'une expression est constante, on montre que sa dérivée est nulle, puis on calcule la constante en un point bien choisi (souvent $x = 1$ pour le logarithme).
\end{attention}

\begin{exemplebox}[Vérification numérique]
Avec $a = 8$ et $b = 4$ : $\ln 8 + \ln 4 = 3\ln 2 + 2\ln 2 = 5\ln 2$ (nous justifierons bientôt la formule des puissances utilisée ici). Et $\ln(8 \times 4) = \ln 32 = \ln(2^5) = 5\ln 2$. Les deux membres coïncident, comme prévu.
\end{exemplebox}

\courssub{Logarithme de l'inverse et du quotient}

Toutes les propriétés qui suivent sont des \textbf{conséquences directes} de la relation fondamentale. C'est la beauté de la théorie : une seule démonstration par dérivation, puis tout découle par de simples calculs algébriques.

\begin{propriete}[Logarithme de l'inverse]
Pour tout réel $a > 0$ :
\[
\ln\left(\frac{1}{a}\right) = -\ln a.
\]
\end{propriete}

\textbf{Démonstration.} Appliquons la relation fondamentale au produit $a \times \dfrac{1}{a} = 1$ :
\[
\ln\left(a \times \frac{1}{a}\right) = \ln a + \ln\left(\frac{1}{a}\right).
\]
Mais le membre de gauche vaut $\ln 1 = 0$. Donc $\ln a + \ln\left(\dfrac{1}{a}\right) = 0$, d'où $\ln\left(\dfrac{1}{a}\right) = -\ln a$. \hfill$\blacksquare$

\begin{propriete}[Logarithme d'un quotient]
Pour tous réels $a > 0$ et $b > 0$ :
\[
\ln\left(\frac{a}{b}\right) = \ln a - \ln b.
\]
\end{propriete}

\textbf{Démonstration.} Un quotient est un produit par l'inverse : $\dfrac{a}{b} = a \times \dfrac{1}{b}$. On applique successivement les deux propriétés précédentes :
\[
\ln\left(\frac{a}{b}\right) = \ln\left(a \times \frac{1}{b}\right) = \ln a + \ln\left(\frac{1}{b}\right) = \ln a - \ln b. \qquad \blacksquare
\]

\begin{exemplebox}[Application directe]
$\ln\left(\dfrac{1}{5}\right) = -\ln 5$ et $\ln\left(\dfrac{12}{3}\right) = \ln 12 - \ln 3$. Vérifions : $\ln 12 - \ln 3 = \ln(4 \times 3) - \ln 3 = \ln 4 + \ln 3 - \ln 3 = \ln 4$, et effectivement $\dfrac{12}{3} = 4$.
\end{exemplebox}

\courssub{Logarithme d'une puissance et d'une racine}

\begin{propriete}[Logarithme d'une puissance entière]
Pour tout réel $a > 0$ et tout entier $n \in \mathbb{N}$ :
\[
\ln(a^n) = n \ln a.
\]
Cette formule s'étend à $n \in \mathbb{Z}$ tout entier, et même à $n \in \mathbb{Q}$. (On admettra qu'elle reste valable pour tout exposant réel, une fois les puissances réelles définies.)
\end{propriete}

\textbf{Démonstration (pour $n \in \mathbb{N}$, par récurrence).}
\begin{itemize}
\item \emph{Initialisation} : pour $n = 0$, $\ln(a^0) = \ln 1 = 0 = 0 \times \ln a$. La propriété est vraie.
\item \emph{Hérédité} : supposons $\ln(a^n) = n\ln a$ pour un entier $n$ fixé. Alors, par la relation fondamentale :
\[
\ln(a^{n+1}) = \ln(a^n \times a) = \ln(a^n) + \ln a = n\ln a + \ln a = (n+1)\ln a.
\]
\item \emph{Conclusion} : par récurrence, $\ln(a^n) = n\ln a$ pour tout $n \in \mathbb{N}$. \hfill$\blacksquare$
\end{itemize}

\emph{Extension à $\mathbb{Z}$} : si $n$ est un entier négatif, écrivons $n = -m$ avec $m \in \mathbb{N}$. Alors
\[
\ln(a^n) = \ln\left(\frac{1}{a^m}\right) = -\ln(a^m) = -m\ln a = n\ln a.
\]

\emph{Extension à $\mathbb{Q}$} : pour $n = \dfrac{p}{q}$ avec $p \in \mathbb{Z}$ et $q \in \mathbb{N}^*$, on remarque que $\left(a^{p/q}\right)^q = a^p$, donc $q\ln\left(a^{p/q}\right) = p\ln a$, d'où $\ln\left(a^{p/q}\right) = \dfrac{p}{q}\ln a$.

\begin{propriete}[Logarithme d'une racine carrée]
Pour tout réel $a > 0$ :
\[
\ln(\sqrt{a}) = \frac{1}{2}\ln a.
\]
\end{propriete}

\textbf{Démonstration.} Puisque $\sqrt{a} \times \sqrt{a} = a$, la relation fondamentale donne :
\[
\ln a = \ln(\sqrt{a} \times \sqrt{a}) = \ln(\sqrt{a}) + \ln(\sqrt{a}) = 2\ln(\sqrt{a}),
\]
d'où $\ln(\sqrt{a}) = \dfrac{1}{2}\ln a$. \hfill$\blacksquare$

\begin{exemplebox}[Calculs types]
$\ln(2^{10}) = 10\ln 2$ ; $\ln\left(\dfrac{1}{9}\right) = -\ln 9 = -2\ln 3$ ; $\ln(\sqrt{5}) = \dfrac{1}{2}\ln 5$ ; $\ln(\sqrt[3]{7}) = \ln\left(7^{1/3}\right) = \dfrac{1}{3}\ln 7$.
\end{exemplebox}

L'arbre ci-dessous résume la structure logique de toute cette section : une seule racine, la relation fondamentale, et tout le reste en découle.

\begin{popfigure}
\begin{center}
\begin{tikzpicture}[
  grow=down,
  level distance=2.1cm,
  every node/.style={draw, rounded corners=3pt, fill=popblueL, text=popdark, font=\small, align=center, inner sep=4pt},
  edge from parent/.style={draw=popblue, thick, -{Stealth}},
  level 1/.style={sibling distance=5.2cm},
  level 2/.style={sibling distance=3.4cm}
]
\node[fill=popblue, text=white, font=\small\bfseries, align=center] {Relation fondamentale\\ $\ln(ab)=\ln a+\ln b$}
  child { node[align=center] {$b=\dfrac{1}{a}$\\[2pt] $\ln\left(\dfrac{1}{a}\right)=-\ln a$}
    child { node[align=center] {$\dfrac{a}{b}=a\times\dfrac{1}{b}$\\[2pt] $\ln\left(\dfrac{a}{b}\right)=\ln a-\ln b$} }
  }
  child { node[align=center] {Récurrence sur $n\in\mathbb{N}$\\[2pt] $\ln(a^n)=n\ln a$}
    child { node[align=center] {Extension à $\mathbb{Z}$, $\mathbb{Q}$\\[2pt] (admis : $\mathbb{R}$)} }
    child { node[align=center] {$a=\sqrt{a}\times\sqrt{a}$\\[2pt] $\ln(\sqrt{a})=\dfrac{1}{2}\ln a$} }
  };
\end{tikzpicture}
\end{center}
\legende{Arbre de démonstration : toutes les propriétés algébriques de $\ln$ naissent de la relation fondamentale $\ln(ab)=\ln a+\ln b$.}
\end{popfigure}

\courssub{Simplification d'expressions logarithmiques}

En pratique, au Baccalauréat, on te demandera de faire deux types de manipulations inverses l'une de l'autre :
\begin{itemize}
\item \textbf{Regrouper} : écrire une somme ou une différence de logarithmes sous la forme d'un \emph{seul} logarithme ;
\item \textbf{Développer} : écrire le logarithme d'une expression compliquée comme une combinaison de logarithmes simples (par exemple en fonction de $\ln 2$ et $\ln 3$).
\end{itemize}

\begin{methode}[Simplifier une expression avec des logarithmes]
\begin{enumerate}
\item Repérer les coefficients devant les logarithmes et les faire « rentrer » en exposants : $n\ln a = \ln(a^n)$.
\item Regrouper les termes \emph{additionnés} en un produit : $\ln a + \ln b = \ln(ab)$.
\item Regrouper les termes \emph{soustraits} en un quotient : $\ln a - \ln b = \ln\left(\dfrac{a}{b}\right)$.
\item Simplifier l'argument obtenu (réduire la fraction, factoriser en puissances de nombres premiers).
\item Éventuellement, redévelopper pour faire apparaître $\ln 2$, $\ln 3$, etc.
\end{enumerate}
\end{methode}

\begin{exemplebox}[Exemple chiffré fondamental]
Simplifions $A = \ln 8 + \ln 4 - \ln 2$.
\[
A = \ln(8 \times 4) - \ln 2 = \ln\left(\frac{8 \times 4}{2}\right) = \ln 16 = \ln(2^4) = 4\ln 2.
\]
On a regroupé en un seul logarithme, puis réécrit le résultat en fonction de $\ln 2$.
\end{exemplebox}

\begin{exoresolu}[Simplifications complètes]
\textbf{1.} Écrire $B = 2\ln 3 + \ln 5 - \dfrac{1}{2}\ln 25$ sous la forme d'un seul logarithme, puis simplifier.

\textbf{2.} Exprimer $C = \ln 48$ en fonction de $\ln 2$ et $\ln 3$.
\tcblower
\textbf{1.} On applique la méthode étape par étape :
\begin{align*}
B &= \ln(3^2) + \ln 5 - \ln\left(25^{1/2}\right) & &\text{(les coefficients deviennent des exposants)}\\
&= \ln 9 + \ln 5 - \ln(\sqrt{25})\\
&= \ln(9 \times 5) - \ln 5 & &\text{(regroupement en produit)}\\
&= \ln\left(\frac{45}{5}\right) & &\text{(regroupement en quotient)}\\
&= \ln 9 = \ln(3^2) = 2\ln 3.
\end{align*}

\textbf{2.} On décompose $48$ en facteurs premiers : $48 = 16 \times 3 = 2^4 \times 3$. Donc :
\[
C = \ln(2^4 \times 3) = \ln(2^4) + \ln 3 = 4\ln 2 + \ln 3.
\]
\end{exoresolu}

\begin{attention}[Le piège le plus fréquent : les sommes]
Le logarithme ne possède \textbf{aucune propriété} pour les sommes :
\[
\ln(a+b) \neq \ln a + \ln b \qquad \text{et} \qquad \ln(a+b) \neq \ln a \times \ln b.
\]
Par exemple, $\ln(2+3) = \ln 5 \approx 1,609$, alors que $\ln 2 + \ln 3 = \ln 6 \approx 1,792$. Ce n'est pas la même chose !

De même, $\ln(a) \times \ln(b) \neq \ln(ab)$ : ne confonds pas le \emph{produit} des logarithmes avec le logarithme du produit. Enfin, $\dfrac{\ln a}{\ln b} \neq \ln\left(\dfrac{a}{b}\right)$ : le quotient des logarithmes n'est pas le logarithme du quotient. Retiens : $\ln$ simplifie les produits, les quotients et les puissances \emph{à l'intérieur} de son argument, et rien d'autre.
\end{attention}

\begin{attention}[Conditions d'existence]
Toutes les formules de cette section exigent des arguments \textbf{strictement positifs}. Écrire $\ln(ab) = \ln a + \ln b$ n'a de sens que si $a > 0$ et $b > 0$. Ainsi, $\ln\left((-2)\times(-3)\right) = \ln 6$ existe, mais on ne peut \emph{pas} écrire $\ln(-2) + \ln(-3)$, qui n'a aucun sens. Dans les exercices, précise toujours les conditions d'existence avant de manipuler les formules.
\end{attention}

\courssub{Complément culturel : l'unicité du logarithme}

\begin{savaistu}[Pourquoi les logarithmes ont changé le monde]
C'est précisément la transformation \emph{produit} $\rightarrow$ \emph{somme} qui a rendu les logarithmes révolutionnaires. Au XVII\textsuperscript{e} siècle, les astronomes comme Kepler devaient multiplier à la main des nombres à dix chiffres pour calculer les orbites des planètes : un travail de plusieurs mois, plein d'erreurs. Grâce aux \textbf{tables de logarithmes} de Napier (1614) et Briggs, on remplaçait chaque multiplication par une simple addition : on lisait $\ln a$ et $\ln b$ dans la table, on additionnait, puis on cherchait dans la table le nombre dont le logarithme était cette somme. Laplace aurait dit que les logarithmes avaient « doublé la vie des astronomes ». Les règles à calcul, utilisées par les ingénieurs jusque dans les années 1970, reposent sur le même principe. Aujourd'hui encore, les échelles logarithmiques (pH en chimie, décibels en acoustique, magnitude des séismes) exploitent cette même propriété.
\end{savaistu}

\begin{savaistu}[Caractérisation de la fonction ln (complément admis)]
On peut démontrer — nous l'admettons ici — que la fonction $\ln$ est l'\textbf{unique} fonction $f$ dérivable sur $]0\,;\,+\infty[$ vérifiant à la fois :
\[
f(ab) = f(a) + f(b) \quad \text{pour tous } a, b > 0, \qquad \text{et} \qquad f'(1) = 1.
\]
Autrement dit, la relation fonctionnelle n'est pas un simple « bonus » : elle \emph{caractérise} entièrement le logarithme népérien. Toute fonction dérivable qui transforme les produits en sommes est de la forme $x \mapsto k\ln x$, et la condition $f'(1) = 1$ fixe $k = 1$. C'est ce point de vue que les mathématiciens adoptent pour définir les logarithmes de manière purement algébrique.
\end{savaistu}

\begin{exercice}[À toi de jouer]
\textbf{1.} Simplifier sous la forme d'un seul logarithme : $D = \ln 12 + \ln 3 - \ln 4 - \ln 9$.

\textbf{2.} Exprimer en fonction de $\ln 2$ et $\ln 5$ : $E = \ln 200 - \ln 8 + \dfrac{1}{2}\ln 25$.

\textbf{3.} Montrer que pour tout $x > 0$ : $\ln\left(x^2\right) - \ln\left(\dfrac{1}{x}\right) = 3\ln x$.
\end{exercice}

\begin{corrige}[Exercice]
\textbf{1.} $D = \ln\left(\dfrac{12 \times 3}{4 \times 9}\right) = \ln\left(\dfrac{36}{36}\right) = \ln 1 = 0$.

\textbf{2.} $E = \ln 200 - \ln 8 + \ln(\sqrt{25}) = \ln(2^3 \times 5^2) - \ln(2^3) + \ln 5 = 3\ln 2 + 2\ln 5 - 3\ln 2 + \ln 5 = 3\ln 5$.

\textbf{3.} Pour $x > 0$ : $\ln(x^2) - \ln\left(\dfrac{1}{x}\right) = 2\ln x - (-\ln x) = 2\ln x + \ln x = 3\ln x$.
\end{corrige}

\begin{aretenir}[L'essentiel de la section]
Pour tous $a, b > 0$ et tout entier $n$ :
\[
\ln(ab) = \ln a + \ln b \ ; \quad \ln\left(\frac{1}{a}\right) = -\ln a \ ; \quad \ln\left(\frac{a}{b}\right) = \ln a - \ln b \ ; \quad \ln(a^n) = n\ln a \ ; \quad \ln(\sqrt{a}) = \frac{1}{2}\ln a.
\]
Toutes ces formules découlent de la première, démontrée en montrant que $x \mapsto \ln(ax) - \ln x$ a une dérivée nulle. Et surtout : $\ln(a+b)$ ne se simplifie \textbf{pas}.
\end{aretenir}

Maintenant que nous maîtrisons l'algèbre du logarithme, il est temps d'étudier son comportement aux bornes de son ensemble de définition : que vaut $\ln x$ lorsque $x$ se rapproche de $0$, ou lorsque $x$ devient immense ? C'est l'objet de la section suivante, consacrée aux limites de la fonction $\ln$.

\coursec{Limites de la fonction ln}

Dans la section précédente, nous avons établi les propriétés algébriques de la fonction logarithme népérien : elle transforme les produits en sommes, les quotients en différences et les puissances en produits. Ces propriétés vont maintenant devenir nos meilleurs outils pour répondre à une question essentielle : \emph{comment se comporte $\ln x$ lorsque $x$ devient très grand, ou très proche de $0$ ?}

Cette question n'est pas anodine. Lorsqu'un ingénieur de CAMTEL modélise l'atténuation d'un signal dans une fibre optique, ou lorsqu'un chimiste mesure le pH d'une solution très diluée, ce sont des logarithmes qui interviennent, et il faut savoir prédire leur comportement aux « extrêmes ». Tu vas découvrir dans cette section un phénomène fascinant : la fonction $\ln$ grandit \emph{sans limite}, mais avec une lenteur déconcertante.

\courssub{Limite de $\ln x$ en $+\infty$}

\textbf{Intuition.} Prenons la calculatrice et testons quelques valeurs :
\[
\ln 10 \approx 2{,}30 \quad ; \quad \ln 100 \approx 4{,}61 \quad ; \quad \ln 1000 \approx 6{,}91 \quad ; \quad \ln 10^6 \approx 13{,}82.
\]
Chaque fois que $x$ est multiplié par $10$, $\ln x$ n'augmente que d'environ $2{,}30$ (ce qui n'est autre que $\ln 10$, conformément à la propriété $\ln(10x) = \ln x + \ln 10$). La fonction croît donc de plus en plus lentement... mais croît-elle quand même « jusqu'au bout », ou finit-elle par se bloquer devant un plafond ? La réponse est sans appel : il n'y a \textbf{aucun plafond}.

\begin{propriete}[Limite de $\ln$ en $+\infty$ — démonstration exigible]
\[
\lim_{x \to +\infty} \ln x = +\infty.
\]
Autrement dit : pour tout réel $A$, aussi grand soit-il, il existe un rang au-delà duquel $\ln x > A$.
\end{propriete}

\begin{experience}[Démonstration guidée : $\ln x$ dépasse toute borne]
On veut montrer que $\ln x$ finit par dépasser n'importe quel réel $A$ fixé.
\begin{enumerate}
\item \textbf{Une suite qui file vers l'infini.} Considérons les puissances de $2$ : $2, 4, 8, 16, \ldots, 2^n, \ldots$ Par les propriétés algébriques de la section précédente :
\[
\ln(2^n) = n \ln 2.
\]
Or $\ln 2 > 0$ (car $2 > 1$ et $\ln$ est strictement croissante avec $\ln 1 = 0$). Donc $n \ln 2 \to +\infty$ quand $n \to +\infty$. La suite $\bigl(\ln(2^n)\bigr)$ tend vers $+\infty$.
\item \textbf{La croissance fait le reste.} Soit $A$ un réel quelconque. D'après l'étape 1, il existe un entier $n_0$ tel que $\ln(2^{n_0}) > A$. Posons $x_0 = 2^{n_0}$. Comme $\ln$ est \cle{strictement croissante} sur $]0\,;+\infty[$, pour tout $x > x_0$ :
\[
\ln x > \ln(x_0) = \ln(2^{n_0}) > A.
\]
\item \textbf{Conclusion.} Tout réel $A$ est dépassé à partir d'un certain rang : c'est exactement la définition de $\lim\limits_{x \to +\infty} \ln x = +\infty$. $\blacksquare$
\end{enumerate}
\end{experience}

\begin{exemplebox}[Quelle taille pour dépasser 100 ?]
On cherche à partir de quelle valeur de $x$ on a $\ln x > 100$. En anticipant sur la fonction exponentielle (fonction réciproque de $\ln$, strictement croissante, que nous étudierons bientôt) :
\[
\ln x > 100 \iff x > e^{100}.
\]
Or $e^{100} \approx 2{,}7 \times 10^{43}$ : c'est un nombre astronomique, bien plus grand que le nombre d'habitants de la Terre ! Il faut donc aller \emph{très loin} pour que $\ln x$ dépasse $100$... mais on y arrive toujours. C'est toute la philosophie de cette limite : $\ln$ monte \textbf{lentement mais sûrement} vers $+\infty$.
\end{exemplebox}

\begin{attention}[Le piège du « plafond imaginaire »]
Beaucoup d'élèves, voyant la courbe de $\ln$ s'aplatir, concluent à tort que $\ln x$ tend vers une limite finie (comme $x \mapsto \frac{1}{x}$ tend vers $0$). C'est \textbf{faux} : la courbe s'aplatit mais ne cesse jamais de monter. Une tangente de pente de plus en plus faible n'empêche pas la fonction de croître indéfiniment. Retiens : $\ln x \to +\infty$, mais avec une \cle{lenteur extrême}.
\end{attention}

\courssub{Limite de $\ln x$ en $0^+$}

\textbf{Intuition.} Que se passe-t-il lorsque $x$ s'approche de $0$ par valeurs positives ? Testons :
\[
\ln 0{,}1 \approx -2{,}30 \quad ; \quad \ln 0{,}01 \approx -4{,}61 \quad ; \quad \ln 10^{-6} \approx -13{,}82.
\]
Les valeurs plongent vers le bas. C'est logique : prendre l'inverse transforme « très proche de $0$ » en « très grand », et nous savons déjà ce qui se passe en $+\infty$.

\begin{propriete}[Limite de $\ln$ en $0^+$ — démonstration exigible]
\[
\lim_{\substack{x \to 0 \\ x > 0}} \ln x = -\infty.
\]
Graphiquement, la droite d'équation $x = 0$ (l'axe des ordonnées) est \cle{asymptote verticale} à la courbe de $\ln$.
\end{propriete}

\begin{experience}[Démonstration guidée : le changement de variable $X = \dfrac{1}{x}$]
\begin{enumerate}
\item \textbf{Poser le changement de variable.} Pour $x > 0$, posons $X = \dfrac{1}{x}$. Quand $x \to 0^+$, on a $X \to +\infty$.
\item \textbf{Réécrire $\ln x$ en fonction de $X$.} Par la propriété du logarithme d'un quotient :
\[
\ln x = \ln\left(\frac{1}{X}\right) = \ln 1 - \ln X = -\ln X.
\]
\item \textbf{Composer les limites.} Quand $x \to 0^+$ :
\[
\lim_{x \to 0^+} \ln x = \lim_{X \to +\infty} (-\ln X) = -\infty,
\]
car nous avons démontré que $\lim\limits_{X \to +\infty} \ln X = +\infty$. $\blacksquare$
\end{enumerate}
\end{experience}

\begin{methode}[Le changement de variable, un réflexe à acquérir]
Face à une limite de $\ln$ en un point « inhabituel » :
\begin{enumerate}
\item Identifier où tend la quantité \emph{à l'intérieur} du logarithme.
\item Poser un changement de variable ($X = \frac{1}{x}$, $X = x - a$, etc.) pour se ramener aux deux limites de référence : en $+\infty$ ou en $0^+$.
\item Conclure par composition des limites.
\end{enumerate}
\end{methode}

\begin{popfigure}
\begin{center}
\begin{tikzpicture}[scale=0.9]
\draw[->, thick, popdark] (-0.6,0) -- (6.5,0) node[below left] {$x$};
\draw[->, thick, popdark] (0,-3.2) -- (0,3.2) node[below left] {$y$};
\draw[popblue, very thick, domain=0.06:6.2, samples=200] plot (\x,{ln(\x)});
\draw[popmuted, dashed, thick] (0,-3.2) -- (0,3.2);
\node[popmuted, right] at (0.15,2.6) {asymptote $x = 0$};
\node[popblue, right] at (4.6,1.35) {$y = \ln x$};
\filldraw[popink] (1,0) circle (1.5pt) node[below right] {\small $1$};
\filldraw[popink] (2.718,1) circle (1.5pt) node[above right] {\small $e$};
\draw[popmuted, dashed] (2.718,0) node[below] {\small $e$} -- (2.718,1) -- (0,1) node[left] {\small $1$};
\node[popblue, align=center] at (1.1,-2.5) {\small plongée vers $-\infty$\\ \small quand $x \to 0^+$};
\node[popblue, align=center] at (5.3,2.4) {\small montée lente\\ \small vers $+\infty$};
\end{tikzpicture}
\end{center}
\legende{Courbe de $\ln$ : asymptote verticale d'équation $x = 0$ en $0^+$, et branche infinie qui monte sans borne (mais de plus en plus lentement) en $+\infty$.}
\end{popfigure}

\courssub{Croissances comparées : la lenteur de $\ln$ face aux puissances}

Nous savons que $\ln x \to +\infty$ et que $x \to +\infty$. Mais lequel « gagne » la course ? La réponse, fondamentale pour lever les indéterminations, est que \textbf{toute puissance de $x$ écrase $\ln x$}.

\begin{propriete}[Croissances comparées — résultats admis]
Ces résultats sont \textbf{admis} à ce stade (leur démonstration deviendra accessible après l'étude de la fonction exponentielle) mais leur \emph{utilisation} est exigible au Baccalauréat.
\begin{enumerate}
\item En $+\infty$, la puissance l'emporte sur le logarithme :
\[
\lim_{x \to +\infty} \frac{\ln x}{x} = 0 \qquad \text{et, plus généralement, pour tout } n \in \mathbb{N}^* : \quad \lim_{x \to +\infty} \frac{\ln x}{x^n} = 0.
\]
\item En $0^+$, la puissance l'emporte encore :
\[
\lim_{\substack{x \to 0 \\ x > 0}} x \ln x = 0 \qquad \text{et, plus généralement : } \quad \lim_{\substack{x \to 0 \\ x > 0}} x^n \ln x = 0 \quad (n \in \mathbb{N}^*).
\]
\end{enumerate}
\end{propriete}

\begin{aretenir}[La hiérarchie des infinis]
En $+\infty$, on retiendra l'ordre de « force » suivant :
\[
\ln x \;\ll\; x \;\ll\; x^2 \;\ll\; \cdots \;\ll\; e^x.
\]
Concrètement : dans une somme ou un produit où $\ln x$ affronte une puissance de $x$, c'est \textbf{toujours la puissance qui impose le comportement}. On dit que la courbe de $\ln$ admet en $+\infty$ une \cle{branche parabolique} de direction $(Ox)$ : elle monte vers $+\infty$ mais « moins vite que toute droite ».
\end{aretenir}

\begin{popfigure}
\begin{center}
\begin{tikzpicture}
\begin{axis}[
  width=\linewidth, height=6.5cm,
  xmin=0, xmax=30, ymin=0, ymax=6,
  axis lines=left,
  xlabel={$x$}, ylabel={$y$},
  legend style={at={(0.97,0.35)}, anchor=east, font=\small, draw=popline, fill=popcream},
  tick label style={font=\small},
  label style={font=\small}
]
\addplot[popcurve, domain=0.05:30, samples=200] {ln(x)};
\addlegendentry{$y = \ln x$}
\addplot[poporange, very thick, domain=0:30, samples=100] {sqrt(max(0,x))};
\addlegendentry{$y = \sqrt{x}$}
\addplot[popgreen, very thick, domain=0:6, samples=50] {x};
\addlegendentry{$y = x$}
\end{axis}
\end{tikzpicture}
\end{center}
\legende{Comparaison de $\ln x$, $\sqrt{x}$ et $x$ : dès que $x$ grandit, $\ln x$ est distancé par toutes les puissances, même par $\sqrt{x}$. Voilà l'illustration visuelle de $\lim\limits_{x\to+\infty}\dfrac{\ln x}{x} = 0$.}
\end{popfigure}

\begin{exemplebox}[Exemple chiffré : $\lim\limits_{x \to +\infty}(\ln x - x)$]
Il s'agit d'une forme indéterminée « $+\infty - \infty$ ». On \textbf{factorise par le terme prépondérant}, ici $x$ :
\[
\ln x - x = x\left(\frac{\ln x}{x} - 1\right).
\]
Or, par croissances comparées, $\lim\limits_{x \to +\infty} \dfrac{\ln x}{x} = 0$, donc la parenthèse tend vers $0 - 1 = -1$, et $x \to +\infty$. Par produit :
\[
\lim_{x \to +\infty}(\ln x - x) = -\infty.
\]
Vérification numérique : pour $x = 10^6$, $\ln x - x \approx 13{,}8 - 1\,000\,000 \approx -999\,986$ : la puissance écrase effectivement le logarithme.
\end{exemplebox}

\courssub{Une limite de référence : le taux d'accroissement en $1$}

\begin{propriete}[Limite de référence en $0$]
\[
\lim_{h \to 0} \frac{\ln(1+h)}{h} = 1.
\]
\end{propriete}

\begin{experience}[Pourquoi cette limite vaut 1 ?]
Reconnaissons un \cle{taux d'accroissement}. La fonction $\ln$ est dérivable en $1$ (nous le démontrerons dans la section 4) avec $\ln'(x) = \frac{1}{x}$, donc $\ln'(1) = 1$. Comme $\ln 1 = 0$ :
\[
\frac{\ln(1+h)}{h} = \frac{\ln(1+h) - \ln 1}{h} \xrightarrow[h \to 0]{} \ln'(1) = 1.
\]
Cette limite est donc le \textbf{nombre dérivé de $\ln$ en $1$}. Graphiquement, la tangente à la courbe de $\ln$ au point $(1\,;0)$ a pour coefficient directeur $1$ : au voisinage de $1$, on a l'approximation $\ln(1+h) \approx h$.
\end{experience}

\begin{exemplebox}[Application : $\lim\limits_{x \to 0} \dfrac{\ln(1+3x)}{x}$]
On fait apparaître la forme de référence en écrivant :
\[
\frac{\ln(1+3x)}{x} = 3 \times \frac{\ln(1+3x)}{3x}.
\]
Posons $h = 3x$ : quand $x \to 0$, $h \to 0$, donc $\dfrac{\ln(1+h)}{h} \to 1$. Conclusion :
\[
\lim_{x \to 0} \frac{\ln(1+3x)}{x} = 3 \times 1 = 3.
\]
\end{exemplebox}

\courssub{Limites de fonctions composées avec $\ln$}

\begin{propriete}[Composition avec $\ln$]
Soit $u$ une fonction définie au voisinage d'un point $a$ (fini ou infini), à valeurs strictement positives.
\begin{center}
\begin{tabularx}{\linewidth}{|l|X|}
\hline
\rowcolor{popblueL} \textbf{Si $u(x)$ tend vers...} & \textbf{Alors $\ln\bigl(u(x)\bigr)$ tend vers...} \\
\hline
$+\infty$ & $+\infty$ \\
\hline
$0$ par valeurs positives ($0^+$) & $-\infty$ \\
\hline
un réel $\ell > 0$ & $\ln \ell$ (continuité de $\ln$) \\
\hline
\end{tabularx}
\end{center}
\end{propriete}

\begin{exoresolu}[Trois limites composées]
Déterminer les limites suivantes :
\[
\text{a) } \lim_{x \to +\infty} \ln\left(\frac{x+1}{x-2}\right) \qquad \text{b) } \lim_{\substack{x \to 3 \\ x > 3}} \ln(x^2 - 9) \qquad \text{c) } \lim_{x \to +\infty} \ln(x^2 + 1).
\]
\tcblower
\textbf{a)} Posons $u(x) = \dfrac{x+1}{x-2}$. Quand $x \to +\infty$, $u(x) = \dfrac{1+\frac{1}{x}}{1-\frac{2}{x}} \to 1$, et $u(x) > 0$ pour $x > 2$. Par continuité de $\ln$ en $1$ :
\[
\lim_{x \to +\infty} \ln\left(\frac{x+1}{x-2}\right) = \ln 1 = 0.
\]
\textbf{b)} Posons $u(x) = x^2 - 9 = (x-3)(x+3)$. Quand $x \to 3^+$, $u(x) \to 0$ avec $u(x) > 0$ (car $x > 3$). On est dans le cas « $u \to 0^+$ » :
\[
\lim_{\substack{x \to 3 \\ x > 3}} \ln(x^2 - 9) = -\infty.
\]
\textbf{c)} Posons $u(x) = x^2 + 1 \to +\infty$ quand $x \to +\infty$. On est dans le cas « $u \to +\infty$ » :
\[
\lim_{x \to +\infty} \ln(x^2 + 1) = +\infty.
\]
\end{exoresolu}

\begin{attention}[Toujours vérifier le signe de l'intérieur]
Avant de composer, assure-toi que $u(x) > 0$ au voisinage du point considéré : $\ln$ n'est définie que sur $]0\,;+\infty[$. Écrire $\lim\limits_{x \to 3} \ln(x^2 - 9)$ sans préciser $x > 3$ est une erreur : pour $x < 3$ proche de $3$, $x^2 - 9 < 0$ et l'expression n'existe pas !
\end{attention}

\courssub{Techniques de levée d'indétermination}

Face à une forme indéterminée impliquant $\ln$ (types « $\infty - \infty$ », « $\frac{\infty}{\infty}$ », « $0 \times \infty$ »), deux stratégies dominent.

\begin{methode}[Lever une indétermination avec $\ln$]
\begin{enumerate}
\item \textbf{Factoriser par le terme prépondérant.} Dans une somme, on met en facteur le terme le plus « fort » selon la hiérarchie $\ln x \ll x \ll x^n \ll e^x$, puis on utilise les croissances comparées.
\item \textbf{Poser $X = \ln x$ (ou $X = \frac{1}{x}$).} Ce changement de variable transforme l'expression en une forme polynomiale ou exponentielle plus facile à traiter. Attention à bien convertir : $x = e^X$, et quand $x \to +\infty$, $X \to +\infty$.
\item \textbf{Faire apparaître une limite de référence} : $\dfrac{\ln x}{x}$, $x\ln x$ ou $\dfrac{\ln(1+h)}{h}$.
\end{enumerate}
\end{methode}

\begin{exoresolu}[Deux indéterminations classiques]
Déterminer :
\[
\text{a) } \lim_{x \to +\infty} \frac{\ln x + 2}{x} \qquad \text{b) } \lim_{x \to +\infty} \frac{(\ln x)^2}{x}.
\]
\tcblower
\textbf{a)} On sépare puis on applique les croissances comparées :
\[
\frac{\ln x + 2}{x} = \frac{\ln x}{x} + \frac{2}{x} \xrightarrow[x \to +\infty]{} 0 + 0 = 0.
\]
\textbf{b)} Forme « $\frac{\infty}{\infty}$ ». Posons $X = \ln x$, donc $x = e^X$ avec $X \to +\infty$ quand $x \to +\infty$ :
\[
\frac{(\ln x)^2}{x} = \frac{X^2}{e^X}.
\]
Or la croissance comparée exponentielle/puissance (que nous détaillerons dans la leçon sur l'exponentielle) donne $\lim\limits_{X \to +\infty} \dfrac{e^X}{X^2} = +\infty$, donc par passage à l'inverse :
\[
\lim_{x \to +\infty} \frac{(\ln x)^2}{x} = 0.
\]
Conclusion morale : \textbf{même élevé au carré (ou à toute puissance), $\ln x$ reste écrasé par $x$}.
\end{exoresolu}

\begin{exoresolu}[Le changement de variable $X = \sqrt{x}$ en action]
Déterminer $\lim\limits_{x \to +\infty} \bigl(\ln x - 3\sqrt{x}\bigr)$.
\tcblower
Forme « $+\infty - \infty$ ». Factorisons par le terme prépondérant $\sqrt{x}$ :
\[
\ln x - 3\sqrt{x} = \sqrt{x}\left(\frac{\ln x}{\sqrt{x}} - 3\right).
\]
Pour la parenthèse, posons $X = \sqrt{x}$ (donc $x = X^2$, $X \to +\infty$) :
\[
\frac{\ln x}{\sqrt{x}} = \frac{\ln(X^2)}{X} = \frac{2\ln X}{X} \xrightarrow[X \to +\infty]{} 2 \times 0 = 0.
\]
La parenthèse tend donc vers $-3$ et $\sqrt{x} \to +\infty$. Par produit :
\[
\lim_{x \to +\infty} \bigl(\ln x - 3\sqrt{x}\bigr) = -\infty.
\]
Encore une victoire de la puissance sur le logarithme !
\end{exoresolu}

\begin{savaistu}[Pourquoi le pH et les décibels utilisent-ils des logarithmes ?]
L'échelle des pH en chimie va de $0$ à $14$, mais la concentration en ions $H_3O^+$ qu'elle mesure varie de $1$ à $10^{-14}$ mol/L : un rapport de $1$ à $100\,000\,000\,000\,000$ ! De même, l'échelle des décibels comprime l'intensité sonore, et l'échelle de Richter comprime l'énergie des séismes. C'est précisément la \textbf{lenteur de croissance de $\ln$} — le fait que $\ln x$ « écrase » les grandes valeurs — qui rend ces échelles pratiques. Les mathématiciens disent que le logarithme « comprime les ordres de grandeur » : un outil indispensable dès que les phénomènes s'étalent sur des rapports immenses, comme en sismologie, un domaine suivi de près au Cameroun avec la surveillance du mont Cameroun, volcan actif.
\end{savaistu}

\begin{aretenir}[Bilan des limites à connaître par cœur]
\begin{center}
\begin{tabularx}{\linewidth}{|X|X|}
\hline
\rowcolor{popblueL} \textbf{Limite} & \textbf{Résultat} \\
\hline
$\lim\limits_{x \to +\infty} \ln x$ & $+\infty$ (montée lente, sans borne) \\
\hline
$\lim\limits_{x \to 0^+} \ln x$ & $-\infty$ (asymptote verticale $x = 0$) \\
\hline
$\lim\limits_{x \to +\infty} \dfrac{\ln x}{x}$ & $0$ (la puissance l'emporte) \\
\hline
$\lim\limits_{x \to 0^+} x \ln x$ & $0$ (la puissance l'emporte) \\
\hline
$\lim\limits_{h \to 0} \dfrac{\ln(1+h)}{h}$ & $1$ (nombre dérivé en $1$) \\
\hline
\end{tabularx}
\end{center}
Face à une indétermination : \textbf{factoriser par le terme dominant} ou \textbf{poser $X = \ln x$}, puis se ramener à ce tableau.
\end{aretenir}

\begin{exercice}[À toi de jouer]
Déterminer les limites suivantes, en justifiant chaque étape :
\[
\text{a) } \lim_{x \to +\infty} \frac{x - \ln x}{x} \qquad
\text{b) } \lim_{\substack{x \to 0 \\ x > 0}} \bigl(x + x\ln x\bigr) \qquad
\text{c) } \lim_{x \to +\infty} \ln\left(\frac{2x^2+1}{x^2+3}\right) \qquad
\text{d) } \lim_{x \to 0} \frac{\ln(1 - 2x)}{x}.
\]
\end{exercice}

\begin{corrige}[Exercice « À toi de jouer »]
\textbf{a)} $\dfrac{x - \ln x}{x} = 1 - \dfrac{\ln x}{x} \to 1 - 0 = 1$.

\textbf{b)} $x + x\ln x$ : on a $x \to 0$ et $x\ln x \to 0$ (croissance comparée en $0^+$), donc la somme tend vers $0$.

\textbf{c)} $\dfrac{2x^2+1}{x^2+3} = \dfrac{2+\frac{1}{x^2}}{1+\frac{3}{x^2}} \to 2$ quand $x \to +\infty$, et $2 > 0$, donc par continuité de $\ln$ en $2$ : la limite vaut $\ln 2$.

\textbf{d)} $\dfrac{\ln(1-2x)}{x} = -2 \times \dfrac{\ln(1-2x)}{-2x}$. Avec $h = -2x \to 0$, le quotient tend vers $1$, donc la limite vaut $-2$.
\end{corrige}

Nous disposons désormais d'une vision complète du comportement de $\ln$ aux bornes de son ensemble de définition. Il reste à étudier ce qui se passe \emph{entre} ces bornes : la dérivabilité, le sens de variation et le tracé précis de la courbe — ce sera l'objet de la section suivante.

\coursec{Dérivation, variations et courbe représentative}

Dans la section précédente, nous avons établi les limites de la fonction logarithme népérien : $\displaystyle\lim_{x\to 0^+}\ln x = -\infty$ et $\displaystyle\lim_{x\to +\infty}\ln x = +\infty$. Nous savons donc comment la fonction se comporte aux bornes de son domaine. Mais que se passe-t-il \emph{entre} ces bornes ? La fonction monte-t-elle ? Vite ou lentement ? Sa courbe est-elle bombée ou creuse ? Pour répondre à toutes ces questions, l'outil roi est la \cle{dérivation}. C'est d'ailleurs ainsi que la fonction $\ln$ a été \emph{construite} : rappelle-toi, $\ln$ est la primitive de $x\mapsto \dfrac{1}{x}$ qui s'annule en $1$. La dérivée de $\ln$ est donc connue avant même la fonction elle-même !

\courssub{Dérivée de la fonction ln et stricte croissance}

\textbf{Intuition.} Dire que $\ln$ est une primitive de $x\mapsto \dfrac{1}{x}$, c'est dire que la pente de la tangente à la courbe de $\ln$ au point d'abscisse $x$ vaut exactement $\dfrac{1}{x}$. Or pour tout $x>0$, on a $\dfrac{1}{x}>0$ : la tangente « monte » toujours. La fonction $\ln$ est donc constamment en train de croître. De plus, quand $x$ devient grand, $\dfrac{1}{x}$ devient petit : la courbe monte de moins en moins vite, sans jamais s'arrêter de monter.

\begin{propriete}[Dérivée de la fonction ln]
La fonction $\ln$ est dérivable sur $]0\,;\,+\infty[$ et, pour tout réel $x>0$ :
\[ (\ln)'(x) = \frac{1}{x}. \]
Comme $\dfrac{1}{x}>0$ pour tout $x>0$, la fonction $\ln$ est \cle{strictement croissante} sur $]0\,;\,+\infty[$.
\end{propriete}

\begin{experience}[Démonstration guidée — dérivabilité de ln]
Cette démonstration est classique et formatrice ; sais-tu la refaire seul ?
\begin{enumerate}
\item Par définition, $\ln$ est \textbf{la} primitive de $x\mapsto \dfrac{1}{x}$ sur $]0\,;\,+\infty[$ qui s'annule en $1$ : $\ln x = \displaystyle\int_1^x \frac{1}{t}\,\mathrm{d}t$.
\item La fonction $t\mapsto \dfrac{1}{t}$ est continue sur $]0\,;\,+\infty[$. D'après le théorème fondamental de l'analyse, la fonction $x\mapsto \displaystyle\int_1^x \frac{1}{t}\,\mathrm{d}t$ est dérivable sur $]0\,;\,+\infty[$, de dérivée $x\mapsto \dfrac{1}{x}$.
\item Donc $(\ln)'(x) = \dfrac{1}{x}$ pour tout $x>0$. Comme $\dfrac{1}{x}>0$ sur $]0\,;\,+\infty[$, $\ln$ est strictement croissante sur cet intervalle. $\square$
\end{enumerate}
\end{experience}

\begin{propriete}[Conséquences de la stricte croissance]
Pour tous réels $a>0$ et $b>0$ :
\[ \ln a = \ln b \iff a = b \qquad\text{et}\qquad \ln a < \ln b \iff a < b. \]
Autrement dit, $\ln$ \emph{conserve} l'ordre : comparer deux logarithmes revient à comparer leurs arguments. Ces équivalences sont la clé de la résolution des équations et inéquations avec $\ln$ (section 6).
\end{propriete}

\begin{exemplebox}[Comparer sans calculatrice]
Comparer $\ln 7$ et $\ln 9$, puis résoudre $\ln x \geq \ln 5$.

\medskip
Comme $7<9$ et que $\ln$ est strictement croissante : $\ln 7 < \ln 9$.

Pour l'inéquation $\ln x \geq \ln 5$ : la condition d'existence est $x>0$, et par stricte croissance, $\ln x \geq \ln 5 \iff x \geq 5$. L'ensemble des solutions est $[5\,;\,+\infty[$.
\end{exemplebox}

\courssub{Dérivée de la composée : la formule $(\ln u)' = u'/u$}

En pratique, au Baccalauréat, on dérive rarement $\ln x$ tout seul : on dérive $\ln\bigl(u(x)\bigr)$, par exemple $\ln(2x+3)$ ou $\ln(x^2+1)$. C'est une simple application de la formule de dérivation des fonctions composées.

\begin{propriete}[Dérivée de $\ln u$]
Soit $u$ une fonction dérivable et \textbf{strictement positive} sur un intervalle $I$. Alors la fonction $f : x\mapsto \ln\bigl(u(x)\bigr)$ est dérivable sur $I$ et :
\[ f'(x) = \bigl(\ln u\bigr)'(x) = \frac{u'(x)}{u(x)}. \]
\end{propriete}

\begin{experience}[Démonstration guidée — pourquoi $u'/u$ ?]
\begin{enumerate}
\item Rappel : si $v$ est dérivable sur $I$ et $g$ dérivable sur $v(I)$, alors $(g\circ v)' = v'\times (g'\circ v)$.
\item Ici $g = \ln$, donc $g'(t) = \dfrac{1}{t}$, et $v = u$. La condition « $\ln$ dérivable sur $u(I)$ » impose $u(x)>0$ sur $I$.
\item On obtient : $\bigl(\ln u\bigr)'(x) = u'(x)\times \dfrac{1}{u(x)} = \dfrac{u'(x)}{u(x)}$. $\square$
\end{enumerate}
\end{experience}

\begin{exemplebox}[Deux dérivations types]
\textbf{a)} Soit $f(x) = \ln(2x+3)$. On pose $u(x)=2x+3$, strictement positive pour $x>-\dfrac{3}{2}$, et $u'(x)=2$. Donc, pour tout $x>-\dfrac{3}{2}$ :
\[ f'(x) = \frac{2}{2x+3}. \]
\textbf{b)} Soit $g(x) = \ln(x^2+1)$. Ici $u(x)=x^2+1>0$ pour \emph{tout} réel $x$, donc $g$ est définie et dérivable sur $\mathbb{R}$ tout entier, et $u'(x)=2x$ :
\[ g'(x) = \frac{2x}{x^2+1}. \]
Le signe de $g'$ est celui de $2x$ : $g$ est décroissante sur $]-\infty\,;\,0]$ et croissante sur $[0\,;\,+\infty[$.
\end{exemplebox}

\begin{attention}[La dérivée ne donne pas le domaine !]
Erreur très fréquente au Bac : pour étudier $f(x)=\ln\bigl(u(x)\bigr)$, des élèves dérivent d'abord et « déduisent » le domaine de la fraction $\dfrac{u'}{u}$. C'est un raisonnement faux et dangereux.
\begin{itemize}
\item Le domaine de $f$ s'obtient en résolvant l'inéquation $u(x)>0$ (\textbf{strictement}).
\item Exemple piège : $f(x)=\ln(x^2-1)$. Le domaine est $x^2-1>0$, c'est-à-dire $]-\infty\,;\,-1[\,\cup\,]1\,;\,+\infty[$, alors que $f'(x)=\dfrac{2x}{x^2-1}$ laisse croire à tort qu'on pourrait évaluer $f$ « presque partout ». Non : $f(0)=\ln(-1)$ n'existe pas !
\end{itemize}
Retiens : \textbf{domaine d'abord, dérivation ensuite}.
\end{attention}

\begin{methode}[Étudier une fonction contenant $\ln u$]
\begin{enumerate}
\item \textbf{Domaine} : résoudre $u(x)>0$ (inéquation stricte). Écrire $D_f$ proprement.
\item \textbf{Limites} aux bornes de $D_f$ : utiliser les limites usuelles de $\ln$ et les théorèmes sur les composées (si $u(x)\to 0^+$, alors $\ln u(x)\to -\infty$ ; si $u(x)\to +\infty$, alors $\ln u(x)\to +\infty$).
\item \textbf{Dérivée} : $f'(x)=\dfrac{u'(x)}{u(x)}$ (ou dériver une expression plus complexe avec les règles usuelles).
\item \textbf{Signe de $f'$} : comme $u(x)>0$ sur $D_f$, le signe de $\dfrac{u'}{u}$ est le signe de $u'(x)$. Dresser le tableau de variations.
\item \textbf{Courbe} : points remarquables, tangentes, asymptotes éventuelles.
\end{enumerate}
\end{methode}

\courssub{Tableau de variations complet et courbe représentative de $\ln$}

Réunissons tout ce que nous savons sur $\ln$ : dérivée $\dfrac{1}{x}>0$, limites $-\infty$ en $0^+$ et $+\infty$ en $+\infty$, valeurs remarquables $\ln 1 = 0$ et $\ln e = 1$.

\begin{center}
\begin{tabular}{|c|ccc|}
\hline
\rowcolor{popblueL}
$x$ & $0$ & \hspace{2.5cm} & $+\infty$ \\
\hline
$(\ln)'(x)=\dfrac{1}{x}$ & & $+$ & \\
\hline
$\ln x$ & $-\infty$ & $\nearrow$ & $+\infty$ \\
\hline
\end{tabular}
\end{center}

\begin{aretenir}[Éléments remarquables de la courbe de $\ln$]
\begin{itemize}
\item La courbe passe par les points $(1\,;\,0)$ et $(e\,;\,1)$, car $\ln 1 = 0$ et $\ln e = 1$.
\item Comme $\displaystyle\lim_{x\to 0^+}\ln x = -\infty$, la droite d'équation $x=0$ (l'axe des ordonnées) est \cle{asymptote verticale} à la courbe.
\item Comme $\displaystyle\lim_{x\to +\infty}\ln x = +\infty$, il n'y a \textbf{pas d'asymptote horizontale}.
\item Le coefficient directeur de la tangente en $x=1$ vaut $(\ln)'(1)=\dfrac{1}{1}=1$. La tangente au point $(1\,;\,0)$ a donc pour équation :
\[ y = (x-1)\times 1 + 0 \quad\text{c'est-à-dire}\quad y = x-1. \]
\end{itemize}
\end{aretenir}

\begin{popfigure}
\begin{center}
\begin{tikzpicture}
\begin{axis}[
    width=0.85\linewidth, height=7.5cm,
    axis lines=middle,
    xlabel={$x$}, ylabel={$y$},
    xmin=-0.5, xmax=7.5, ymin=-3.2, ymax=3.2,
    xtick={1,2.718,5}, xticklabels={$1$,$e$,$5$},
    ytick={-2,-1,1,2},
    grid=both, grid style={popline!40},
    legend style={at={(0.98,0.35)},anchor=east,draw=popline,fill=popcream}
]
\addplot[popcurve,domain=0.045:7.3,samples=300]{ln(x)};
\addlegendentry{$y=\ln x$}
\addplot[poporange,thick,domain=-0.4:4.2,samples=2]{x-1};
\addlegendentry{tangente $y=x-1$}
\addplot[popmuted,dashed,thick] coordinates {(0,-3.2) (0,3.2)};
\addplot[only marks,mark=*,popink,mark size=1.8pt] coordinates {(1,0) (2.718,1)};
\node[popink,anchor=south west] at (axis cs:1,0) {\small $(1\,;\,0)$};
\node[popink,anchor=south west] at (axis cs:2.718,1) {\small $(e\,;\,1)$};
\node[popmuted,anchor=west] at (axis cs:0.08,2.6) {\small $x=0$};
\end{axis}
\end{tikzpicture}
\end{center}
\legende{Courbe de $y=\ln x$, sa tangente $y=x-1$ au point $(1\,;\,0)$ et l'asymptote verticale $x=0$. La courbe reste sous sa tangente : c'est la concavité.}
\end{popfigure}

\textbf{Complément utile au Bac : la concavité.} La dérivée seconde se calcule facilement : pour tout $x>0$,
\[ (\ln)''(x) = \left(\frac{1}{x}\right)' = -\frac{1}{x^2} < 0. \]
La fonction $\ln$ est donc \cle{concave} sur $]0\,;\,+\infty[$ : sa courbe est toujours située \emph{en dessous} de chacune de ses tangentes. Ce n'est pas un savoir exigible du programme, mais c'est un atout précieux pour justifier rigoureusement certaines inégalités.

\begin{savaistu}[L'inégalité $\ln x \leq x - 1$ : un classique des sujets de Bac]
Puisque $\ln$ est concave, sa courbe est sous sa tangente en $(1\,;\,0)$ : pour tout $x>0$,
\[ \ln x \leq x - 1, \]
avec égalité si et seulement si $x=1$. Cette inégalité, d'apparence anodine, est extraordinairement puissante : elle intervient dans les sujets de Baccalauréat (études de fonctions, suites, probabilités) et dans toutes les mathématiques modernes, de la théorie de l'information à l'économie. Retiens-la, et sache la redémontrer en étudiant les variations de $h(x)=x-1-\ln x$ : $h'(x)=1-\dfrac{1}{x}$, donc $h$ admet un minimum en $x=1$ valant $h(1)=0$.
\end{savaistu}

\courssub{Étude complète d'une fonction type : $f(x) = \ln x / x$}

Appliquons la méthode à une fonction très classique des sujets de Baccalauréat.

\begin{exoresolu}[Étude complète de $f(x) = \ln x / x$]
Soit $f$ la fonction définie par $f(x)=\dfrac{\ln x}{x}$.
\begin{enumerate}
\item Déterminer le domaine de définition de $f$.
\item Calculer les limites de $f$ aux bornes de son domaine et interpréter graphiquement.
\item Calculer $f'(x)$, étudier son signe et dresser le tableau de variations de $f$.
\item Tracer l'allure de la courbe représentative de $f$.
\end{enumerate}
\tcblower
\textbf{1. Domaine.} Il faut $\ln x$ défini, c'est-à-dire $x>0$, et le dénominateur non nul, c'est-à-dire $x\neq 0$. Donc $D_f = ]0\,;\,+\infty[$.

\textbf{2. Limites.}
\begin{itemize}
\item En $0^+$ : $\ln x \to -\infty$ et $\dfrac{1}{x}\to +\infty$, donc par produit $\displaystyle\lim_{x\to 0^+} f(x) = -\infty$. La droite $x=0$ est asymptote verticale.
\item En $+\infty$ : c'est une forme indéterminée « $\dfrac{\infty}{\infty}$ », levée par la limite de croissance comparée vue en section 3 : $\displaystyle\lim_{x\to +\infty}\frac{\ln x}{x} = 0$. L'axe des abscisses ($y=0$) est asymptote horizontale en $+\infty$.
\end{itemize}

\textbf{3. Dérivée et variations.} $f$ est un quotient $\dfrac{u}{v}$ avec $u(x)=\ln x$, $u'(x)=\dfrac{1}{x}$, $v(x)=x$, $v'(x)=1$ :
\[ f'(x) = \frac{\dfrac{1}{x}\times x - \ln x \times 1}{x^2} = \frac{1-\ln x}{x^2}. \]
Comme $x^2>0$ sur $D_f$, le signe de $f'$ est celui de $1-\ln x$ :
\[ 1-\ln x > 0 \iff \ln x < 1 \iff x < e \quad(\text{stricte croissance de }\ln). \]
Donc $f$ est strictement croissante sur $]0\,;\,e]$ et strictement décroissante sur $[e\,;\,+\infty[$. Elle admet un \textbf{maximum} en $x=e$ :
\[ f(e) = \frac{\ln e}{e} = \frac{1}{e} \approx 0{,}368. \]

\begin{center}
\begin{tabular}{|c|ccccc|}
\hline
\rowcolor{popblueL}
$x$ & $0$ & & $e$ & & $+\infty$ \\
\hline
$1-\ln x$ & & $+$ & $0$ & $-$ & \\
\hline
$f(x)$ & $-\infty$ & $\nearrow$ & $\dfrac{1}{e}$ & $\searrow$ & $0$ \\
\hline
\end{tabular}
\end{center}

\textbf{4. Courbe.} Voir la figure ci-dessous : la courbe part de $-\infty$ le long de l'axe des ordonnées, coupe l'axe des abscisses en $x=1$ (car $\ln 1 = 0$), culmine au point $\left(e\,;\,\dfrac{1}{e}\right)$, puis redescend en s'approchant de l'axe des abscisses.
\end{exoresolu}

\begin{popfigure}
\begin{center}
\begin{tikzpicture}
\begin{axis}[
    width=0.85\linewidth, height=7cm,
    axis lines=middle,
    xlabel={$x$}, ylabel={$y$},
    xmin=-0.5, xmax=12, ymin=-2.2, ymax=1,
    xtick={1,2.718,6,10}, xticklabels={$1$,$e$,$6$,$10$},
    ytick={-2,-1,0.368}, yticklabels={$-2$,$-1$,$\frac{1}{e}$},
    grid=both, grid style={popline!40}
]
\addplot[popcurve,domain=0.1:11.8,samples=300]{ln(x)/x};
\addplot[popmuted,dashed,thick] coordinates {(0,-2.2) (0,1)};
\addplot[poporange,dashed,thick] coordinates {(2.718,0) (2.718,0.368)};
\addplot[poporange,dashed,thick] coordinates {(0,0.368) (2.718,0.368)};
\addplot[only marks,mark=*,popink,mark size=1.8pt] coordinates {(2.718,0.368) (1,0)};
\node[popink,anchor=south] at (axis cs:2.718,0.42) {\small maximum $\left(e\,;\,\frac{1}{e}\right)$};
\end{axis}
\end{tikzpicture}
\end{center}
\legende{Courbe de $f(x)=\dfrac{\ln x}{x}$ : maximum $\dfrac{1}{e}$ atteint en $x=e$, asymptote verticale $x=0$ et asymptote horizontale $y=0$ en $+\infty$.}
\end{popfigure}

\begin{exoresolu}[Second entraînement : $f(x) = x - \ln x$]
Soit $f(x)=x-\ln x$. Étudier ses variations et en déduire le signe de $f(x)$ sur $]0\,;\,+\infty[$.
\tcblower
\textbf{Domaine} : $D_f=]0\,;\,+\infty[$ à cause de $\ln x$.

\textbf{Limites} : en $0^+$, $-\ln x \to +\infty$ donc $f(x)\to +\infty$. En $+\infty$, forme indéterminée : on factorise, $f(x)=x\left(1-\dfrac{\ln x}{x}\right)$. Or $\dfrac{\ln x}{x}\to 0$, donc $f(x)\to +\infty$.

\textbf{Dérivée} : $f'(x)=1-\dfrac{1}{x}=\dfrac{x-1}{x}$. Comme $x>0$, le signe de $f'$ est celui de $x-1$ : $f$ est décroissante sur $]0\,;\,1]$ et croissante sur $[1\,;\,+\infty[$.

\textbf{Minimum} : $f(1)=1-\ln 1 = 1$. Le minimum de $f$ vaut $1>0$, donc pour tout $x>0$ : $f(x)\geq 1 > 0$, c'est-à-dire $x-\ln x > 0$, soit $\ln x < x$. On retrouve ainsi l'inégalité $\ln x \leq x-1$ (puisque $x-\ln x \ge 1 \iff \ln x \le x-1$) par une autre méthode !
\end{exoresolu}

\begin{exercice}[À toi de jouer]
Soit $g$ la fonction définie par $g(x)=\ln(x^2+x+1)$.
\begin{enumerate}
\item Montrer que $g$ est définie sur $\mathbb{R}$ (penser au discriminant).
\item Calculer $g'(x)$ et dresser le tableau de variations de $g$.
\item Déterminer le minimum de $g$.
\end{enumerate}
\end{exercice}

\begin{corrige}[Exercice : $g(x) = \ln(x^2 + x + 1)$]
\textbf{1.} Le discriminant de $x^2+x+1$ est $\Delta = 1-4 = -3 < 0$ : le trinôme ne s'annule pas et est du signe de $a=1>0$, donc $x^2+x+1>0$ pour tout réel $x$. Ainsi $D_g=\mathbb{R}$.

\textbf{2.} Avec $u(x)=x^2+x+1$, $u'(x)=2x+1$ :
\[ g'(x) = \frac{2x+1}{x^2+x+1}. \]
Le dénominateur est strictement positif, donc le signe de $g'$ est celui de $2x+1$ : $g$ est décroissante sur $\left]-\infty\,;\,-\dfrac{1}{2}\right]$ et croissante sur $\left[-\dfrac{1}{2}\,;\,+\infty\right[$.

\textbf{3.} Le minimum est atteint en $x=-\dfrac{1}{2}$ :
\[ g\left(-\tfrac{1}{2}\right) = \ln\left(\tfrac{1}{4}-\tfrac{1}{2}+1\right) = \ln\tfrac{3}{4}. \]
Comme $\dfrac{3}{4}<1$, ce minimum est négatif : $\ln\dfrac{3}{4}\approx -0{,}29$.
\end{corrige}

\begin{aretenir}[L'essentiel de la section]
\begin{itemize}
\item $(\ln)'(x)=\dfrac{1}{x}>0$ sur $]0\,;\,+\infty[$ : $\ln$ est strictement croissante, d'où $\ln a = \ln b \iff a=b$ et $\ln a < \ln b \iff a<b$ (pour $a,b>0$).
\item $\bigl(\ln u\bigr)' = \dfrac{u'}{u}$, valable uniquement là où $u$ est dérivable et \textbf{strictement positive}.
\item La courbe de $\ln$ passe par $(1\,;\,0)$ et $(e\,;\,1)$, admet la tangente $y=x-1$ en $(1\,;\,0)$ et l'asymptote verticale $x=0$.
\item $(\ln)''(x)=-\dfrac{1}{x^2}<0$ : $\ln$ est concave, d'où l'inégalité $\ln x \leq x-1$.
\item Toute étude de fonction avec $\ln$ commence par la résolution de $u(x)>0$ : le domaine d'abord, la dérivée ensuite.
\end{itemize}
\end{aretenir}

\coursec{Primitives de la forme $u'/u$}

Dans la section précédente, nous avons établi que la fonction logarithme népérien est \emph{dérivable} sur $]0\,;+\infty[$ et que, pour tout $x > 0$,
\[ \ln'(x) = \frac{1}{x}. \]
Nous avons même généralisé ce résultat : si $u$ est une fonction dérivable et strictement positive sur un intervalle $I$, alors la fonction composée $\ln u$ est dérivable sur $I$ et
\[ (\ln u)' = \frac{u'}{u}. \]
Renversons maintenant le point de vue. Au lieu de \emph{dériver} $\ln u$, demandons-nous : quelle fonction a pour dérivée $\dfrac{u'}{u}$ ? La réponse tombe sous le sens : c'est $\ln u$. Autrement dit, \textbf{la recherche de primitives est l'opération inverse de la dérivation}, et la formule de dérivation de $\ln u$ lue « de droite à gauche » devient une formule de primitives extrêmement puissante. C'est toute l'ambition de cette section : apprendre à \cle{reconnaître} la forme $\dfrac{u'}{u}$ dans une expression, à \cle{ajuster} les constantes, et à en déduire primitives et intégrales. Ce savoir-faire est l'un des plus rentables du programme de Terminale : il revient chaque année au Baccalauréat, dans les calculs d'aires, les équations différentielles et les problèmes d'évolution.

\courssub{Le résultat fondamental}

\courssub{Intuition : dériver, c'est descendre ; primitiver, c'est remonter}

Imagine un téléphérique qui relie deux quartiers de Yaoundé situés à des altitudes différentes. La dérivation, c'est le trajet descendant : à partir d'une fonction $F$ (en haut), on obtient sa dérivée $f$ (en bas). La recherche de primitives, c'est le trajet remontant : connaissant $f$, on cherche $F$ telle que $F' = f$. Or nous venons de voir que la dérivation envoie $\ln u$ sur $\dfrac{u'}{u}$. Le trajet inverse envoie donc $\dfrac{u'}{u}$ sur $\ln u$. Il ne reste qu'à formaliser proprement, en faisant attention aux conditions de validité.

\begin{definition}[Primitive d'une fonction]
Soit $f$ une fonction définie sur un intervalle $I$. On appelle \cle{primitive} de $f$ sur $I$ toute fonction $F$ dérivable sur $I$ telle que, pour tout $x \in I$ : $F'(x) = f(x)$.
\end{definition}

Rappelons deux faits essentiels : une fonction qui admet une primitive sur $I$ en admet une infinité, et deux primitives d'une même fonction sur un \emph{intervalle} diffèrent d'une constante. C'est pourquoi on écrira souvent $\ln u + C$, où $C$ désigne une constante réelle quelconque.

\begin{propriete}[Théorème fondamental : primitive de $u'/u$]
Soit $u$ une fonction dérivable sur un intervalle $I$.
\begin{itemize}
\item Si $u(x) > 0$ pour tout $x \in I$, alors la fonction $\dfrac{u'}{u}$ admet pour primitives sur $I$ les fonctions
\[ x \longmapsto \ln\bigl(u(x)\bigr) + C, \quad C \in \mathbb{R}. \]
\item Plus généralement, si $u$ \textbf{ne s'annule pas} sur $I$ (donc garde un signe constant, d'après le théorème des valeurs intermédiaires appliqué à la fonction continue $u$), alors $\dfrac{u'}{u}$ admet pour primitives sur $I$ les fonctions
\[ x \longmapsto \ln\bigl|u(x)\bigr| + C, \quad C \in \mathbb{R}. \]
\end{itemize}
\end{propriete}

\begin{experience}[Démonstration guidée (exigible et formatrice)]
Nous allons démontrer les deux cas. L'outil central est la dérivation des fonctions composées : $(v \circ u)' = u' \times (v' \circ u)$.
\begin{enumerate}
\item \textbf{Cas $u > 0$ sur $I$.} Posons $F = \ln u$, c'est-à-dire $F = \ln \circ u$. La fonction $\ln$ est dérivable sur $]0\,;+\infty[$ et $u$ est dérivable sur $I$ à valeurs dans $]0\,;+\infty[$, donc $F$ est dérivable sur $I$ et, pour tout $x \in I$ :
\[ F'(x) = u'(x) \times \ln'\bigl(u(x)\bigr) = u'(x) \times \frac{1}{u(x)} = \frac{u'(x)}{u(x)}. \]
Donc $F = \ln u$ est bien une primitive de $\dfrac{u'}{u}$ sur $I$.
\item \textbf{Cas $u < 0$ sur $I$.} La fonction $\ln u$ n'est alors pas définie. L'idée est de poser $G = \ln(-u)$ : comme $u < 0$, on a $-u > 0$, donc $G$ est bien définie et dérivable sur $I$. La dérivée de $-u$ est $-u'$, donc :
\[ G'(x) = \frac{-u'(x)}{-u(x)} = \frac{u'(x)}{u(x)}. \]
Miracle des signes : $G = \ln(-u)$ est aussi une primitive de $\dfrac{u'}{u}$ !
\item \textbf{Synthèse.} Si $u > 0$, $|u| = u$ et $\ln|u| = \ln u$ ; si $u < 0$, $|u| = -u$ et $\ln|u| = \ln(-u)$. Dans les deux cas, $\ln|u|$ est dérivable sur $I$ et
\[ \bigl(\ln|u|\bigr)' = \frac{u'}{u}. \]
La formule $\ln|u|$ unifie donc les deux situations. $\blacksquare$
\end{enumerate}
\end{experience}

\begin{aretenir}[La formule à connaître par cœur]
Si $u$ est dérivable et \textbf{ne s'annule pas} sur l'intervalle $I$, alors :
\[ \boxed{\;\int \frac{u'(x)}{u(x)}\,\mathrm{d}x = \ln\bigl|u(x)\bigr| + C\;} \]
En particulier, lorsque $u > 0$ sur $I$ (cas le plus fréquent au Bac) : une primitive de $\dfrac{u'}{u}$ est $\ln u$.
\end{aretenir}

\begin{attention}[La condition « $u$ ne s'annule pas » n'est pas négociable]
Trois erreurs reviennent sans cesse dans les copies :
\begin{itemize}
\item \textbf{Oublier de vérifier que $u \neq 0$ sur $I$.} La fonction $\dfrac{u'}{u}$ n'est même pas \emph{définie} là où $u$ s'annule. Par exemple, $\dfrac{2x}{x^2 - 1}$ n'a de sens que sur $]-\infty\,;-1[$, $]-1\,;1[$ ou $]1\,;+\infty[$ : il faut choisir l'intervalle de travail et le dire.
\item \textbf{Oublier la valeur absolue} quand $u$ peut être négatif. Sur $]-\infty\,;-1[$, $x^2 - 1 > 0$ donc $\ln(x^2-1)$ convient ; mais pour $\dfrac{1}{x}$ sur $]-\infty\,;0[$, la primitive est $\ln(-x) = \ln|x|$, et non $\ln x$ qui n'existe pas !
\item \textbf{Croire que la primitive de $\dfrac{1}{u}$ est $\ln|u|$.} C'est faux en général : il faut le $u'$ au numérateur. Par exemple, une primitive de $\dfrac{1}{x^2+1}$ n'est \emph{pas} $\ln(x^2+1)$ (c'est $\arctan x$, hors programme).
\end{itemize}
\end{attention}

\courssub{Reconnaître la forme $u'/u$ : la stratégie du dénominateur}

Tout l'art consiste à repérer, dans une fraction, que le numérateur est la dérivée du dénominateur — éventuellement à un facteur constant près.

\begin{methode}[Reconnaître et traiter la forme $u'/u$]
Face à une fonction $f$ écrite comme une fraction (rationnelle, avec racines, avec $\ln$, etc.) :
\begin{enumerate}
\item \textbf{Identifier le dénominateur} et poser $u(x) = $ dénominateur.
\item \textbf{Calculer $u'(x)$} et le comparer au numérateur.
\item \textbf{Trois cas possibles :}
\begin{itemize}
\item le numérateur est exactement $u'$ : une primitive est $\ln|u|$ ;
\item le numérateur est $k \times u'$ ($k$ constante) : une primitive est $k\ln|u|$ ;
\item le numérateur n'est pas proportionnel à $u'$ : la forme $u'/u$ ne s'applique pas directement (il faudra une autre technique : décomposition, forme $u' \times u^n$, etc.).
\end{itemize}
\item \textbf{Vérifier le signe de $u$} sur l'intervalle considéré pour décider si la valeur absolue peut être enlevée.
\item \textbf{Contrôler le résultat en dérivant} : la dérivée de la primitive trouvée doit redonner $f$. Cette vérification, rapide, est une excellente habitude de copie.
\end{enumerate}
\end{methode}

\begin{exemplebox}[Exemple chiffré fondamental]
Déterminons une primitive de $f(x) = \dfrac{2x}{x^2 + 1}$ sur $\mathbb{R}$.
\begin{enumerate}
\item Le dénominateur est $u(x) = x^2 + 1$. Pour tout réel $x$, $x^2 + 1 \geq 1 > 0$ : $u$ est strictement positive sur $\mathbb{R}$, pas de problème de domaine.
\item $u'(x) = 2x$ : c'est \emph{exactement} le numérateur. La forme $u'/u$ s'applique directement.
\item Conclusion : une primitive de $f$ sur $\mathbb{R}$ est
\[ F(x) = \ln(x^2 + 1). \]
\item Vérification : $F'(x) = \dfrac{2x}{x^2+1} = f(x)$. $\checkmark$
\end{enumerate}
\end{exemplebox}

\courssub{Ajustement par une constante multiplicative}

En pratique, le numérateur est rarement \emph{exactement} $u'$ : il lui est \emph{proportionnel}. Comme la primitivation est linéaire (une primitive de $k \times f$ est $k \times F$), on ajuste simplement par un facteur.

\begin{propriete}[Primitive de $k \cdot u'/u$]
Si $u$ est dérivable et ne s'annule pas sur $I$, et si $k$ est une constante réelle, alors les primitives de $x \mapsto k\,\dfrac{u'(x)}{u(x)}$ sur $I$ sont les fonctions
\[ x \longmapsto k\,\ln\bigl|u(x)\bigr| + C, \quad C \in \mathbb{R}. \]
\end{propriete}

La démonstration est immédiate : $\bigl(k\ln|u|\bigr)' = k \times \dfrac{u'}{u}$.

\begin{exemplebox}[Ajustement de la constante, étape par étape]
Déterminons une primitive de $g(x) = \dfrac{x}{x^2 + 4}$ sur $\mathbb{R}$.
\begin{enumerate}
\item $u(x) = x^2 + 4 > 0$ sur $\mathbb{R}$ ; $u'(x) = 2x$.
\item Le numérateur est $x = \dfrac{1}{2} \times 2x = \dfrac{1}{2}u'(x)$. On réécrit donc :
\[ g(x) = \frac{1}{2} \times \frac{2x}{x^2+4} = \frac{1}{2}\,\frac{u'(x)}{u(x)}. \]
\item Une primitive est $G(x) = \dfrac{1}{2}\ln(x^2 + 4)$.
\item Vérification : $G'(x) = \dfrac{1}{2} \times \dfrac{2x}{x^2+4} = \dfrac{x}{x^2+4} = g(x)$. $\checkmark$
\end{enumerate}
Autre cas : $h(x) = \dfrac{6x + 3}{x^2 + x}$ sur $]0\,;+\infty[$. Ici $u(x) = x^2 + x$, $u'(x) = 2x + 1$, et le numérateur est $6x + 3 = 3(2x+1) = 3u'(x)$. Sur $]0\,;+\infty[$, $u(x) = x(x+1) > 0$, donc une primitive est $H(x) = 3\ln(x^2 + x)$.
\end{exemplebox}

\begin{attention}[On n'invente pas le $u'$ : on l'ajuste, on ne le crée pas]
L'ajustement ne fonctionne que pour un facteur \textbf{constant}. On peut écrire $x = \frac{1}{2}\times 2x$, mais on ne peut pas « faire apparaître » un $x$ manquant : $\dfrac{1}{x^2+1}$ n'est pas de la forme $k\,\dfrac{u'}{u}$, car il faudrait multiplier par $x$, qui n'est pas une constante. De même, $\dfrac{x}{x+1}$ n'est pas de la forme $u'/u$ (ici, on écrirait plutôt $\dfrac{x}{x+1} = 1 - \dfrac{1}{x+1}$, et là le second morceau est bien de la forme $u'/u$ avec $u = x+1$).
\end{attention}

\courssub{Tableau-schéma de reconnaissance}

Le tableau suivant rassemble les situations classiques à reconnaître instantanément le jour du Bac. Entraîne-toi à le refaire de mémoire, puis à vérifier chaque ligne en dérivant la primitive proposée.

\begin{popfigure}
\begin{center}
\begin{tabularx}{\linewidth}{|l|l|l|X|}
\hline
\rowcolor{popblueL} \textbf{Expression $f(x)$} & \textbf{$u(x)$} & \textbf{$u'(x)$} & \textbf{Primitive $F(x)$} \\
\hline
$\dfrac{2x}{x^2+1}$ & $x^2+1$ & $2x$ & $\ln(x^2+1)$ \\
\hline
$\dfrac{x}{x^2+4}$ & $x^2+4$ & $2x$ & $\dfrac{1}{2}\ln(x^2+4)$ \\
\hline
$\dfrac{3}{3x-5}$ & $3x-5$ & $3$ & $\ln|3x-5|$ \\
\hline
$\dfrac{1}{x}$ \;($x<0$) & $x$ & $1$ & $\ln(-x) = \ln|x|$ \\
\hline
$\dfrac{2x+1}{x^2+x-3}$ & $x^2+x-3$ & $2x+1$ & $\ln|x^2+x-3|$ \\
\hline
$\dfrac{x}{\sqrt{x^2+1}}$ & $\sqrt{x^2+1}$ & $\dfrac{x}{\sqrt{x^2+1}}$ & $\ln\!\bigl(\sqrt{x^2+1}\bigr) = \tfrac{1}{2}\ln(x^2+1)$ \\
\hline
$\dfrac{1}{x\ln x}$ \;($x>1$) & $\ln x$ & $\dfrac{1}{x}$ & $\ln(\ln x)$ \\
\hline
$\dfrac{1}{x\ln x}$ \;($0<x<1$) & $\ln x$ & $\dfrac{1}{x}$ & $\ln(-\ln x) = \ln|\ln x|$ \\
\hline
\end{tabularx}
\end{center}
\legende{Tableau de reconnaissance de la forme $u'/u$ : à chaque ligne, le numérateur est la dérivée du dénominateur, à un facteur constant près.}
\end{popfigure}

Les deux dernières lignes méritent un arrêt. Pour $f(x) = \dfrac{1}{x \ln x}$, le réflexe « dénominateur » consiste à écrire $f(x) = \dfrac{1/x}{\ln x}$ et à poser $u(x) = \ln x$ : alors $u'(x) = \dfrac{1}{x}$, qui est exactement le numérateur ! D'où une primitive $F(x) = \ln|\ln x|$. Ce résultat, classique au Baccalauréat (souvent donné comme complément ou question intermédiaire), illustre la puissance de la méthode : le « $u$ » peut lui-même contenir un logarithme.

\begin{exoresolu}[Primitives de la forme $u'/u$ : entraînement complet]
\textbf{Énoncé.} Déterminer une primitive de chacune des fonctions suivantes sur l'intervalle indiqué, en justifiant soigneusement :
\begin{enumerate}
\item $f_1(x) = \dfrac{4x}{x^2 + 3}$ sur $\mathbb{R}$ ;
\item $f_2(x) = \dfrac{5}{2x - 7}$ sur $I = \left]\dfrac{7}{2}\,;+\infty\right[$ ;
\item $f_3(x) = \dfrac{2x - 1}{x^2 - x - 6}$ sur $J = ]3\,;+\infty[$ ;
\item $f_4(x) = \dfrac{1}{x \ln x}$ sur $K = ]1\,;+\infty[$.
\end{enumerate}
\tcblower
\textbf{Solution détaillée.}
\begin{enumerate}
\item Posons $u(x) = x^2 + 3$. Pour tout réel $x$, $u(x) \geq 3 > 0$, et $u'(x) = 2x$. Le numérateur est $4x = 2 \times 2x = 2u'(x)$, donc $f_1 = 2\,\dfrac{u'}{u}$. Une primitive est :
\[ F_1(x) = 2\ln(x^2 + 3). \]
Vérification : $F_1'(x) = 2 \times \dfrac{2x}{x^2+3} = \dfrac{4x}{x^2+3}$. $\checkmark$
\item Posons $u(x) = 2x - 7$. Sur $I$, $x > \dfrac{7}{2}$ donc $u(x) > 0$ : pas besoin de valeur absolue. On a $u'(x) = 2$ et le numérateur est $5 = \dfrac{5}{2} \times 2$, donc $f_2 = \dfrac{5}{2}\,\dfrac{u'}{u}$. Une primitive est :
\[ F_2(x) = \frac{5}{2}\ln(2x - 7). \]
Vérification : $F_2'(x) = \dfrac{5}{2} \times \dfrac{2}{2x-7} = \dfrac{5}{2x-7}$. $\checkmark$
\item Posons $u(x) = x^2 - x - 6 = (x-3)(x+2)$. Sur $J = ]3\,;+\infty[$, $x - 3 > 0$ et $x + 2 > 0$, donc $u(x) > 0$. On a $u'(x) = 2x - 1$ : c'est exactement le numérateur. Une primitive est :
\[ F_3(x) = \ln(x^2 - x - 6). \]
Vérification : $F_3'(x) = \dfrac{2x-1}{x^2-x-6}$. $\checkmark$
\item Écrivons $f_4(x) = \dfrac{1/x}{\ln x}$ et posons $u(x) = \ln x$. Sur $K = ]1\,;+\infty[$, $\ln x > 0$, et $u'(x) = \dfrac{1}{x}$ : exactement le numérateur. Une primitive est :
\[ F_4(x) = \ln(\ln x). \]
Vérification : $F_4'(x) = \dfrac{1/x}{\ln x} = \dfrac{1}{x\ln x}$. $\checkmark$
\end{enumerate}
\end{exoresolu}

\courssub{Lien avec le calcul intégral}

La formule de primitives se traduit immédiatement en une formule d'intégration, grâce au théorème fondamental du calcul intégral : si $F$ est une primitive de $f$ continue sur $[a\,;b]$, alors $\displaystyle\int_a^b f(x)\,\mathrm{d}x = F(b) - F(a)$.

\begin{propriete}[Intégrale de $u'/u$]
Soit $u$ une fonction dérivable, dont la dérivée est continue, et qui \textbf{ne s'annule pas} sur le segment $[a\,;b]$. Alors :
\[ \int_a^b \frac{u'(x)}{u(x)}\,\mathrm{d}x = \Bigl[\ln\bigl|u(x)\bigr|\Bigr]_a^b = \ln\bigl|u(b)\bigr| - \ln\bigl|u(a)\bigr| = \ln\left|\frac{u(b)}{u(a)}\right|. \]
\end{propriete}

La dernière égalité vient de la propriété algébrique $\ln A - \ln B = \ln\dfrac{A}{B}$ (avec $A, B > 0$), combinée à $|\ln|p| - \ln|q||$ qui se réécrit proprement. Cette forme « condensée » est très efficace en calcul.

\begin{exemplebox}[Calcul d'intégrale direct]
Calculons $I = \displaystyle\int_0^1 \frac{2x}{x^2 + 1}\,\mathrm{d}x$.
La fonction $u(x) = x^2 + 1$ est strictement positive sur $[0\,;1]$ et $u'(x) = 2x$. Donc :
\[ I = \Bigl[\ln(x^2+1)\Bigr]_0^1 = \ln 2 - \ln 1 = \ln 2. \]
Autre exemple : $J = \displaystyle\int_1^2 \frac{1}{x}\,\mathrm{d}x = \bigl[\ln x\bigr]_1^2 = \ln 2$. Retrouver deux fois $\ln 2$ n'est pas un hasard : le changement de variable $t = x^2 + 1$... mais cela, c'est une autre histoire (hors programme, mais délicieuse).
\end{exemplebox}

\begin{exoresolu}[Aire sous une courbe hyperbolique — contexte Bac]
\textbf{Énoncé.} Soit $f$ la fonction définie sur $[1\,;3]$ par $f(x) = \dfrac{4x + 2}{x^2 + x}$. On admet que $f(x) > 0$ sur $[1\,;3]$. Calculer, en unités d'aire, l'aire $\mathcal{A}$ du domaine délimité par la courbe de $f$, l'axe des abscisses et les droites d'équations $x = 1$ et $x = 3$. Donner la valeur exacte, puis une valeur approchée à $10^{-2}$ près.
\tcblower
\textbf{Solution détaillée.} L'aire cherchée est $\mathcal{A} = \displaystyle\int_1^3 f(x)\,\mathrm{d}x$.
Posons $u(x) = x^2 + x = x(x+1)$. Sur $[1\,;3]$, $u(x) > 0$, et $u'(x) = 2x + 1$. Le numérateur est $4x + 2 = 2(2x+1) = 2u'(x)$, donc $f = 2\,\dfrac{u'}{u}$ et une primitive est $F(x) = 2\ln(x^2+x)$. Alors :
\begin{align*}
\mathcal{A} &= \Bigl[2\ln(x^2+x)\Bigr]_1^3 = 2\ln(12) - 2\ln(2) \\
&= 2\ln\left(\frac{12}{2}\right) = 2\ln 6.
\end{align*}
Numériquement, $\mathcal{A} = 2\ln 6 \approx 2 \times 1{,}7918 \approx 3{,}58$ unités d'aire.
\end{exoresolu}

\begin{savaistu}[Pourquoi cette formule est partout]
La forme $\dfrac{u'}{u}$ porte un nom en sciences : c'est la \emph{dérivée logarithmique} de $u$. Elle mesure le \emph{taux de variation relatif} d'une grandeur. C'est exactement elle qui apparaît dans les modèles de croissance : quand on dit qu'une population de bactéries croît « de $5\,\%$ par heure », on dit que $\dfrac{u'}{u} = 0{,}05$, et la solution est $u(t) = u_0\,e^{0{,}05t}$ — dont le logarithme est une fonction affine du temps. Les ingénieurs de CAMTEL qui modélisent l'atténuation d'un signal dans une fibre optique, ou les pharmaciens qui suivent l'élimination d'un médicament dans le sang, manipulent sans le dire des primitives de $\dfrac{u'}{u}$. Historiquement, c'est Grégoire de Saint-Vincent qui, en 1647, remarqua que l'aire sous l'hyperbole $y = \dfrac{1}{x}$ possède la propriété $\mathcal{A}(ab) = \mathcal{A}(a) + \mathcal{A}(b)$ : il venait, sans le savoir, de découvrir le logarithme népérien comme primitive de $\dfrac{1}{x}$, vingt ans avant que Newton ne le baptise « logarithme naturel ».
\end{savaistu}

\begin{methode}[Réflexes de copie pour le jour du Bac]
\begin{enumerate}
\item Devant toute fraction à intégrer, tester d'abord la forme $\dfrac{u'}{u}$ : c'est le cas le plus fréquent et le plus rapide.
\item Toujours annoncer l'intervalle et le signe de $u$ avant d'écrire $\ln u$ ou $\ln|u|$.
\item Écrire l'ajustement de constante explicitement : « $x = \frac{1}{2}(2x)$ » — le correcteur doit voir que tu sais ce que tu fais.
\item Terminer par la vérification mentale : dériver la primitive trouvée. Cinq secondes qui sauvent des points.
\end{enumerate}
\end{methode}

\begin{exercice}[À toi de jouer]
Déterminer une primitive de chacune des fonctions suivantes sur l'intervalle proposé :
\begin{enumerate}
\item $f(x) = \dfrac{6x^2}{2x^3 + 5}$ sur $\mathbb{R}$ ;
\item $g(x) = \dfrac{1}{x\ln x}$ sur $\left]0\,;\dfrac{1}{2}\right]$ (attention au signe de $\ln x$ !) ;
\item $h(x) = \dfrac{x - 2}{x^2 - 4x + 5}$ sur $\mathbb{R}$ (vérifier d'abord que le dénominateur ne s'annule pas).
\end{enumerate}
Puis calculer $\displaystyle\int_1^{e} \frac{1}{x\ln x \cdot \ln x}\,\mathrm{d}x$... non, plutôt : calculer $\displaystyle\int_2^{e+1} \frac{1}{x - 1}\,\mathrm{d}x$ et donner la valeur exacte.
\end{exercice}

\begin{corrige}[Exercice : primitives de la forme $u'/u$]
\begin{enumerate}
\item $u(x) = 2x^3 + 5$, $u'(x) = 6x^2$ : forme exacte. Mais $u$ s'annule en $x = -\sqrt[3]{5/2}$ ! Sur $\mathbb{R}$ il faut donc préciser : sur chaque intervalle où $u \neq 0$, une primitive est $F(x) = \ln|2x^3 + 5|$. (Sur $]-\sqrt[3]{5/2}\,;+\infty[$, $u > 0$ et $F(x) = \ln(2x^3+5)$.)
\item Sur $\left]0\,;\tfrac{1}{2}\right]$, $\ln x < 0$, donc $|\ln x| = -\ln x$ et une primitive est $G(x) = \ln(-\ln x) = \ln|\ln x|$.
\item $x^2 - 4x + 5 = (x-2)^2 + 1 \geq 1 > 0$ sur $\mathbb{R}$. $u'(x) = 2x - 4$ et $x - 2 = \frac{1}{2}(2x-4)$, donc $H(x) = \frac{1}{2}\ln(x^2 - 4x + 5)$.
\item $\displaystyle\int_2^{e+1}\frac{\mathrm{d}x}{x-1} = \bigl[\ln(x-1)\bigr]_2^{e+1} = \ln e - \ln 1 = 1$.
\end{enumerate}
\end{corrige}

En résumé : la formule « une primitive de $\dfrac{u'}{u}$ est $\ln|u|$ » est le pendant intégral de la dérivée de $\ln u$ étudiée dans la section précédente. Maîtriser sa reconnaissance, l'ajustement des constantes et la gestion de la valeur absolue, c'est tenir entre ses mains l'un des outils les plus puissants du calcul intégral de Terminale. Nous allons maintenant mettre tous ces acquis — propriétés algébriques, limites, dérivation, primitives — au service de la résolution d'équations, d'inéquations, de systèmes et de problèmes concrets : c'est l'objet de la dernière section de cette leçon.

\coursec{Équations, inéquations, systèmes et applications}

Dans la section précédente, nous avons appris à \emph{remonter} d'une fonction vers ses primitives : toute fonction de la forme $\dfrac{u'}{u}$ admet pour primitive $\ln|u|$. Nous allons maintenant utiliser le logarithme népérien comme un \emph{outil de résolution} : grâce à ses propriétés algébriques et à sa stricte croissance, il transforme des équations apparemment compliquées en équations simples, et il permet de répondre à des questions très concrètes : en combien d'années un capital placé à la BEAC double-t-il ? Comment mesure-t-on l'acidité de l'eau distribuée par la CAMWATER ? Pourquoi l'échelle des séismes est-elle logarithmique ?

\courssub{Le principe fondamental : l'injectivité de la fonction ln}

\textbf{Intuition.}\par
La fonction $\ln$ est \emph{strictement croissante} sur $]0\,;+\infty[$ : deux nombres strictement positifs \emph{distincts} ont toujours des logarithmes \emph{distincts}. Autrement dit, le logarithme ne peut pas « confondre » deux nombres : si deux logarithmes sont égaux, c'est que les nombres de départ étaient égaux. C'est cette propriété qui nous autorise à « supprimer le ln » des deux côtés d'une équation.

\begin{propriete}[Équivalence fondamentale]
Pour tous réels $a$ et $b$ \emph{strictement positifs} :
\[ \ln a = \ln b \iff a = b \qquad\text{et}\qquad \ln a < \ln b \iff a < b. \]
De plus, $\ln$ réalise une \cle{bijection} de $]0\,;+\infty[$ sur $\mathbb{R}$ : tout réel $a$ est le logarithme d'\emph{un unique} réel strictement positif.
\end{propriete}

\begin{experience}[Pourquoi ces équivalences sont-elles vraies ?]
Raisonnons pas à pas.
\begin{enumerate}
\item \textbf{Sens direct ($\Rightarrow$).} Si $a = b$, alors évidemment $\ln a = \ln b$.
\item \textbf{Sens réciproque ($\Leftarrow$).} Supposons $\ln a = \ln b$ et montrons que $a = b$. Raisonnons par l'absurde : si $a \neq b$, alors soit $a < b$, soit $a > b$. Comme $\ln$ est \emph{strictement croissante}, on aurait $\ln a < \ln b$ ou $\ln a > \ln b$ : dans les deux cas, contradiction avec $\ln a = \ln b$. Donc $a = b$.
\item \textbf{Pour l'inégalité.} Le sens direct vient de la croissance de $\ln$ ; le sens réciproque se démontre de la même façon par l'absurde.
\item \textbf{Bijectivité.} L'injectivité vient d'être prouvée. La \emph{surjectivité} (tout réel est atteint) est \textbf{admise} à ce niveau : elle repose sur le fait que $\ln$ est continue sur $]0\,;+\infty[$ avec $\lim_{0^+}\ln = -\infty$ et $\lim_{+\infty}\ln = +\infty$ (théorème des valeurs intermédiaires). Nous la préciserons avec la fonction exponentielle.
\end{enumerate}
\end{experience}

\begin{attention}[La condition d'existence, toujours en premier !]
L'équivalence $\ln a = \ln b \iff a = b$ n'a de sens que si $a > 0$ \textbf{et} $b > 0$. Avant toute résolution, on \textbf{doit} déterminer l'ensemble sur lequel tous les logarithmes de l'équation existent. Une « solution » trouvée après transformation mais qui ne vérifie pas ces conditions doit être \emph{rejetée}. C'est l'erreur la plus fréquente au Baccalauréat sur ce chapitre.
\end{attention}

\begin{methode}[Organigramme de résolution d'une équation ou inéquation avec ln]
Toute résolution suit le même schéma en quatre étapes :
\begin{enumerate}
\item \textbf{Conditions d'existence} : écrire que chaque quantité sous un ln est strictement positive ; en déduire le domaine $D$.
\item \textbf{Transformation} : utiliser les propriétés algébriques ($\ln a + \ln b = \ln(ab)$, etc.) ou poser $X = \ln x$.
\item \textbf{Résolution} : résoudre l'équation transformée (souvent une équation du $1^{\text{re}}$ ou $2^{\text{ème}}$ degré).
\item \textbf{Vérification} : ne retenir que les solutions appartenant à $D$, puis conclure par l'ensemble $S$.
\end{enumerate}
\end{methode}

\begin{popfigure}
\begin{center}
\begin{tikzpicture}[
  node distance=0.55cm,
  etape/.style={rectangle, rounded corners=4pt, draw=popblue, fill=popblueL, text width=6.2cm, align=center, minimum height=1cm, font=\small},
  fleche/.style={-{Stealth[length=3mm]}, thick, popdark}
]
\node[etape, align=center] (e1) {\textbf{Étape 1 — Conditions d'existence}\\ Chaque argument de ln est $>0$ ; domaine $D$};
\node[etape, below=of e1, align=center] (e2) {\textbf{Étape 2 — Transformation}\\ Propriétés algébriques ou changement $X=\ln x$};
\node[etape, below=of e2, align=center] (e3) {\textbf{Étape 3 — Résolution}\\ Équation transformée ($1^{\text{re}}$ ou $2^{\text{ème}}$ degré)};
\node[etape, below=of e3, draw=poporange, fill=poporangeL, align=center] (e4) {\textbf{Étape 4 — Vérification et conclusion}\\ Ne garder que les solutions dans $D$ ; écrire $S$};
\draw[fleche] (e1) -- (e2);
\draw[fleche] (e2) -- (e3);
\draw[fleche] (e3) -- (e4);
\end{tikzpicture}
\end{center}
\legende{Organigramme de résolution : la vérification finale (étape 4) est obligatoire.}
\end{popfigure}

\courssub{Équations du type $\ln x = a$ et $\ln u(x) = \ln v(x)$}

\textbf{Intuition.}\par
Puisque $\ln$ est bijective de $]0\,;+\infty[$ sur $\mathbb{R}$, l'équation $\ln x = a$ possède, quel que soit le réel $a$, \emph{une unique solution}. Graphiquement, c'est l'abscisse du point d'intersection de la courbe de $\ln$ avec la droite horizontale d'équation $y = a$.

\begin{propriete}[Solution de $\ln x = a$]
Pour tout réel $a$, l'équation $\ln x = a$ admet une \cle{unique solution} dans $]0\,;+\infty[$. On note cette solution $x = e^{a}$ (ou $\mathrm{e}^a$). En particulier :
\[ \ln x = 0 \iff x = 1 \qquad\text{et}\qquad \ln x = 1 \iff x = e. \]
\end{propriete}

\begin{savaistu}[Ouverture : la fonction exponentielle]
Le nombre $e^{a}$ est, par définition, \emph{l'unique réel strictement positif dont le logarithme vaut $a$}. Cela définit une nouvelle fonction, la \cle{fonction exponentielle} $x \mapsto e^{x}$, réciproque de $\ln$ : on a $\ln(e^{x}) = x$ pour tout réel $x$ et $e^{\ln x} = x$ pour tout $x > 0$. Elle sera étudiée en détail dans la leçon suivante ; ici, nous l'utilisons simplement comme une \emph{notation} pour la solution de $\ln x = a$.
\end{savaistu}

\begin{popfigure}
\begin{center}
\begin{tikzpicture}
\begin{axis}[
  width=0.85\linewidth, height=6.5cm,
  axis lines=middle,
  xlabel={$x$}, ylabel={$y$},
  xmin=-0.5, xmax=9, ymin=-2.5, ymax=2.5,
  xtick={1,2.718,7.389},
  xticklabels={$1$,$e$,$e^{2}$},
  ytick={1,2},
  yticklabels={$1$,$2$},
  grid=both, grid style={popline!40},
  legend style={at={(0.03,0.97)}, anchor=north west, font=\small}
]
\addplot[popcurve, domain=0.05:9, samples=200] {ln(x)};
\addlegendentry{$y=\ln x$}
\addplot[poporange, thick, dashed, domain=0:9] {2};
\addlegendentry{$y=2$}
\addplot[only marks, mark=*, poppink, mark size=2.5pt] coordinates {(7.389,2)};
\draw[popmuted, dashed] (axis cs:7.389,0) -- (axis cs:7.389,2);
\node[popdark, font=\small, anchor=south west] at (axis cs:7.389,2) {$(e^{2}\,;\,2)$};
\end{axis}
\end{tikzpicture}
\end{center}
\legende{La solution de $\ln x = 2$ est l'abscisse $x = e^{2} \approx 7,39$ du point d'intersection de la courbe avec la droite $y = 2$.}
\end{popfigure}

\begin{propriete}[Équations et inéquations du type $\ln u = \ln v$]
Soient $u$ et $v$ deux fonctions. Sous les conditions $u(x) > 0$ et $v(x) > 0$ :
\[ \ln u(x) = \ln v(x) \iff u(x) = v(x) \qquad\text{et}\qquad \ln u(x) < \ln v(x) \iff u(x) < v(x). \]
Pour les inéquations, le sens de l'inégalité est \emph{conservé} car $\ln$ est \emph{croissante}.
\end{propriete}

\begin{exoresolu}[Équation du type $\ln u = \ln v$]
Résoudre dans $\mathbb{R}$ l'équation $\ln(3x - 1) = \ln(x + 5)$.
\tcblower
\textbf{Étape 1 — Conditions d'existence.}
\[ \begin{cases} 3x - 1 > 0 \\ x + 5 > 0 \end{cases} \iff \begin{cases} x > \dfrac{1}{3} \\ x > -5 \end{cases} \iff x > \dfrac{1}{3}. \]
Le domaine est $D = \left]\dfrac{1}{3}\,;+\infty\right[$.

\textbf{Étapes 2 et 3 — Transformation et résolution.} Sur $D$, par injectivité de $\ln$ :
\[ \ln(3x-1) = \ln(x+5) \iff 3x - 1 = x + 5 \iff 2x = 6 \iff x = 3. \]

\textbf{Étape 4 — Vérification.} $3 \in D$ car $3 > \dfrac{1}{3}$. De plus $3\times 3 - 1 = 8$ et $3 + 5 = 8$ : on a bien $\ln 8 = \ln 8$.
\[ S = \{3\}. \]
\end{exoresolu}

\begin{exoresolu}[Inéquation du type $\ln u < \ln v$]
Résoudre dans $\mathbb{R}$ l'inéquation $\ln(2x + 3) < \ln(6 - x)$.
\tcblower
\textbf{Étape 1 — Conditions.} $2x + 3 > 0 \iff x > -\dfrac{3}{2}$ et $6 - x > 0 \iff x < 6$. Donc $D = \left]-\dfrac{3}{2}\,;6\right[$.

\textbf{Étapes 2 et 3.} $\ln$ étant croissante, elle conserve l'ordre :
\[ 2x + 3 < 6 - x \iff 3x < 3 \iff x < 1. \]

\textbf{Étape 4.} On intersecte avec le domaine : $S = D \cap \left]-\infty\,;1\right[ = \left]-\dfrac{3}{2}\,;1\right[$.
\end{exoresolu}

\begin{attention}[Ne jamais « croiser » avec le signe de ln]
On ne peut pas écrire $\ln u < \ln v \Rightarrow u < v$ sans avoir d'abord prouvé que $u > 0$ et $v > 0$. Par exemple, l'inéquation $\ln(x^2) < \ln 1$ exige $x \neq 0$ : sa solution est $]-1\,;0[\,\cup\,]0\,;1[$, et non $]-1\,;1[$. Le point $x = 0$ doit être \emph{exclu}.
\end{attention}

\courssub{Équations et inéquations avec inconnue auxiliaire $X = \ln x$}

\textbf{Intuition.}\par
Certaines équations contiennent à la fois $\ln x$ et $(\ln x)^2$ : ce sont des équations du second degré « déguisées ». L'astuce consiste à poser $X = \ln x$ : l'équation devient une équation classique en $X$, que l'on sait résoudre depuis la Première. On revient ensuite à $x$ grâce à $x = e^{X}$.

\begin{methode}[Résolution d'une équation en $\ln x$]
\begin{enumerate}
\item \textbf{Condition d'existence} : $x > 0$.
\item \textbf{Changement d'inconnue} : poser $X = \ln x$ (avec $X \in \mathbb{R}$, sans autre restriction).
\item \textbf{Résoudre} l'équation en $X$ (souvent du second degré : discriminant $\Delta$).
\item \textbf{Revenir à $x$} : pour chaque solution $X$, écrire $x = e^{X}$. Comme $e^{X} > 0$ pour tout réel $X$, la condition $x > 0$ est \emph{toujours} satisfaite : aucune solution n'est rejetée à cette étape.
\end{enumerate}
\end{methode}

\begin{exemplebox}[L'exemple classique du Baccalauréat]
Résolvons $(\ln x)^2 - \ln x - 2 = 0$.

\textbf{Condition} : $x > 0$.

Posons $X = \ln x$. L'équation devient :
\[ X^2 - X - 2 = 0. \]
Discriminant : $\Delta = (-1)^2 - 4 \times 1 \times (-2) = 1 + 8 = 9 > 0$. Deux racines :
\[ X_1 = \frac{1 - 3}{2} = -1 \qquad\text{et}\qquad X_2 = \frac{1 + 3}{2} = 2. \]
Retour à $x$ : $\ln x = -1 \iff x = e^{-1} = \dfrac{1}{e}$ et $\ln x = 2 \iff x = e^{2}$.
\[ S = \left\{ \dfrac{1}{e}\,;\, e^{2} \right\}. \]
Vérification pour $x = e^{2}$ : $(\ln e^{2})^2 - \ln e^{2} - 2 = 4 - 2 - 2 = 0$. Correct.
\end{exemplebox}

\begin{exoresolu}[Inéquation du second degré en $\ln x$]
Résoudre dans $\mathbb{R}$ l'inéquation $(\ln x)^2 - \ln x - 2 < 0$.
\tcblower
\textbf{Condition} : $x > 0$. Posons $X = \ln x$ : on obtient $X^2 - X - 2 < 0$.

Le trinôme $X^2 - X - 2$ a pour racines $-1$ et $2$ (voir exemple précédent). Un trinôme est du signe de $a = 1 > 0$ \emph{à l'extérieur} des racines, donc négatif \emph{entre} les racines :
\[ X^2 - X - 2 < 0 \iff -1 < X < 2 \iff -1 < \ln x < 2. \]
Comme $\ln$ est croissante et bijective, on applique la fonction exponentielle (qui conserve l'ordre) :
\[ e^{-1} < x < e^{2}. \]
Ces bornes sont strictement positives : la condition $x > 0$ est satisfaite.
\[ S = \left] \dfrac{1}{e}\,;\, e^{2} \right[. \]
\end{exoresolu}

\begin{attention}[Deux pièges classiques]
\begin{itemize}
\item Ne concluez jamais « $X = -1$ donc $x = -1$ » : on doit revenir à $x$ par $x = e^{X}$, et $e^{-1}$ est bien positif même si $X$ est négatif.
\item Pour une inéquation, n'oubliez pas que $\ln x < 0 \iff 0 < x < 1$ : la borne $0$ doit apparaître, car $x$ doit rester strictement positif.
\end{itemize}
\end{attention}

\courssub{Systèmes à deux inconnues faisant intervenir $\ln x$ et $\ln y$}

\textbf{Intuition.}\par
Un système où apparaissent $\ln x$ et $\ln y$ se ramène, en posant $X = \ln x$ et $Y = \ln y$, à un système linéaire de deux équations à deux inconnues, que l'on résout par substitution ou combinaison. On revient ensuite à $(x\,;y)$ par $x = e^{X}$ et $y = e^{Y}$.

\begin{methode}[Résolution d'un système en $\ln x$ et $\ln y$]
\begin{enumerate}
\item Écrire les conditions d'existence : $x > 0$ et $y > 0$.
\item Poser $X = \ln x$ et $Y = \ln y$ ; résoudre le système linéaire en $(X\,;Y)$.
\item Conclure : $(x\,;y) = (e^{X}\,;\,e^{Y})$. Les conditions $x > 0$, $y > 0$ sont automatiquement satisfaites.
\end{enumerate}
\end{methode}

\begin{exoresolu}[Système avec logarithmes]
Résoudre dans $\mathbb{R}^2$ le système :
\[ \begin{cases} \ln x + \ln y = \ln 6 \\ \ln x - \ln y = 0 \end{cases} \]
\tcblower
\textbf{Conditions} : $x > 0$ et $y > 0$.

Posons $X = \ln x$ et $Y = \ln y$. Le système devient :
\[ \begin{cases} X + Y = \ln 6 \\ X - Y = 0 \end{cases} \]
Par addition membre à membre : $2X = \ln 6$, donc $X = \dfrac{\ln 6}{2} = \ln\sqrt{6}$. La deuxième équation donne $Y = X = \ln\sqrt{6}$.

Retour aux inconnues : $x = e^{\ln\sqrt{6}} = \sqrt{6}$ et $y = \sqrt{6}$.

\textbf{Vérification} : $\ln\sqrt{6} + \ln\sqrt{6} = \ln(\sqrt{6}\times\sqrt{6}) = \ln 6$ et $\ln\sqrt{6} - \ln\sqrt{6} = 0$.
\[ S = \left\{ \left(\sqrt{6}\,;\,\sqrt{6}\right) \right\}. \]
\end{exoresolu}

\courssub{Applications : le logarithme, outil pour « inverser » les puissances}

\textbf{L'idée maîtresse.}\par
Beaucoup de phénomènes (placements bancaires, croissance démographique, désintégration radioactive) sont modélisés par une puissance $a^{n}$. Pour trouver $n$ connaissant $a^{n}$, on « descend l'exposant » grâce au logarithme :
\[ a^{n} = b \iff n\,\ln a = \ln b \iff n = \frac{\ln b}{\ln a}. \]
C'est l'une des applications les plus puissantes du logarithme népérien.

\begin{exemplebox}[Temps de doublement d'un capital]
Mme Etoundi place un capital à intérêts composés au taux annuel de $7\,\%$. Au bout de $n$ années, son capital est multiplié par $1{,}07^{n}$. En combien d'années le capital aura-t-il \emph{doublé} ?

On résout :
\[ 1{,}07^{n} = 2 \iff n\,\ln 1{,}07 = \ln 2 \iff n = \frac{\ln 2}{\ln 1{,}07} \approx \frac{0{,}693}{0{,}0677} \approx 10,2. \]
Le capital double au bout d'environ $10,2$ ans, c'est-à-dire lors de la \textbf{$11^{\text{ème}}$ année} (car $n$ doit être entier : $1{,}07^{10} \approx 1,97 < 2$ et $1{,}07^{11} \approx 2,10 > 2$).
\end{exemplebox}

\begin{savaistu}[Le pH de l'eau de la CAMWATER]
Le \cle{pH} d'une solution mesure son acidité : $\mathrm{pH} = -\log[\mathrm{H}^{+}]$, où $[\mathrm{H}^{+}]$ est la concentration en ions hydrogène et $\log$ désigne le \emph{logarithme décimal}, lié au logarithme népérien par :
\[ \log x = \frac{\ln x}{\ln 10}. \]
L'eau potable distribuée par la CAMWATER doit avoir un pH compris entre $6,5$ et $8,5$. Comme le pH est une échelle logarithmique, une eau de pH $6$ est \emph{dix fois} plus acide qu'une eau de pH $7$ ! La même idée d'échelle logarithmique se retrouve en sismologie (magnitude de Richter : $+1$ de magnitude correspond à une amplitude multipliée par $10$) et en acoustique (le décibel : $\mathrm{dB} = 10\log\dfrac{I}{I_0}$). Ces échelles permettent d'« écraser » des grandeurs qui varient sur des ordres de grandeur immenses pour les rendre lisibles.
\end{savaistu}

\begin{exoresolu}[Application : croissance d'une population]
La population d'une ville du littoral camerounais croît de $4\,\%$ par an. En combien d'années aura-t-elle été multipliée par $1,5$ ? (On donnera la réponse à $0,1$ près.)
\tcblower
Au bout de $n$ années, la population est multipliée par $1,04^{n}$. On résout :
\[ 1,04^{n} = 1,5 \iff n = \frac{\ln 1,5}{\ln 1,04} \approx \frac{0,4055}{0,0392} \approx 10,3. \]
La population aura été multipliée par $1,5$ au bout d'environ $10,3$ ans, donc au cours de la $11^{\text{ème}}$ année.
\end{exoresolu}

\courssub{Synthèse méthodologique de la section}

\begin{aretenir}[Les réflexes à acquérir]
\begin{itemize}
\item \textbf{Toujours commencer} par les conditions d'existence : chaque argument d'un logarithme doit être strictement positif.
\item \textbf{Injectivité} : $\ln u = \ln v \iff u = v$ (sous conditions) ; $\ln u < \ln v \iff u < v$ car $\ln$ est \emph{croissante} (le sens de l'inégalité est conservé).
\item \textbf{Équation $\ln x = a$} : unique solution $x = e^{a}$.
\item \textbf{Changement d'inconnue} : poser $X = \ln x$ (ou $X = \ln x$, $Y = \ln y$ pour un système), résoudre, puis revenir par $x = e^{X}$, toujours positif.
\item \textbf{Puissances} : $a^{n} = b \iff n = \dfrac{\ln b}{\ln a}$ — c'est l'outil des problèmes de doublement, de pH et d'échelles logarithmiques.
\item \textbf{Toujours finir} par la vérification : toute solution hors du domaine est rejetée.
\end{itemize}
\end{aretenir}

\begin{exercice}[Pour s'entraîner]
\begin{enumerate}
\item Résoudre dans $\mathbb{R}$ : \textbf{a)} $\ln(2x - 4) = \ln(x + 1)$ ; \textbf{b)} $\ln(x^2 - 3) = 0$.
\item Résoudre dans $\mathbb{R}$ l'inéquation $\ln(3 - x) \geq \ln(2x)$.
\item Résoudre $(\ln x)^2 + \ln x - 6 = 0$.
\item Résoudre le système : $\begin{cases} \ln x + 2\ln y = \ln 8 \\ 2\ln x - \ln y = \ln 2 \end{cases}$
\item Un capital placé à $5\,\%$ l'an à intérêts composés doit tripler. En combien d'années ?
\end{enumerate}
\end{exercice}

\begin{corrige}[Exercice]
\begin{enumerate}
\item \textbf{a)} Conditions : $x > 2$ et $x > -1$, donc $D = ]2\,;+\infty[$. Alors $2x - 4 = x + 1 \iff x = 5 \in D$. $S = \{5\}$.

\textbf{b)} Condition : $x^2 - 3 > 0 \iff x \in \left]-\infty\,;-\sqrt{3}\right[ \cup \left]\sqrt{3}\,;+\infty\right[$. L'équation s'écrit $x^2 - 3 = 1$, soit $x^2 = 4$, $x = \pm 2$. Les deux valeurs vérifient $|2| > \sqrt{3}$ : $S = \{-2\,;\,2\}$.

\item Conditions : $x < 3$ et $x > 0$, donc $D = ]0\,;3[$. Comme $\ln$ est croissante : $3 - x \geq 2x \iff x \leq 1$. $S = ]0\,;1]$.

\item Condition : $x > 0$. Posons $X = \ln x$ : $X^2 + X - 6 = 0$, $\Delta = 1 + 24 = 25$, $X = 2$ ou $X = -3$. Donc $x = e^{2}$ ou $x = e^{-3}$. $S = \{e^{-3}\,;\,e^{2}\}$.

\item Conditions : $x > 0$, $y > 0$. Posons $X = \ln x$, $Y = \ln y$ :
\[ \begin{cases} X + 2Y = \ln 8 = 3\ln 2 \\ 2X - Y = \ln 2 \end{cases} \]
De la deuxième équation, $Y = 2X - \ln 2$. En reportant dans la première :
$X + 2(2X - \ln 2) = 3\ln 2 \iff 5X - 2\ln 2 = 3\ln 2 \iff 5X = 5\ln 2 \iff X = \ln 2$.
Alors $Y = 2\ln 2 - \ln 2 = \ln 2$.
D'où $x = e^{\ln 2} = 2$ et $y = 2$.
Vérification : $\ln 2 + 2\ln 2 = 3\ln 2 = \ln 8$ ; $2\ln 2 - \ln 2 = \ln 2$.
$S = \{(2\,;\,2)\}$.

\item On résout $1,05^{n} = 3 \iff n = \dfrac{\ln 3}{\ln 1,05} \approx \dfrac{1,0986}{0,0488} \approx 22,5$. Le capital triple au bout d'environ $22,5$ ans, donc lors de la $23^{\text{ème}}$ année.
\end{enumerate}
\end{corrige}

\newpage

\coursec{Activité d'intégration}

\coursec{Activité d'intégration : Gestion d'un forage à Nkongsamba}

\courssub{Objectifs de l'activité}

Cette activité d'intégration te place dans une situation sanitaire et économique authentique : la gestion d'un forage par une coopérative agricole de Nkongsamba. À travers trois tâches complémentaires, tu vas mobiliser l'ensemble des savoir-faire de la leçon sur la fonction logarithme népérien :

\begin{itemize}
\item utiliser les propriétés algébriques de $\ln$ pour transformer une expression (passage du logarithme décimal au logarithme népérien) ;
\item résoudre une équation de la forme $e^{at} = b$ et une inéquation de la forme $e^{at} > b$ en passant au logarithme népérien ;
\item résoudre une équation et une inéquation faisant intervenir $\ln x$, en vérifiant les conditions d'existence ;
\item interpréter les résultats obtenus dans un contexte concret (conformité d'une eau potable, croissance d'une population bactérienne, placement à intérêts composés) ;
\item comparer des croissances exponentielles et rédiger une synthèse argumentée.
\end{itemize}

\courssub{Situation}

La coopérative agricole \og Vallée du Moungo \fg{} de Nkongsamba, dans la région du Littoral, exploite un forage qui alimente en eau potable plusieurs quartiers et irrigue des plantations de bananiers. Trois préoccupations occupent actuellement le comité de gestion, qui te sollicite en tant qu'expert(e) en mathématiques.

\textbf{Préoccupation 1 : la qualité sanitaire de l'eau.} Le laboratoire de la coopérative mesure l'acidité de l'eau du forage grâce au pH, défini par
\[ \mathrm{pH} = -\log\left[\mathrm{H}^{+}\right], \]
où $\left[\mathrm{H}^{+}\right]$ désigne la concentration en ions hydrogène, exprimée en mol/L, et $\log$ le logarithme décimal. On rappelle la relation de conversion :
\[ \log x = \frac{\ln x}{\ln 10} \quad (x > 0). \]
La dernière analyse donne un pH de $6{,}8$. La norme camerounaise de potabilité exige que le pH soit compris entre $6{,}5$ et $8{,}5$.

\begin{popfigure}
\begin{tikzpicture}[scale=0.9]
\draw[popblue, thick] (0,0) -- (14,0);
\foreach \x/\label in {0/0, 1/1, 2/2, 3/3, 4/4, 5/5, 6/6, 7/7, 8/8, 9/9, 10/10, 11/11, 12/12, 13/13, 14/14} {
    \draw (\x,0.1) -- (\x,-0.1) node[below] {\small $\label$};
}
\draw[popgreen, thick] (6.5,0.3) -- (8.5,0.3) node[midway, above] {\textbf{Zone de potabilité $[6,5; 8,5]$}};
\draw[poporange, very thick] (6.8, -0.3) -- (6.8, 0.6) node[above] {\textbf{pH mesuré $= 6,8$}};
\draw[popred, dashed] (0,-0.5) -- (14,-0.5);
\draw[popred, dashed] (0,0.5) -- (14,0.5);
\node[popred] at (1,-0.5) {Acide};
\node[popblue] at (13,0.5) {Basique};
\end{tikzpicture}
\legende{Échelle de pH et position de la mesure du forage.}
\end{popfigure}

\textbf{Préoccupation 2 : la prolifération bactérienne dans le réservoir.} Le réservoir de stockage, mal protégé de la chaleur, abrite une population de bactéries dont l'effectif au temps $t$ (en heures) est modélisé par
\[ N(t) = N_0\, e^{0{,}3t}, \]
où $N_0$ est l'effectif initial ($N_0 > 0$). Le comité veut connaître le temps de doublement de cette population et le moment à partir duquel l'effectif dépassera $1000\,N_0$, seuil au-delà duquel l'eau est déclarée impropre à la consommation.

\begin{popfigure}
\begin{tikzpicture}[scale=0.7, domain=0:8]
\begin{axis}[
    axis lines=middle,
    xlabel={$t$ (heures)}, ylabel={$N(t)/N_0$},
    xmin=0, xmax=8, ymin=0, ymax=12,
    xtick={0,1,2,3,4,5,6,7,8}, ytick={1,2,4,8},
    width=\linewidth, height=6cm,
    grid=major, grid style={popline},
    popcurve/.style={popblue, thick, samples=100}
]
\addplot[popcurve] {exp(0.3*x)};
\addplot[poporange, dashed] coordinates {(0,2) (2.31,2) (2.31,0)} node[below right, pos=0.5] {$t_2 \approx 2,31$ h};
\addplot[popred, dashed] coordinates {(0,1000) (23,1000) (23,0)} node[below right, pos=0.5] {Seuil $1000$};
\addplot[only marks, mark=*, poporange] coordinates {(2.31,2)};
\addplot[only marks, mark=*, popred] coordinates {(23.03,1000)};
\end{axis}
\end{tikzpicture}
\legende{Croissance bactérienne $N(t)=N_0 e^{0,3t}$ : temps de doublement et seuil critique.}
\end{popfigure}

\textbf{Préoccupation 3 : le placement de la recette.} La coopérative place une recette de $2\,000\,000$ FCFA à la CCA-Bank à un taux annuel de $6\,\%$, à intérêts composés. Le capital au bout de $n$ années est donc
\[ C(n) = 2\,000\,000 \times (1{,}06)^n \quad (n \in \mathbb{N}). \]
Le trésorier souhaite savoir au bout de combien d'années le capital aura doublé, et comparer ce temps de doublement à celui de la population bactérienne.

\courssub{Consignes}

La classe est répartie en trois groupes. Chaque groupe traite une tâche, puis présente ses résultats ; la classe confronte les méthodes et rédige une synthèse commune.

\textbf{Tâche 1 (Groupe 1) : le pH de l'eau.}
\begin{enumerate}
\item En utilisant la relation $\log x = \dfrac{\ln x}{\ln 10}$, montrer que $\mathrm{pH} = -\dfrac{\ln\left[\mathrm{H}^{+}\right]}{\ln 10}$.
\item Déterminer la concentration $\left[\mathrm{H}^{+}\right]$ correspondant à un pH de $6{,}8$ (on donnera une valeur approchée en écriture scientifique).
\item Traduire la norme $6{,}5 \leq \mathrm{pH} \leq 8{,}5$ en un encadrement de $\left[\mathrm{H}^{+}\right]$, puis vérifier que l'eau du forage est conforme.
\end{enumerate}

\textbf{Tâche 2 (Groupe 2) : la croissance bactérienne.}
\begin{enumerate}
\item Résoudre dans $\mathbb{R}_+$ l'équation $e^{0{,}3t} = 2$. En déduire le temps de doublement de la population, arrondi à la minute près.
\item Résoudre l'inéquation $N(t) > 1000\,N_0$. Au bout de combien d'heures entières l'eau devient-elle impropre à la consommation ?
\item Combien de temps s'écoule-t-il entre le premier doublement et le dépassement du seuil critique ? Commenter.
\end{enumerate}

\textbf{Tâche 3 (Groupe 3) : le placement.}
\begin{enumerate}
\item Montrer que le temps de doublement du capital est la plus petite solution entière de l'inéquation $(1{,}06)^n \geq 2$, puis résoudre cette inéquation à l'aide de $\ln$.
\item Au bout de combien d'années le capital atteindra-t-il $4\,000\,000$ FCFA ? Vérifier la cohérence avec la question précédente.
\item Comparer le temps de doublement du capital (en années) à celui des bactéries (en heures). Quelle conclusion économique et sanitaire en tires-tu ?
\end{enumerate}

\textbf{Synthèse commune.} Rédiger en quelques lignes un avis au comité de gestion : l'eau est-elle conforme ? Quel danger représente le réservoir ? Le placement est-il rentable par rapport aux risques ?

\courssub{Ressources à mobiliser}

\begin{aretenir}[Boîte à outils de la leçon]
\begin{itemize}
\item \textbf{Lien $\log$--$\ln$ :} $\log x = \dfrac{\ln x}{\ln 10}$ pour tout $x > 0$.
\item \textbf{Propriétés algébriques de $\ln$ :} $\ln(ab) = \ln a + \ln b$ ; $\ln\left(\dfrac{a}{b}\right) = \ln a - \ln b$ ; $\ln(a^n) = n\ln a$ ($n \in \mathbb{Z}$) ; $\ln(\sqrt{a}) = \frac{1}{2}\ln a$ ($a > 0$).
\item \textbf{Réciproque et monotonie :} pour $a > 0$ et $b > 0$, $\ln a = \ln b \iff a = b$ ; et $\ln a < \ln b \iff a < b$ (stricte croissance de $\ln$ sur $]0;+\infty[$).
\item \textbf{Équations exponentielles :} $e^{x} = b$ (avec $b > 0$) $\iff x = \ln b$. Plus généralement $e^{at} = b \iff t = \dfrac{\ln b}{a}$ ($a \neq 0$).
\item \textbf{Inéquations exponentielles :} si $a > 0$, $e^{at} > b \iff t > \dfrac{\ln b}{a}$ ; si $a < 0$, le sens change. De même $(1{,}06)^n \geq 2 \iff n\ln(1{,}06) \geq \ln 2$ car $\ln(1{,}06) > 0$.
\item \textbf{Condition d'existence :} $\ln x$ n'existe que si $x > 0$ (ensemble de définition $D_{\ln} = ]0;+\infty[$).
\item \textbf{Primitive de la forme $u'/u$ :} Si $u$ est dérivable et strictement positive sur un intervalle $I$, une primitive de $\dfrac{u'}{u}$ sur $I$ est $\ln \circ u$ (démonstration exigible au Bac).
\end{itemize}
\end{aretenir}

\courssub{Démarche et corrigé commenté}

\begin{exoresolu}[Corrigé complet de l'activité]
\textbf{Tâche 1.} 1) Exprimer le pH avec $\ln$. 2) Calculer $\left[\mathrm{H}^{+}\right]$ pour $\mathrm{pH} = 6{,}8$. 3) Encadrer $\left[\mathrm{H}^{+}\right]$ selon la norme et conclure.

\textbf{Tâche 2.} 1) Résoudre $e^{0{,}3t} = 2$. 2) Résoudre $N(t) > 1000\,N_0$. 3) Comparer les durées.

\textbf{Tâche 3.} 1) Résoudre $(1{,}06)^n \geq 2$. 2) Temps pour atteindre $4\,000\,000$ FCFA. 3) Comparaison et conclusion.
\tcblower
\textbf{Tâche 1 : le pH de l'eau.}

\textbf{1)} Par définition, $\mathrm{pH} = -\log\left[\mathrm{H}^{+}\right]$. Or $\log x = \dfrac{\ln x}{\ln 10}$ pour tout $x > 0$, donc en appliquant à $x = \left[\mathrm{H}^{+}\right] > 0$ :
\[ \mathrm{pH} = -\frac{\ln\left[\mathrm{H}^{+}\right]}{\ln 10}. \]

\textbf{2)} On résout l'équation $-\dfrac{\ln\left[\mathrm{H}^{+}\right]}{\ln 10} = 6{,}8$, d'inconnue $\left[\mathrm{H}^{+}\right] > 0$ (condition d'existence satisfaite car une concentration est strictement positive) :
\begin{align*}
\ln\left[\mathrm{H}^{+}\right] &= -6{,}8\,\ln 10 \\
\left[\mathrm{H}^{+}\right] &= e^{-6{,}8\,\ln 10} = \left(e^{\ln 10}\right)^{-6{,}8} = 10^{-6{,}8}.
\end{align*}
Valeur approchée : $10^{-6{,}8} = 10^{-7} \times 10^{0{,}2} \approx 1{,}58 \times 10^{-7}\ \text{mol/L}$.
\emph{(On peut aussi écrire $\left[\mathrm{H}^{+}\right] = \num{1,58e-7}\ \text{mol/L}$.)}

\textbf{3)} La norme s'écrit $6{,}5 \leq -\dfrac{\ln\left[\mathrm{H}^{+}\right]}{\ln 10} \leq 8{,}5$.
Comme $\ln 10 > 0$, le nombre $-\ln 10$ est strictement négatif. En multipliant les trois membres par $-\ln 10$, \textbf{le sens des inégalités change} :
\[ -8{,}5\,\ln 10 \leq \ln\left[\mathrm{H}^{+}\right] \leq -6{,}5\,\ln 10. \]
La fonction exponentielle est strictement croissante sur $\mathbb{R}$, on peut l'appliquer sans changer le sens :
\[ e^{-8{,}5\,\ln 10} \leq \left[\mathrm{H}^{+}\right] \leq e^{-6{,}5\,\ln 10} \iff 10^{-8{,}5} \leq \left[\mathrm{H}^{+}\right] \leq 10^{-6{,}5}. \]
Numériquement : $10^{-8{,}5} \approx 3{,}16 \times 10^{-9}$ et $10^{-6{,}5} \approx 3{,}16 \times 10^{-7}$.
L'encadrement est donc : $\boxed{3{,}16 \times 10^{-9} \leq \left[\mathrm{H}^{+}\right] \leq 3{,}16 \times 10^{-7}\ \text{mol/L}}$.
Comme $1{,}58 \times 10^{-7} \in \left[3{,}16 \times 10^{-9}\,;\,3{,}16 \times 10^{-7}\right]$, \textbf{l'eau du forage est conforme à la norme de potabilité}.

\textbf{Tâche 2 : la croissance bactérienne.}

\textbf{1)} L'équation $e^{0{,}3t} = 2$ admet une solution unique car $t \mapsto e^{0{,}3t}$ est bijective de $\mathbb{R}$ vers $]0;+\infty[$.
\begin{align*}
e^{0{,}3t} = 2 &\iff 0{,}3t = \ln 2 \quad (\text{car } \ln \text{ est la réciproque de } \exp) \\
&\iff t = \frac{\ln 2}{0{,}3}.
\end{align*}
Calcul : $\ln 2 \approx 0{,}693147$, donc $t \approx \dfrac{0{,}693147}{0{,}3} \approx 2{,}31049$ h.
Conversion en minutes : $0{,}31049 \times 60 \approx 18{,}6$ min $\approx 19$ min.
Le temps de doublement est \textbf{2 h 19 min} (arrondi à la minute près).

\textbf{2)} $N(t) > 1000\,N_0 \iff N_0\,e^{0{,}3t} > 1000\,N_0$.
Comme $N_0 > 0$, on divise par $N_0$ sans changer le sens : $e^{0{,}3t} > 1000$.
La fonction $\ln$ est strictement croissante sur $]0;+\infty[$, donc :
\[ 0{,}3t > \ln 1000 \iff t > \frac{\ln 1000}{0{,}3}. \]
Or $\ln 1000 = \ln(10^3) = 3\ln 10$, donc $t > \dfrac{3\ln 10}{0{,}3} = 10\ln 10$.
Calcul : $\ln 10 \approx 2{,}302585$, donc $t > 23{,}02585$ h.
L'eau devient impropre à la consommation \textbf{au bout de 24 heures entières} (dès le début de la $24^{\text{e}}$ heure, le seuil est dépassé ; on arrondit \emph{au-dessus} pour un seuil de sécurité).

\textbf{3)} Durée entre le premier doublement ($t_2 \approx 2{,}31$ h) et le dépassement du seuil ($t_{\text{seuil}} \approx 23{,}03$ h) :
$\Delta t \approx 23{,}03 - 2{,}31 = 20{,}72$ h, soit environ \textbf{20 h 43 min}.
\textbf{Commentaire :} On note que $1000 \approx 2^{10}$ (car $2^{10} = 1024$). Atteindre $1000\,N_0$ correspond donc à environ $10$ doublements successifs. $10 \times 2{,}31 \approx 23{,}1$ h, ce qui est cohérent avec $23{,}03$ h. La croissance exponentielle est fulgurante : le réservoir doit être nettoyé \textbf{quotidiennement}.

\textbf{Tâche 3 : le placement.}

\textbf{1)} Le capital double lorsque $C(n) \geq 2 \times 2\,000\,000$, c'est-à-dire $(1{,}06)^n \geq 2$.
La fonction $n \mapsto (1{,}06)^n$ est strictement croissante. On passe au $\ln$ (qui est strictement croissante) : comme $\ln(1{,}06) > 0$, le sens de l'inégalité est conservé.
\begin{align*}
(1{,}06)^n \geq 2 &\iff \ln\left((1{,}06)^n\right) \geq \ln 2 \\
&\iff n\ln(1{,}06) \geq \ln 2 \\
&\iff n \geq \frac{\ln 2}{\ln(1{,}06)}.
\end{align*}
Calcul : $\ln 2 \approx 0{,}693147$, $\ln(1{,}06) \approx 0{,}0582689$, rapport $\approx 11{,}8957$.
$n$ étant un entier (nombre d'années), la plus petite solution est $n = 12$.
\textbf{Le capital double au bout de 12 ans.}

\textbf{2)} Atteindre $4\,000\,000$ FCFA, c'est quadrupler le capital initial : $(1{,}06)^n \geq 4 = 2^2$.
\[ n \geq \frac{\ln 4}{\ln(1{,}06)} = \frac{2\ln 2}{\ln(1{,}06)} \approx 2 \times 11{,}8957 \approx 23{,}79. \]
La plus petite solution entière est $n = 24$.
\textbf{Vérification :} $24 = 2 \times 12$. Quadrupler demande deux doublements, donc deux fois le temps de doublement : cohérent.

\textbf{3)} Les bactéries doublent en $2{,}31$ \textbf{heures}, le capital en $12$ \textbf{années} ($\approx 105\,120$ heures).
Le rapport des vitesses est vertigineux : en une seule journée ($24$ h), la population bactérienne est multipliée par $2^{10} \approx 1000$, tandis que le capital ne croît que de $6\,\%$ par an.
\textbf{Conclusion :} Le placement est sûr mais lent ; en revanche, le risque sanitaire évolue à une vitesse incomparable et exige une surveillance quotidienne du réservoir.

\textbf{Synthèse commune (exemple de rédaction).}
\og L'eau du forage est conforme à la norme de potabilité (pH $= 6{,}8 \in [6{,}5\,;\,8{,}5]$). Cependant, la population bactérienne du réservoir double toutes les $2$ h $19$ min et rend l'eau impropre en moins de $24$ h : un nettoyage quotidien s'impose. Le placement à $6\,\%$ double en $12$ ans ; il est rentable à long terme mais ne doit pas faire oublier l'urgence sanitaire. \fg
\end{exoresolu}

\courssub{Bilan de l'activité}

Cette activité a permis de mettre en évidence plusieurs acquis essentiels de la leçon :

\begin{itemize}
\item \textbf{Le passage au logarithme népérien est la clé} de toute équation ou inéquation exponentielle : $e^{at} = b \iff t = \dfrac{\ln b}{a}$, et le signe du coefficient $a$ (ou celui de $\ln$ de la base) détermine le sens de l'inégalité.
\item \textbf{Les propriétés algébriques de $\ln$} simplifient les calculs : $\ln 1000 = 3\ln 10$, $\ln(2^2) = 2\ln 2$, et la conversion $\log x = \dfrac{\ln x}{\ln 10}$ relie les deux logarithmes.
\item \textbf{La stricte croissance de $\ln$} garantit l'équivalence des inégalités lors du passage au logarithme, à condition de vérifier que les quantités sont strictement positives (conditions d'existence).
\item \textbf{Le temps de doublement} est un concept transversal : il vaut $\dfrac{\ln 2}{a}$ pour une croissance $e^{at}$ ($a>0$), et $\dfrac{\ln 2}{\ln(1+t)}$ pour des intérêts composés au taux $t$. Il ne dépend pas de la quantité initiale.
\item \textbf{L'interprétation des résultats} est une compétence à part entière : arrondir correctement (à l'entier supérieur pour un seuil à atteindre), convertir les unités (heures, minutes, années) et formuler une recommandation argumentée sont des savoir-faire attendus au Baccalauréat comme dans la vie professionnelle.
\end{itemize}

\begin{attention}[Pièges rencontrés pendant l'activité]
\begin{itemize}
\item Ne pas oublier que multiplier une inégalité par $-\ln 10$ (nombre négatif) en change le sens.
\item Ne pas écrire $\ln(1000\,N_0) = \ln 1000 \times \ln N_0$ : c'est $\ln 1000 + \ln N_0$ (propriété $\ln(ab)=\ln a+\ln b$).
\item Pour $n \geq 11{,}90$ avec $n$ entier, la réponse est $n = 12$ et non $n = 11$ : on arrondit toujours \emph{au-dessus} pour un seuil à atteindre ou dépasser.
\item Vérifier systématiquement la condition d'existence $x > 0$ avant d'appliquer $\ln x$ (ici concentrations, effectifs, capital sont $>0$).
\end{itemize}
\end{attention}

\begin{savaistu}[Le pH et le logarithme népérien dans la vie courante]
L'échelle de pH, inventée par le chimiste danois Sørensen en 1909, est une échelle \textbf{logarithmique} : une variation de 1 unité de pH correspond à une multiplication par 10 de la concentration en ions $\mathrm{H}^+$. C'est pourquoi on utilise le logarithme (décimal ou népérien) pour manipuler des grandeurs qui varient sur plusieurs ordres de grandeur : échelle de Richter (séismes), décibels (son), magnitude (astronomie). Au Cameroun, le contrôle de la qualité de l'eau (MINSANTE, CAMWATER) repose sur ces mesures logarithmiques.
\end{savaistu}

\newpage

\coursec{Méthodes et exercices résolus}

\coursec{Fonction logarithme népérien}

\courssub{Définition et premières propriétés}

\begin{definition}[Fonction logarithme népérien]
On appelle \textbf{logarithme népérien} la fonction $\ln$ définie sur $]0\,;+\infty[$ par :
\[ \ln x = \int_1^x \frac{\mathrm{d}t}{t}. \]
C'est l'unique fonction dérivable sur $]0\,;+\infty[$ telle que $\ln' x = \dfrac{1}{x}$ et $\ln 1 = 0$.
Le nombre $e$ (constante d'Euler) est l'unique réel tel que $\ln e = 1$ ; on a $e \approx 2,718$.
\end{definition}

\begin{propriete}[Propriétés fondamentales de $\ln$]
\begin{itemize}
\item \textbf{Ensemble de définition :} $\mathcal{D}_{\ln} = ]0\,;+\infty[$.
\item \textbf{Dérivée :} $\forall x>0,\quad \ln' x = \dfrac{1}{x} > 0$.
\item \textbf{Variations :} $\ln$ est \cle{strictement croissante} sur $]0\,;+\infty[$.
\item \textbf{Limites aux bornes :}
\[ \lim_{x\to 0^+} \ln x = -\infty \qquad\text{et}\qquad \lim_{x\to +\infty} \ln x = +\infty. \]
\item \textbf{Valeurs remarquables :} $\ln 1 = 0$ ; $\ln e = 1$ ; $\ln(e^a) = a$ ; $e^{\ln a} = a\ (a>0)$.
\item \textbf{Courbe représentative :} Elle passe par $(1,0)$ et $(e,1)$, est strictement croissante, convexe, avec l'axe des ordonnées ($x=0$) comme asymptote verticale.
\end{itemize}
\end{propriete}

\begin{popfigure}
\begin{tikzpicture}[>=Stealth]
\begin{axis}[
    axis lines=middle,
    xlabel=$x$, ylabel=$y$,
    xmin=-0.5, xmax=5,
    ymin=-2.5, ymax=2,
    xtick={1,2.718}, xticklabels={$1$,$e$},
    ytick={-1,1}, yticklabels={$-1$,$1$},
    domain=0.01:5, samples=200,
    restrict y to domain=-2.5:2,
    width=\linewidth, height=0.6\linewidth,
    axis line style={popdark},
    tick style={popdark},
    xlabel style={anchor=north west},
    ylabel style={anchor=south east},
]
\addplot[popcurve, thick] {ln(x)};
\draw[popmuted, dashed] (axis cs:0,-2.5) -- (axis cs:0,2);
\node[popdark, anchor=south west] at (axis cs:1,0) {$(1,0)$};
\node[popdark, anchor=south west] at (axis cs:2.718,1) {$(e,1)$};
\end{axis}
\end{tikzpicture}
\legende{Courbe représentative de la fonction logarithme népérien $y=\ln x$.}
\end{popfigure}

\begin{experience}[Découverte des propriétés algébriques (démonstration guidée)]
Soient $a>0$ et $b>0$. On veut montrer $\ln(ab)=\ln a+\ln b$.
\begin{enumerate}
\item Poser $f(x)=\ln(ax)-\ln x$ sur $]0\,;+\infty[$. Calculer $f'(x)$.
\item En déduire que $f$ est constante. Calculer cette constante en $x=1$.
\item Conclure pour $\ln(ab)$. Adapter la méthode pour $\ln(a/b)$, $\ln(a^n)$, $\ln(\sqrt{a})$.
\end{enumerate}
\end{experience}

\begin{aretenir}[Propriétés algébriques de $\ln$]
Pour tout $a>0$, $b>0$ et tout $n\in\mathbb{Z}$ (et même $n\in\mathbb{Q}$ pour les racines) :
\[
\ln(ab)=\ln a+\ln b \quad;\quad
\ln\!\left(\frac{a}{b}\right)=\ln a-\ln b \quad;\quad
\ln(a^n)=n\ln a \quad;\quad
\ln(\sqrt{a})=\frac{1}{2}\ln a.
\]
\textbf{Conséquences immédiates :} $\ln 1=0$ ; $\ln e=1$ ; $\ln(e^a)=a$ ; $e^{\ln a}=a\ (a>0)$.
\end{aretenir}

\begin{exemplebox}[Calculs directs]
$\ln 4 = 2\ln 2$ ; $\ln\sqrt{e} = \frac{1}{2}$ ; $\ln\frac{1}{e^3} = -3$ ; $\ln(\frac{e}{2}) = 1-\ln 2$.
\end{exemplebox}

\courssub{Propriétés algébriques de $\ln$}

\textbf{Méthode : Simplifier une expression contenant des logarithmes}\par
\begin{methode}[Simplifier une expression contenant des logarithmes]
\begin{enumerate}
\item Repérer les \cle{produits}, \cle{quotients}, \cle{puissances} et \cle{racines carrées} à l'intérieur des logarithmes.
\item Appliquer les formules fondamentales, valables pour $a>0$, $b>0$ et $n$ entier :
\[ \ln(ab)=\ln a+\ln b \;;\qquad \ln\!\left(\frac{a}{b}\right)=\ln a-\ln b \;;\qquad \ln(a^n)=n\ln a \;;\qquad \ln(\sqrt{a})=\frac{1}{2}\ln a. \]
\item Factoriser $\ln 2$, $\ln 3$, etc., pour regrouper les termes semblables.
\item Conclure en donnant le résultat sous la forme demandée (souvent $a\ln 2+b\ln 3$ avec $a,b\in\mathbb{Z}$).
\end{enumerate}
\textbf{Réflexes :} $\ln 1=0$ ; $\ln e=1$ ; pour aller dans l'autre sens (regrouper), $p\ln a+q\ln b=\ln(a^p b^q)$.
\end{methode}

\begin{exoresolu}[Simplification d'une expression logarithmique (type Bac)]
On pose $A=\ln 48-2\ln 6+\ln 3$ et $B=\ln\!\left(\dfrac{9}{4}\right)+\ln\!\left(\dfrac{8}{3}\right)-\ln 6$.\\
Écrire $A$ et $B$ sous la forme $a\ln 2+b\ln 3$, où $a$ et $b$ sont des entiers relatifs.
\tcblower
\textbf{Calcul de $A$.} On décompose chaque nombre en facteurs premiers : $48=2^4\times 3$ et $6=2\times 3$.
\begin{align*}
A &= \ln(2^4\times 3)-2\ln(2\times 3)+\ln 3 \\
&= 4\ln 2+\ln 3-2(\ln 2+\ln 3)+\ln 3 \\
&= 4\ln 2+\ln 3-2\ln 2-2\ln 3+\ln 3 \\
&= 2\ln 2+0\times\ln 3 = 2\ln 2.
\end{align*}
\textbf{Calcul de $B$.} On utilise $\ln\!\left(\dfrac{a}{b}\right)=\ln a-\ln b$ :
\begin{align*}
B &= \ln 9-\ln 4+\ln 8-\ln 3-\ln 6 \\
&= 2\ln 3-2\ln 2+3\ln 2-\ln 3-(\ln 2+\ln 3) \\
&= 2\ln 3-2\ln 2+3\ln 2-\ln 3-\ln 2-\ln 3 \\
&= 0\times\ln 2+0\times\ln 3 = 0.
\end{align*}
\textbf{Conclusion :} $A=2\ln 2$ (c'est-à-dire $A=\ln 4$) et $B=0$ (c'est-à-dire que l'argument vaut $1$).
\end{exoresolu}

\courssub{Limites de la fonction $\ln$}

\begin{propriete}[Limites usuelles et croissances comparées]
\begin{itemize}
\item $\displaystyle\lim_{x\to 0^+} \ln x = -\infty$ ; $\displaystyle\lim_{x\to +\infty} \ln x = +\infty$.
\item \textbf{Croissances comparées (à connaître) :}
\[ \lim_{x\to +\infty} \frac{\ln x}{x} = 0 \quad;\quad \lim_{x\to 0^+} x\ln x = 0 \quad;\quad \lim_{x\to +\infty} \frac{\ln x}{x^n} = 0\ (n\geqslant 1). \]
\item \textbf{Composées :} Si $\lim u(x) = \ell \in ]0\,;+\infty[$, alors $\lim \ln(u(x)) = \ln \ell$ (continuité).
Si $\lim u(x) = 0^+$, alors $\lim \ln(u(x)) = -\infty$ ; si $\lim u(x) = +\infty$, alors $\lim \ln(u(x)) = +\infty$.
\item \textbf{Limite remarquable en $1$ :} $\displaystyle\lim_{x\to 1} \frac{\ln x}{x-1} = 1$ (taux d'accroissement de $\ln$ en $1$).
\end{itemize}
\end{propriete}

\textbf{Méthode : Calculer les limites d'une fonction contenant $\ln$}\par
\begin{methode}[Lever les indéterminations avec $\ln$]
\begin{enumerate}
\item Identifier la forme : en $+\infty$, $\ln x\to+\infty$ ; en $0^+$, $\ln x\to-\infty$.
\item Formes indéterminées $\infty-\infty$ ou $\dfrac{\infty}{\infty}$ : \cle{factoriser par le terme dominant} (souvent $x$ ou $\ln x$) puis utiliser les limites de croissances comparées.
\item Pour les \cle{composées} : poser $X=u(x)$, déterminer la limite de $u$, puis appliquer la limite de $\ln$ en ce point.
\item Limite de référence en $1$ : $\displaystyle\lim_{x\to 1}\frac{\ln x}{x-1}=1$.
\end{enumerate}
\end{methode}

\begin{exoresolu}[Calculs de limites (type Bac)]
Calculer les limites suivantes :
\[ \text{a) }\lim_{x\to+\infty}\big(x-\ln x\big) \;;\qquad \text{b) }\lim_{x\to 0^+}\big(x\ln x-x\big) \;;\qquad \text{c) }\lim_{x\to+\infty}\ln\!\left(\frac{2x+1}{x-3}\right). \]
\tcblower
\textbf{a)} Forme indéterminée $\infty-\infty$. On factorise par le terme dominant $x$ :
\[ x-\ln x=x\left(1-\frac{\ln x}{x}\right). \]
Or $\displaystyle\lim_{x\to+\infty}\frac{\ln x}{x}=0$ (croissance comparée), donc $\displaystyle\lim_{x\to+\infty}\left(1-\frac{\ln x}{x}\right)=1$ et $\displaystyle\lim_{x\to+\infty}x=+\infty$.
Par produit : $\displaystyle\lim_{x\to+\infty}(x-\ln x)=+\infty$.

\textbf{b)} On sait que $\displaystyle\lim_{x\to 0^+}x\ln x=0$ (croissance comparée) et $\displaystyle\lim_{x\to 0^+}x=0$.
Par somme : $\displaystyle\lim_{x\to 0^+}(x\ln x-x)=0$.

\textbf{c)} Posons $u(x)=\dfrac{2x+1}{x-3}$. Pour $x\to+\infty$ :
\[ u(x)=\frac{x\left(2+\frac{1}{x}\right)}{x\left(1-\frac{3}{x}\right)}=\frac{2+\frac{1}{x}}{1-\frac{3}{x}}\xrightarrow[x\to+\infty]{}2. \]
La fonction $\ln$ étant continue en $2$ (avec $2>0$), par composition :
\[ \lim_{x\to+\infty}\ln\!\left(\frac{2x+1}{x-3}\right)=\ln 2. \]
\textbf{Conclusion :} a) $+\infty$ ; b) $0$ ; c) $\ln 2$.
\end{exoresolu}

\courssub{Dérivation, variations et courbe représentative}

\begin{propriete}[Dérivée de $\ln u$]
Soit $u$ une fonction dérivable et strictement positive sur un intervalle $I$.
La fonction $\ln u$ est dérivable sur $I$ et :
\[ (\ln u)' = \frac{u'}{u}. \]
\textbf{Preuve :} Dérivée de la composée $x \mapsto \ln(u(x))$ : $\ln'(u(x)) \times u'(x) = \frac{1}{u(x)} \times u'(x)$.
\end{propriete}

\textbf{Méthode : Étudier une fonction contenant $\ln u$ et tracer sa courbe}\par
\begin{methode}[Étude complète de $f(x)$ contenant $\ln u(x)$]
\begin{enumerate}
\item \textbf{Ensemble de définition :} résoudre $u(x)>0$.
\item \textbf{Limites aux bornes} de $\mathcal{D}_f$ ; en déduire les éventuelles \cle{asymptotes} (verticale si limite infinie en un bord fini, horizontale si limite finie en $\pm\infty$).
\item \textbf{Dérivée :} si $f=\ln u$, alors $f'=\dfrac{u'}{u}$ ; si $f$ est une somme ou un produit, dériver terme à terme en utilisant $(\ln x)'=\dfrac{1}{x}$.
\item \textbf{Signe de $f'$ :} attention, sur $\mathcal{D}_f$ on a $u(x)>0$, donc le signe de $\dfrac{u'}{u}$ est celui de $u'$. Dresser le tableau de variations.
\item \textbf{Points remarquables :} intersection avec l'axe des abscisses ($f(x)=0$), valeur en $1$ ou en $e$, tangente éventuelle.
\item \textbf{Tracer la courbe} en plaçant les points et en respectant les asymptotes.
\end{enumerate}
\end{methode}

\begin{exoresolu}[Étude de $f(x)=x-\ln x$ (type Bac)]
Soit $f$ la fonction définie par $f(x)=x-\ln x$ et $\mathcal{C}_f$ sa courbe dans un repère orthonormé.
\begin{enumerate}
\item Déterminer l'ensemble de définition de $f$ et les limites aux bornes.
\item Étudier les variations de $f$ et dresser son tableau de variations.
\item Montrer que l'équation $f(x)=0$ n'admet aucune solution.
\end{enumerate}
\tcblower
\textbf{1. Ensemble de définition et limites.} $\ln x$ n'existe que pour $x>0$, donc $\mathcal{D}_f=\,]0\,;+\infty[$.

En $0^+$ : $\displaystyle\lim_{x\to 0^+}x=0$ et $\displaystyle\lim_{x\to 0^+}\ln x=-\infty$, donc $\displaystyle\lim_{x\to 0^+}f(x)=+\infty$ : la droite $x=0$ (axe des ordonnées) est \textbf{asymptote verticale} à $\mathcal{C}_f$.

En $+\infty$ : forme indéterminée ; on factorise : $f(x)=x\left(1-\dfrac{\ln x}{x}\right)$. Comme $\displaystyle\lim_{x\to+\infty}\dfrac{\ln x}{x}=0$, on obtient $\displaystyle\lim_{x\to+\infty}f(x)=+\infty$.

\textbf{2. Variations.} $f$ est dérivable sur $]0\,;+\infty[$ comme somme de fonctions dérivables :
\[ f'(x)=1-\frac{1}{x}=\frac{x-1}{x}. \]
Sur $]0\,;+\infty[$, $x>0$, donc $f'(x)$ a le signe de $x-1$ : $f'(x)<0$ sur $]0\,;1[$, $f'(1)=0$, $f'(x)>0$ sur $]1\,;+\infty[$.

Minimum : $f(1)=1-\ln 1=1$.

\begin{center}
\begin{tabular}{|c|ccccc|}
\hline
$x$ & $0$ & & $1$ & & $+\infty$ \\
\hline
$f'(x)$ & & $-$ & $0$ & $+$ & \\
\hline
$f(x)$ & $+\infty$ & $\searrow$ & $1$ & $\nearrow$ & $+\infty$ \\
\hline
\end{tabular}
\end{center}

\textbf{3. Équation $f(x)=0$.} D'après le tableau de variations, $f$ admet sur $]0\,;+\infty[$ un minimum égal à $1$. Donc pour tout $x>0$ : $f(x)\geqslant 1>0$.

\textbf{Conclusion :} l'équation $f(x)=0$ n'admet aucune solution dans $]0\,;+\infty[$. (On a au passage démontré l'inégalité classique $\ln x\leqslant x-1$ pour tout $x>0$.)
\end{exoresolu}

\begin{popfigure}
\begin{tikzpicture}[>=Stealth]
\begin{axis}[
    axis lines=middle,
    xlabel=$x$, ylabel=$y$,
    xmin=-0.5, xmax=4,
    ymin=-0.5, ymax=4,
    xtick={1}, xticklabels={$1$},
    ytick={1}, yticklabels={$1$},
    domain=0.01:4, samples=200,
    width=\linewidth, height=0.6\linewidth,
]
\addplot[popcurve, thick] {x - ln(x)};
\draw[popmuted, dashed] (axis cs:0,0) -- (axis cs:0,4);
\node[popdark, anchor=south east] at (axis cs:1,1) {$(1,1)$};
\end{axis}
\end{tikzpicture}
\legende{Courbe de $f(x)=x-\ln x$ : minimum en $(1,1)$, asymptote verticale $x=0$.}
\end{popfigure}

\courssub{Primitives de la forme $\dfrac{u'}{u}$}

\begin{propriete}[Primitive de $\dfrac{u'}{u}$]
Soit $u$ une fonction dérivable ne s'annulant pas sur un intervalle $I$.
Une primitive de $\dfrac{u'}{u}$ sur $I$ est $\ln|u|$.
Si de plus $u>0$ sur $I$, une primitive est $\ln(u)$.
\[ \int \frac{u'(x)}{u(x)}\,\mathrm{d}x = \ln|u(x)| + C \quad\text{sur } I. \]
\end{propriete}

\textbf{Méthode : Déterminer une primitive de la forme $\dfrac{u'}{u}$}\par
\begin{methode}[Primitives et logarithmes]
\begin{enumerate}
\item Mettre le numérateur sous la forme $k\times u'(x)$ en identifiant la fonction $u$ du dénominateur.
\item Appliquer la formule : si $u$ est dérivable et \textbf{ne s'annule pas} sur l'intervalle $I$, alors
\[ \int \frac{u'(x)}{u(x)}\,\mathrm{d}x=\ln|u(x)|+C \quad\text{sur } I. \]
Sur un intervalle où $u>0$, on écrit simplement $\ln\big(u(x)\big)+C$.
\item Si le numérateur est $k\,u'$, sortir le facteur $k$ : $\displaystyle\int k\frac{u'}{u}=k\ln|u|+C$.
\item \textbf{Vérifier} en dérivant le résultat obtenu.
\end{enumerate}
\textbf{Réflexe :} $\displaystyle\int\frac{\mathrm{d}x}{x}=\ln x+C$ sur $]0\,;+\infty[$ ; $\displaystyle\int_1^{e}\frac{\mathrm{d}x}{x}=1$.
\end{methode}

\begin{exoresolu}[Primitives et calcul intégral (type Bac)]
\begin{enumerate}
\item Déterminer une primitive $F$ de $f(x)=\dfrac{1}{x\ln x}$ sur $]1\,;+\infty[$.
\item Déterminer la primitive $G$ de $g(x)=\dfrac{2x+3}{x^2+3x+1}$ sur $[0\,;+\infty[$ qui s'annule en $0$.
\end{enumerate}
\tcblower
\textbf{1.} Posons $u(x)=\ln x$. Alors $u'(x)=\dfrac{1}{x}$, et $f(x)=\dfrac{u'(x)}{u(x)}$.

Sur $]1\,;+\infty[$, $\ln x>0$, donc $u$ est strictement positive : une primitive est
\[ F(x)=\ln\big(\ln x\big). \]
\textbf{Vérification :} $F'(x)=\dfrac{\frac{1}{x}}{\ln x}=\dfrac{1}{x\ln x}=f(x)$. \checkmark

\textbf{2.} Posons $u(x)=x^2+3x+1$. Alors $u'(x)=2x+3$ : le numérateur est exactement $u'$, donc $g=\dfrac{u'}{u}$.

Sur $[0\,;+\infty[$, $u(x)=x^2+3x+1\geqslant 1>0$ (somme de termes positifs, $u(0)=1$). Les primitives de $g$ sont donc :
\[ G(x)=\ln(x^2+3x+1)+C. \]
Condition $G(0)=0$ : $\ln 1+C=0$, soit $C=0$.

\textbf{Conclusion :} $G(x)=\ln(x^2+3x+1)$.

\textbf{Vérification :} $G'(x)=\dfrac{2x+3}{x^2+3x+1}=g(x)$ et $G(0)=\ln 1=0$. \checkmark
\end{exoresolu}

\courssub{Équations, inéquations, systèmes et applications}

\textbf{Méthode : Résoudre une équation ou inéquation logarithmique}\par
\begin{methode}[Résoudre une équation ou inéquation logarithmique]
\begin{enumerate}
\item \textbf{Étape indispensable :} déterminer l'\cle{ensemble de définition} en imposant que chaque argument de $\ln$ soit \textbf{strictement positif}.
\item Se ramener, grâce aux propriétés algébriques, à l'une des formes :
\begin{itemize}
\item $\ln u=\ln v \iff u=v$ (le logarithme est \cle{injectif}) ;
\item $\ln u=k \iff u=e^{k}$ ;
\item $\ln u<\ln v \iff u<v$ car $\ln$ est \cle{strictement croissante} sur $]0\,;+\infty[$ : le sens de l'inégalité est \textbf{conservé}.
\end{itemize}
\item Résoudre l'équation ou l'inéquation obtenue.
\item \textbf{Ne pas oublier :} confronter les solutions trouvées à l'ensemble de définition et rejeter celles qui n'y appartiennent pas.
\end{enumerate}
\textbf{Réflexe :} $\ln x>0 \iff x>1$ et $\ln x<0 \iff 0<x<1$.
\end{methode}

\begin{exoresolu}[Équation et inéquation logarithmiques (type Bac)]
\begin{enumerate}
\item Résoudre dans $\mathbb{R}$ l'équation $(E)$ : $\ln(x+1)+\ln(x-1)=\ln 3$.
\item Résoudre dans $\mathbb{R}$ l'inéquation $(I)$ : $\ln(2x-3)\leqslant \ln(x+1)$.
\end{enumerate}
\tcblower
\textbf{1. Résolution de $(E)$.}

\textbf{Conditions d'existence :} il faut $x+1>0$ et $x-1>0$, c'est-à-dire $x>-1$ et $x>1$. L'ensemble de définition est donc $\mathcal{D}=\,]1\,;+\infty[$.

Sur $\mathcal{D}$, on regroupe : $\ln(x+1)+\ln(x-1)=\ln\big[(x+1)(x-1)\big]=\ln(x^2-1)$.
L'équation devient $\ln(x^2-1)=\ln 3$, d'où, par injectivité de $\ln$ :
\[ x^2-1=3 \iff x^2=4 \iff x=2 \ \text{ou}\ x=-2. \]
\textbf{Confrontation à $\mathcal{D}$ :} $-2\notin\mathcal{D}$ est rejeté ; $2\in\mathcal{D}$ est accepté.

\textbf{Conclusion :} $\mathcal{S}_E=\{2\}$.

\textbf{2. Résolution de $(I)$.}

\textbf{Conditions :} $2x-3>0$ et $x+1>0$, soit $x>\dfrac{3}{2}$. Donc $\mathcal{D}=\left]\dfrac{3}{2}\,;+\infty\right[$.

Sur $\mathcal{D}$, $\ln$ étant strictement croissante :
\[ \ln(2x-3)\leqslant\ln(x+1) \iff 2x-3\leqslant x+1 \iff x\leqslant 4. \]
On intersecte avec $\mathcal{D}$ : $x\in\left]\dfrac{3}{2}\,;4\right]$.

\textbf{Conclusion :} $\mathcal{S}_I=\left]\dfrac{3}{2}\,;4\right]$.
\end{exoresolu}

\textbf{Méthode : Résoudre un système d'inconnue $\ln x$ (changement de variable)}\par
\begin{methode}[Poser $X=\ln x$]
\begin{enumerate}
\item Vérifier la condition $x>0$ (ou les conditions sur les arguments).
\item Poser le \cle{changement de variable} $X=\ln x$ (et $Y=\ln y$ pour un système) pour obtenir une équation polynomiale ou un système linéaire en $X$ (et $Y$).
\item Résoudre dans $\mathbb{R}$ l'équation ou le système en $X$ (et $Y$).
\item Revenir à l'inconnue initiale par $x=e^{X}$ (et $y=e^{Y}$) : toute valeur réelle de $X$ donne une solution, car $e^X>0$ automatiquement.
\item Conclure en donnant les valeurs exactes, éventuellement une valeur approchée.
\end{enumerate}
\end{methode}

\begin{exoresolu}[Système logarithmique (type Bac)]
Résoudre dans $\mathbb{R}^2$ le système :
\[ \begin{cases} \ln x+\ln y=\ln 8 \\ 2\ln x-\ln y=\ln 2 \end{cases} \]
\tcblower
\textbf{Conditions :} $x>0$ et $y>0$.

\textbf{Changement de variable :} posons $X=\ln x$ et $Y=\ln y$. Le système devient :
\[ \begin{cases} X+Y=\ln 8 \\ 2X-Y=\ln 2 \end{cases} \]
\textbf{Résolution :} additionnons les deux lignes : $3X=\ln 8+\ln 2=\ln 16$, donc $X=\dfrac{\ln 16}{3}=\dfrac{4\ln 2}{3}$.

De la première ligne : $Y=\ln 8-X=3\ln 2-\dfrac{4\ln 2}{3}=\dfrac{9\ln 2-4\ln 2}{3}=\dfrac{5\ln 2}{3}$.

\textbf{Retour à $x$ et $y$ :}
\[ x=e^{X}=e^{\frac{4\ln 2}{3}}=2^{4/3}=\sqrt[3]{16} \;;\qquad y=e^{Y}=e^{\frac{5\ln 2}{3}}=2^{5/3}=\sqrt[3]{32}. \]
Ces valeurs sont strictement positives : elles conviennent.

\textbf{Vérification :} $\ln x+\ln y=\dfrac{4\ln 2}{3}+\dfrac{5\ln 2}{3}=3\ln 2=\ln 8$ \checkmark{} ; $2\ln x-\ln y=\dfrac{8\ln 2-5\ln 2}{3}=\ln 2$ \checkmark.

\textbf{Conclusion :} $\mathcal{S}=\left\{\left(\sqrt[3]{16}\,;\sqrt[3]{32}\right)\right\}$, soit approximativement $(2{,}52\,;\,3{,}17)$.
\end{exoresolu}

\textbf{Méthode : Modéliser et résoudre un problème de seuil (croissance/désintégration)}\par
\begin{methode}[Modéliser et résoudre un problème de seuil]
\begin{enumerate}
\item Traduire l'énoncé par une inéquation du type $a^n\geqslant b$ ou $A\,e^{kt}\geqslant B$.
\item Pour $a^n\geqslant b$ (avec $a>0$, $a\neq 1$, $b>0$) : appliquer $\ln$ aux deux membres ($\ln$ est croissante, le sens est conservé) : $n\ln a\geqslant\ln b$.
\item \textbf{Attention au signe de $\ln a$ :} si $a>1$, $\ln a>0$ et on divise sans changer le sens ; si $0<a<1$, $\ln a<0$ et la division \textbf{change le sens} de l'inégalité.
\item Conclure avec le plus petit entier $n$ convenable (arrondir à l'entier supérieur) et rédiger une phrase de conclusion dans le contexte.
\end{enumerate}
\end{methode}

\begin{exoresolu}[Croissance démographique à Yaoundé (type Bac)]
La population d'un quartier de Yaoundé est de \num{50000} habitants et augmente de $4\,\%$ par an. On modélise la population après $n$ années par $P(n)=\num{50000}\times(1{,}04)^n$.

À partir de combien d'années la population dépassera-t-elle \num{100000} habitants ?
\tcblower
On résout l'inéquation $P(n)>\num{100000}$, c'est-à-dire :
\[ \num{50000}\times(1{,}04)^n>\num{100000} \iff (1{,}04)^n>2. \]
La fonction $\ln$ étant strictement croissante, elle conserve le sens de l'inégalité :
\[ \ln\big((1{,}04)^n\big)>\ln 2 \iff n\ln(1{,}04)>\ln 2. \]
Comme $1{,}04>1$, on a $\ln(1{,}04)>0$ : on divise sans changer le sens :
\[ n>\frac{\ln 2}{\ln(1{,}04)}\approx\frac{0{,}6931}{0{,}0392}\approx 17{,}67. \]
Le plus petit entier strictement supérieur à $17{,}67$ est $n=18$.

\textbf{Vérification :} $(1{,}04)^{17}\approx 1{,}948<2$ et $(1{,}04)^{18}\approx 2{,}026>2$. \checkmark

\textbf{Conclusion :} c'est à partir de la $18^{\text{e}}$ année que la population du quartier dépassera \num{100000} habitants (elle aura alors doublé).
\end{exoresolu}

\begin{savaistu}[Le logarithme népérien : histoire et applications au Cameroun]
\begin{itemize}
\item \textbf{Histoire :} John Napier (Neper, 1550--1617), Écossais, invente les logarithmes pour simplifier les calculs astronomiques (produits $\to$ sommes). Henry Briggs (1561--1630) propose la base $10$. Leonhard Euler (1707--1783) impose la base $e$ et la notation $\ln$ (logarithmus naturalis).
\item \textbf{Échelle de Richter (séismes) :} Magnitude $M = \log_{10}(A/A_0)$. Une magnitude 6 libère 10 fois plus d'énergie qu'une magnitude 5.
\item \textbf{pH (chimie, eau potable) :} $\mathrm{pH} = -\log_{10}[\mathrm{H}^+]$. L'eau pure a pH $=7$. Au Cameroun, le contrôle de la qualité de l'eau (CAMWATER, forages ruraux) utilise cette échelle.
\item \textbf{Niveau sonore (décibels) :} $L = 10\log_{10}(I/I_0)$. Une conversation $\approx 60\,\mathrm{dB}$, un concert $\approx 110\,\mathrm{dB}$.
\item \textbf{Démographie et finance :} Modélisation de la croissance de la population (Douala, Yaoundé), calcul du temps de doublement, intérêts composés (banques : Afriland First Bank, UBA, BICEC).
\end{itemize}
\end{savaistu}

\begin{attention}[Les 5 erreurs qui coûtent des points au Bac]
\begin{enumerate}
\item \textbf{Oublier les conditions d'existence.} Avant toute résolution d'équation ou d'inéquation avec $\ln$, écris systématiquement que chaque argument est \emph{strictement positif}, et rejette à la fin les solutions hors de l'ensemble de définition. Une solution $x=-2$ acceptée pour $\ln(x-1)$ fait perdre tous les points de la question.
\item \textbf{Inventer des formules fausses.} $\ln(a+b)\neq\ln a+\ln b$ et $\dfrac{\ln a}{\ln b}\neq\ln a-\ln b$. Les seules formules valables portent sur le produit, le quotient, la puissance et la racine.
\item \textbf{Dériver $\ln u$ en $\dfrac{1}{u}$.} La dérivée de $\ln u$ est $\dfrac{u'}{u}$ : ne pas oublier le facteur $u'$. Par exemple $\big(\ln(x^2+1)\big)'=\dfrac{2x}{x^2+1}$ et non $\dfrac{1}{x^2+1}$.
\item \textbf{Changer le sens d'une inégalité sans raison... ou l'oublier quand il le faut.} Composer par $\ln$ conserve le sens ($\ln$ est croissante) ; mais diviser par $\ln a$ quand $0<a<1$ (donc $\ln a<0$) inverse le sens. Vérifie toujours le signe de la quantité par laquelle tu divises.
\item \textbf{Écrire $\ln|u|$ à tort ou oublier le domaine dans les primitives.} La primitive de $\dfrac{u'}{u}$ est $\ln|u|$ en général, mais sur un intervalle où $u>0$ on écrit $\ln u$. Précise toujours l'intervalle de validité et vérifie ta primitive en la dérivant.
\end{enumerate}
\end{attention}

\newpage

\coursec{Exercices}

\coursec{Fonction logarithme népérien}

\courssub{Définition et premières propriétés}

\begin{definition}[Fonction logarithme népérien]
La fonction \textbf{logarithme népérien}, notée $\ln$, est l'unique fonction dérivable sur $]0\,;\,+\infty[$ qui satisfait :
\[ \forall x>0,\quad \ln'(x)=\dfrac{1}{x} \qquad\text{et}\qquad \ln(1)=0. \]
Son ensemble de définition est $D_{\ln}=]0\,;\,+\infty[$.
\end{definition}

\begin{aretenir}[Conséquences immédiates de la définition]
\begin{itemize}
\item $\ln$ est \textbf{strictement croissante} sur $]0\,;\,+\infty[$ car sa dérivée est strictement positive.
\item $\ln$ est une \textbf{bijection} de $]0\,;\,+\infty[$ sur $\mathbb{R}$.
\item La fonction $\ln$ est l'application réciproque de la fonction exponentielle : $\forall x>0,\; e^{\ln x}=x$ et $\forall y\in\mathbb{R},\; \ln(e^y)=y$.
\item Le nombre $e$ (base des logarithmes népériens) est l'unique réel tel que $\ln e = 1$ ($e\approx 2,718$).
\end{itemize}
\end{aretenir}

\begin{attention}[Pièges classiques sur l'ensemble de définition]
\begin{itemize}
\item $\ln x$ n'existe \textbf{que si $x>0$}. On n'écrit jamais $\ln(-2)$ ni $\ln 0$.
\item Dans un exercice, la \textbf{première étape} est toujours de déterminer les conditions d'existence (C.E.) : arguments des $\ln$ strictement positifs.
\end{itemize}
\end{attention}

\begin{experience}[Découverte de la dérivée par la fonction réciproque]
On sait que $\exp$ est dérivable sur $\mathbb{R}$ avec $\exp'=\exp$ et que $\ln$ est sa bijection réciproque.
Soit $x>0$. Posons $y=\ln x$, alors $x=e^y$.
Dérivons l'égalité $x=e^y$ par rapport à $x$ : $1 = e^y \cdot y' = x \cdot (\ln)'(x)$.
Donc $(\ln)'(x) = \dfrac{1}{x}$.
\end{experience}

\begin{exemplebox}[Calculs directs de valeurs de $\ln$]
Sachant $\ln 2 \approx 0,693$ et $\ln 3 \approx 1,099$ (valeurs à connaître), on utilise les propriétés algébriques (vues section suivante) :
\begin{align*}
\ln 6 &= \ln(2\times 3) = \ln 2 + \ln 3 \approx 1,792 \\
\ln 12 &= \ln(4\times 3) = 2\ln 2 + \ln 3 \approx 2,485 \\
\ln \frac{3}{2} &= \ln 3 - \ln 2 \approx 0,406
\end{align*}
\end{exemplebox}

\begin{exoresolu}[QCM : les fondamentaux de ln]
Pour chaque question, une seule réponse est exacte. Entoure-la et justifie brièvement.
\begin{enumerate}
\item La fonction $\ln$ est définie sur :
\begin{enumerate}
\item[(a)] $\mathbb{R}$ \qquad (b) $[0\,;\,+\infty[$ \qquad (c) $]0\,;\,+\infty[$ \qquad (d) $\mathbb{R}^*$
\end{enumerate}
\item Pour tous réels $a>0$ et $b>0$, $\ln(ab)$ est égal à :
\begin{enumerate}
\item[(a)] $\ln a \times \ln b$ \qquad (b) $\ln a + \ln b$ \qquad (c) $\ln(a+b)$ \qquad (d) $a\ln b$
\end{enumerate}
\item $\displaystyle\lim_{x\to 0^+}\ln x$ est égale à :
\begin{enumerate}
\item[(a)] $0$ \qquad (b) $1$ \qquad (c) $-\infty$ \qquad (d) $+\infty$
\end{enumerate}
\item Une primitive de $f(x)=\dfrac{2x}{x^2+3}$ sur $\mathbb{R}$ est :
\begin{enumerate}
\item[(a)] $\ln(x^2+3)$ \qquad (b) $\dfrac{2}{x^2+3}$ \qquad (c) $2\ln x$ \qquad (d) $\dfrac{1}{x^2+3}$
\end{enumerate}
\end{enumerate}
\tcblower
\textbf{Solutions et justifications :}
\begin{enumerate}
\item \textbf{(c)} $]0\,;\,+\infty[$. Par définition, l'argument du logarithme doit être strictement positif.
\item \textbf{(b)} $\ln a + \ln b$. Propriété fondamentale : le logarithme transforme le produit en somme.
\item \textbf{(c)} $-\infty$. La courbe de $\ln$ a une asymptote verticale $x=0$ et tend vers $-\infty$ à droite de $0$.
\item \textbf{(a)} $\ln(x^2+3)$. On reconnaît la forme $\dfrac{u'}{u}$ avec $u(x)=x^2+3$ ($u'(x)=2x$). Comme $x^2+3>0$ sur $\mathbb{R}$, une primitive est $\ln(x^2+3)$.
\end{enumerate}
\end{exoresolu}

\begin{exoresolu}[Vrai ou faux, en justifiant]
Pour chaque affirmation, dire si elle est vraie ou fausse, en justifiant la réponse.
\begin{enumerate}
\item Pour tout réel $x>0$, $\ln(x^2)=2\ln x$.
\item Pour tout réel $x$, $\ln(x^2)=2\ln x$.
\item $\ln 1 = 0$ et $\ln e = 1$.
\item La fonction $\ln$ est strictement décroissante sur $]0\,;\,+\infty[$.
\item Si $\ln a = \ln b$, alors $a=b$.
\end{enumerate}
\tcblower
\begin{enumerate}
\item \textbf{Vrai}. Pour $x>0$, $x^2>0$ et $\ln(x^2)=2\ln x$ (propriété $\ln(a^n)=n\ln a$).
\item \textbf{Faux}. Si $x<0$, $\ln(x^2)$ est défini (car $x^2>0$) mais $2\ln x$ n'a pas de sens ($\ln x$ non défini). L'égalité correcte pour tout $x\neq 0$ est $\ln(x^2)=2\ln|x|$.
\item \textbf{Vrai}. Par définition $\ln 1=0$ ; $e$ est l'unique nombre tel que $\ln e=1$.
\item \textbf{Faux}. $\ln$ est \textbf{strictement croissante} sur $]0\,;\,+\infty[$ car sa dérivée $1/x > 0$.
\item \textbf{Vrai}. $\ln$ est injective (strictement monotone) sur son ensemble de définition.
\end{enumerate}
\end{exoresolu}

\courssub{Propriétés algébriques de $\ln$}

\begin{propriete}[Propriétés algébriques fondamentales]
Pour tous réels strictement positifs $a$ et $b$, et pour tout réel $n$ (entier, rationnel, puis réel par continuité) :
\begin{align*}
\ln(ab) &= \ln a + \ln b & \text{(produit $\to$ somme)} \\
\ln\!\left(\dfrac{a}{b}\right) &= \ln a - \ln b & \text{(quotient $\to$ différence)} \\
\ln(a^n) &= n\ln a & \text{(puissance $\to$ produit)} \\
\ln(\sqrt[n]{a}) &= \dfrac{1}{n}\ln a & \text{(racine $\to$ division)}
\end{align*}
\end{propriete}

\begin{methode}[Simplifier une expression contenant des $\ln$]
\begin{enumerate}
\item Vérifier que tous les arguments sont $>0$ (C.E.).
\item Appliquer les propriétés pour \textbf{éclater} les produits, quotients, puissances.
\item Regrouper les termes semblables (coefficients devant $\ln 2$, $\ln 3$, $\ln 5$, etc.).
\item Factoriser si nécessaire.
\end{enumerate}
\end{methode}

\begin{attention}[Erreurs fréquentes en calcul littéral]
\begin{itemize}
\item $\ln(a+b) \neq \ln a + \ln b$ (le logarithme n'est \textbf{pas} linéaire).
\item $\ln(a-b) \neq \ln a - \ln b$.
\item $\dfrac{\ln a}{\ln b} \neq \ln\!\left(\dfrac{a}{b}\right)$.
\item $\ln(x^2) = 2\ln|x|$ (valeur absolue indispensable si $x$ peut être négatif).
\end{itemize}
\end{attention}

\begin{exoresolu}[Calculs directs]
Simplifier les expressions suivantes, en écrivant le résultat sous la forme $k\ln 2 + k'\ln 3$ avec $k,k'\in\mathbb{Q}$ :
\[ A=\ln 12+\ln 3-2\ln 6+\ln\sqrt{2} \quad ; \quad B=\ln 48-\ln 27+\ln\dfrac{9}{4}. \]
\tcblower
\textbf{Expression $A$ :}
\begin{align*}
A &= \ln(2^2\times 3) + \ln 3 - 2\ln(2\times 3) + \ln(2^{1/2}) \\
  &= (2\ln 2 + \ln 3) + \ln 3 - 2(\ln 2 + \ln 3) + \tfrac{1}{2}\ln 2 \\
  &= 2\ln 2 + \ln 3 + \ln 3 - 2\ln 2 - 2\ln 3 + \tfrac{1}{2}\ln 2 \\
  &= \boxed{\tfrac{1}{2}\ln 2}
\end{align*}

\textbf{Expression $B$ :}
\begin{align*}
B &= \ln(2^4\times 3) - \ln(3^3) + \ln(3^2) - \ln(2^2) \\
  &= (4\ln 2 + \ln 3) - 3\ln 3 + 2\ln 3 - 2\ln 2 \\
  &= (4-2)\ln 2 + (1-3+2)\ln 3 \\
  &= \boxed{2\ln 2}
\end{align*}
\end{exoresolu}

\courssub{Limites de la fonction $\ln$}

\begin{propriete}[Limites usuelles de $\ln$]
\begin{align*}
\lim_{x\to 0^+} \ln x &= -\infty \\
\lim_{x\to +\infty} \ln x &= +\infty
\end{align*}
La courbe $\mathcal{C}_{\ln}$ a l'axe des ordonnées ($x=0$) comme \textbf{asymptote verticale}.
\end{propriete}

\begin{propriete}[Croissance lente du logarithme]
Le logarithme croît plus lentement que toute fonction puissance $x^\alpha$ ($\alpha>0$) :
\[ \forall \alpha>0,\quad \lim_{x\to +\infty} \dfrac{\ln x}{x^\alpha} = 0. \]
En particulier : $\lim\limits_{x\to +\infty} \dfrac{\ln x}{x} = 0$.
\end{propriete}

\begin{methode}[Calcul de limites de fonctions composées avec $\ln$]
\begin{enumerate}
\item Calculer la limite de l'expression \textbf{intérieure} (argument du $\ln$).
\item Si cette limite est $L \in ]0\,;\,+\infty[$, alors $\lim \ln(u(x)) = \ln L$.
\item Si la limite de l'argument est $0^+$, le $\ln$ tend vers $-\infty$.
\item Si la limite de l'argument est $+\infty$, le $\ln$ tend vers $+\infty$.
\item En cas d'indétermination ($\frac{\ln x}{x}$, $x\ln x$, etc.), utiliser les résultats de croissance ou factoriser.
\end{enumerate}
\end{methode}

\begin{exoresolu}[Calcul de limites]
Calculer les limites suivantes, en levant soigneusement les éventuelles indéterminations :
\begin{enumerate}
\item $\displaystyle\lim_{x\to +\infty}\dfrac{\ln x}{x}$ ;
\item $\displaystyle\lim_{x\to 0^+} x\ln x$ ;
\item $\displaystyle\lim_{x\to +\infty}(x-\ln x)$ ;
\item $\displaystyle\lim_{x\to +\infty}\dfrac{\ln x}{x^2+1}$ ;
\item $\displaystyle\lim_{x\to 1}\dfrac{\ln x}{x-1}$.
\end{enumerate}
\tcblower
\begin{enumerate}
\item Forme $\frac{+\infty}{+\infty}$. Croissance lente : $\lim\limits_{x\to +\infty}\frac{\ln x}{x}=0$.
\item Forme $0 \times (-\infty)$. Posons $X=\frac{1}{x}$ ($X\to +\infty$). $x\ln x = \frac{1}{X}(-\ln X) = -\frac{\ln X}{X} \to 0$. Donc la limite est $\boxed{0}$.
\item $x-\ln x = x\left(1-\frac{\ln x}{x}\right)$. Comme $\frac{\ln x}{x}\to 0$, le facteur $\to 1$. Produit $+\infty \times 1 = \boxed{+\infty}$.
\item Dénominateur $\sim x^2$. $\frac{\ln x}{x^2+1} \sim \frac{\ln x}{x^2} = \frac{1}{x}\frac{\ln x}{x} \to 0 \times 0 = \boxed{0}$.
\item Taux d'accroissement de $\ln$ en $1$ : $\frac{\ln x - \ln 1}{x-1} \to \ln'(1) = \frac{1}{1} = \boxed{1}$.
\end{enumerate}
\end{exoresolu}

\begin{savaistu}[Échelles logarithmiques au Cameroun et dans le monde]
La fonction $\ln$ (ou $\log_{10}$) permet de compresser de grandes gammes de valeurs :
\begin{itemize}
\item \textbf{pH} (chimie, qualité de l'eau du Wouri, sols agricoles) : $\text{pH} = -\log_{10}[\text{H}^+]$. Une variation de 1 unité = facteur 10 en concentration.
\item \textbf{Échelle de Richter} (sismicité, risque volcanique Mont Cameroun) : magnitude $M = \log_{10}(A/A_0)$.
\item \textbf{Décibels (dB)} (acoustique, télécoms MTN/Orange, bruit urbain à Douala/Yaoundé) : $L = 10\log_{10}(I/I_0)$.
\item \textbf{Échelle de magnitude stellaire} (astronomie, observation ciel camerounais).
\end{itemize}
Dans tous ces cas, une échelle logarithmique transforme des rapports multiplicatifs en différences additives, facilitant la lecture et la comparaison.
\end{savaistu}

\courssub{Dérivation, variations et courbe représentative}

\begin{propriete}[Dérivée de $\ln u$]
Soit $u$ une fonction dérivable et strictement positive sur un intervalle $I$.
La fonction $\ln u$ est dérivable sur $I$ et :
\[ (\ln u)' = \dfrac{u'}{u}. \]
\textbf{Démonstration exigible au Bac :} Dérivée de la composée $x \mapsto \ln(u(x))$ par la formule de la chaîne : $\ln'(u(x)) \cdot u'(x) = \frac{1}{u(x)} \cdot u'(x)$.
\end{propriete}

\begin{aretenir}[{Tableau de variations de $\ln x$ sur $]0\,;\,+\infty[$}]
\begin{center}
\begin{tabular}{|c|ccccc|}
\hline
$x$ & $0$ & & $1$ & & $+\infty$ \\
\hline
$\ln'(x)=\frac{1}{x}$ & & $+$ & & $+$ & \\
\hline
$\ln x$ & $-\infty$ & $\nearrow$ & $0$ & $\nearrow$ & $+\infty$ \\
\hline
\end{tabular}
\end{center}
\begin{itemize}
\item Strictement croissante.
\item Passe par $(1,0)$.
\item Tangente en $1$ : $y = x-1$ (car $\ln'(1)=1$).
\item Asymptote verticale $x=0$ (branche infinie en $0^+$).
\item Pas d'asymptote horizontale ni oblique en $+\infty$ (croissance lente mais non bornée).
\end{itemize}
\end{aretenir}

\begin{experience}[Démonstration de l'inégalité fondamentale $\ln x \le x-1$]
Étudions $f(x)=x-1-\ln x$ sur $]0\,;\,+\infty[$.
$f'(x)=1-\frac{1}{x}=\frac{x-1}{x}$.
$f'<0$ sur $]0,1[$, $f'>0$ sur $]1,+\infty[$.
$f$ admet un minimum global en $x=1$, $f(1)=0$.
Donc $\forall x>0,\; f(x)\ge 0 \iff \ln x \le x-1$.
Égalité seulement pour $x=1$.
\end{experience}

\begin{popfigure}
\begin{tikzpicture}[scale=0.8, domain=0.05:5, samples=200]
\draw[->] (-0.5,0) -- (5.5,0) node[right] {$x$};
\draw[->] (0,-3.5) -- (0,2.5) node[above] {$y$};
\draw[popcurve, thick] plot (\x, {ln(\x)});
\draw[poporange, thick, dashed] plot (\x, {\x-1}); % tangente
\draw[popgreen, thick, dashed] (1,-3.5) -- (1,2.5); % asymptote x=0 implied by axis
\draw[poppurple, thick] (0.05, {ln(0.05)}) -- (0.05, -3.5); % visual asymptote
\node[popcurve] at (4,1.6) {$y=\ln x$};
\node[poporange] at (4.5,3) {$y=x-1$};
\fill[popcurve] (1,0) circle (2pt) node[below right] {$(1,0)$};
\draw[dashed, thin] (1,0) -- (1,-3.5);
\draw[dashed, thin] (0,0) -- (1,0);
\end{tikzpicture}
\legende{Courbe de $y=\ln x$ et sa tangente en $(1,0)$ : $y=x-1$. L'axe des ordonnées est asymptote verticale.}
\end{popfigure}

\courssub{Primitives de la forme $\dfrac{u'}{u}$}

\begin{propriete}[Primitive de $\dfrac{u'}{u}$]
Soit $u$ une fonction dérivable et \textbf{ne s'annulant pas} sur un intervalle $I$.
Une primitive de $\dfrac{u'}{u}$ sur $I$ est $\ln|u|$ (ou $\ln u$ si $u>0$ sur $I$).
\[ \int \dfrac{u'(x)}{u(x)}\,dx = \ln|u(x)| + C. \]
\end{propriete}

\begin{attention}[Valeur absolue et ensemble de définition]
\begin{itemize}
\item Si l'intervalle $I$ n'est pas précisé, on écrit \textbf{systématiquement} $\ln|u(x)|$.
\item Si $u(x)>0$ sur $I$ (ex: $x^2+3$, $e^x$, $\cos x$ sur $]-\pi/2,\pi/2[$), on peut écrire $\ln(u(x))$.
\item \textbf{Piège :} $\int \frac{1}{x}\,dx = \ln|x| + C$ sur $\mathbb{R}^*$ (deux intervalles distincts : $]-\infty,0[$ et $]0,+\infty[$). Les constantes $C$ peuvent être différentes de part et d'autre de $0$.
\end{itemize}
\end{attention}

\begin{methode}[Rechercher une primitive de la forme $\frac{u'}{u}$]
\begin{enumerate}
\item Identifier $u(x)$ (souvent le dénominateur ou une fonction à l'intérieur d'un $\ln$).
\item Calculer $u'(x)$.
\item Vérifier que le numérateur est bien $u'(x)$ (ou un multiple constant $k\,u'(x)$).
\item Écrire la primitive : $\ln|u(x)|$ (ou $k\ln|u(x)|$).
\item Préciser l'intervalle de validité (où $u(x)\neq 0$).
\end{enumerate}
\end{methode}

\begin{exoresolu}[Recherche de primitives]
Déterminer une primitive de chacune des fonctions suivantes sur l'intervalle indiqué :
\begin{enumerate}
\item $f(x)=\dfrac{3x^2}{x^3+5}$ sur $\mathbb{R}$ ;
\item $g(x)=\tan x$ sur $\left]-\dfrac{\pi}{2}\,;\,\dfrac{\pi}{2}\right[$ ;
\item $h(x)=\dfrac{1}{x\ln x}$ sur $]1\,;\,+\infty[$ ;
\item $k(x)=\dfrac{1}{3x-2}$ sur $\left]\dfrac{2}{3}\,;\,+\infty\right[$.
\end{enumerate}
\tcblower
\begin{enumerate}
\item $u(x)=x^3+5$, $u'(x)=3x^2$. $u(x)$ s'annule pour $x=-\sqrt[3]{5}$. Sur $\mathbb{R}$, $u$ change de signe. Primitive : $\boxed{\ln|x^3+5|}$.
\item $g(x)=\tan x = \frac{\sin x}{\cos x} = -\frac{-\sin x}{\cos x} = -\frac{u'}{u}$ avec $u(x)=\cos x$. Sur $]-\pi/2,\pi/2[$, $\cos x > 0$. Primitive : $\boxed{-\ln(\cos x)}$ (ou $\ln|\sec x|$).
\item $u(x)=\ln x$, $u'(x)=\frac{1}{x}$. Sur $]1,+\infty[$, $\ln x > 0$. Primitive : $\boxed{\ln(\ln x)}$.
\item $u(x)=3x-2$, $u'(x)=3$. $\frac{1}{3x-2} = \frac{1}{3}\frac{u'}{u}$. Sur $]2/3,+\infty[$, $u>0$. Primitive : $\boxed{\frac{1}{3}\ln(3x-2)}$.
\end{enumerate}
\end{exoresolu}

\courssub{Équations, inéquations, systèmes et applications}

\begin{methode}[Résoudre une équation ou inéquation avec $\ln$]
\begin{enumerate}
\item \textbf{Poser les C.E.} : Tous les arguments des $\ln$ doivent être $>0$. Noter l'ensemble $\mathcal{D}$.
\item \textbf{Transformer} : Utiliser les propriétés algébriques pour n'avoir qu'un seul $\ln$ de chaque côté (équation) ou mettre du même côté (inéquation).
\item \textbf{Utiliser la monotonie} : $\ln$ est strictement croissante.
   \begin{itemize}
   \item $\ln A = \ln B \iff A=B$ (sur $\mathcal{D}$).
   \item $\ln A \le \ln B \iff 0 < A \le B$ (sur $\mathcal{D}$).
   \item $\ln A \ge k \iff A \ge e^k$ (si $k\in\mathbb{R}$).
   \end{itemize}
\item \textbf{Résoudre} l'équation/inéquation algébrique obtenue.
\item \textbf{Selectionner} les solutions appartenant à $\mathcal{D}$.
\end{enumerate}
\end{methode}

\begin{exoresolu}[Équations et inéquations avec conditions d'existence]
\begin{enumerate}
\item Résoudre dans $\mathbb{R}$ : $\ln(x+1)+\ln(x-1)=\ln 8$.
\item Résoudre dans $\mathbb{R}$ : $\ln(2x-3)\leq \ln(x+4)$.
\item Résoudre dans $\mathbb{R}$ : $\ln(3-x) \geq 0$.
\end{enumerate}
\tcblower
\begin{enumerate}
\item \textbf{C.E.} : $x+1>0$ et $x-1>0 \iff x>1$. $\mathcal{D}=]1,+\infty[$.
   $\ln((x+1)(x-1))=\ln 8 \iff \ln(x^2-1)=\ln 8 \iff x^2-1=8 \iff x^2=9 \iff x=3$ ou $x=-3$.
   Seul $x=3 \in \mathcal{D}$. \textbf{Solution : $S=\{3\}$}.
\item \textbf{C.E.} : $2x-3>0 \iff x>\frac{3}{2}$ et $x+4>0 \iff x>-4$. $\mathcal{D}=]\frac{3}{2},+\infty[$.
   $\ln$ croissante $\implies 2x-3 \le x+4 \iff x \le 7$.
   Intersection avec $\mathcal{D}$ : \textbf{$S=]\frac{3}{2}, 7]$}.
\item \textbf{C.E.} : $3-x>0 \iff x<3$. $\mathcal{D}=]-\infty,3[$.
   $\ln(3-x) \ge 0 = \ln 1 \iff 3-x \ge 1 \iff x \le 2$.
   Intersection avec $\mathcal{D}$ : \textbf{$S=]-\infty, 2]$}.
\end{enumerate}
\end{exoresolu}

\begin{methode}[Changement d'inconnue $X=\ln x$]
Pour une équation du type $a(\ln x)^2 + b\ln x + c = 0$ :
\begin{enumerate}
\item Poser $X=\ln x$ (donc $x=e^X$, $x>0$ automatiquement).
\item Résoudre l'équation du second degré en $X$.
\item Revenir à $x$ : $x=e^X$ (toujours $>0$).
\end{enumerate}
\end{methode}

\begin{exoresolu}[Changement d'inconnue et système]
\begin{enumerate}
\item Résoudre $(E)$ : $2(\ln x)^2-5\ln x+2=0$.
\item Résoudre dans $]0\,;\,+\infty[^2$ le système :
\[ \begin{cases} \ln x + \ln y = \ln 6 \\ \ln x - \ln y = \ln 2 \end{cases} \]
\end{enumerate}
\tcblower
\begin{enumerate}
\item Posons $X=\ln x$. $2X^2-5X+2=0$. $\Delta=25-16=9$. $X_1=\frac{5+3}{4}=2$, $X_2=\frac{5-3}{4}=\frac{1}{2}$.
   $x_1=e^2$, $x_2=e^{1/2}=\sqrt{e}$. \textbf{$S=\{e^2\,;\,\sqrt{e}\}$}.
\item Posons $X=\ln x$, $Y=\ln y$. Système linéaire :
   $\begin{cases} X+Y=\ln 6 \\ X-Y=\ln 2 \end{cases}$
   Addition : $2X = \ln 6 + \ln 2 = \ln 12 \implies X = \frac{1}{2}\ln 12 = \ln\sqrt{12} = \ln(2\sqrt{3})$.
   Soustraction : $2Y = \ln 6 - \ln 2 = \ln 3 \implies Y = \frac{1}{2}\ln 3 = \ln\sqrt{3}$.
   $x = e^X = 2\sqrt{3}$, $y = e^Y = \sqrt{3}$.
   Vérification : $xy=6$, $x/y=2$. \textbf{Solution unique : $(x,y)=(2\sqrt{3}\,;\,\sqrt{3})$}.
\end{enumerate}
\end{exoresolu}

\courssub{Je me prépare au Bac : Problèmes de synthèse}

\begin{exercice}[Étude complète de fonction et inégalité fondamentale — Manioc à Bafoussam]
On considère la fonction $f$ définie par $f(x)=x-1-\ln x$, et on note $(\mathcal{C}_f)$ sa courbe représentative dans un repère orthonormé $(O\,;\vec{\imath},\vec{\jmath})$ (unité : \SI{2}{cm}).

\textbf{Partie A : étude de la fonction}
\begin{enumerate}
\item Déterminer l'ensemble de définition $D_f$ de $f$.
\item Calculer les limites de $f$ aux bornes de $D_f$. (En $+\infty$, on factorisera par $x$.)
\item Calculer $f'(x)$ et étudier son signe sur $D_f$.
\item Dresser le tableau de variations complet de $f$.
\item Montrer que $f$ admet un minimum, et préciser sa valeur.
\item En déduire le signe de $f(x)$ sur $D_f$, puis démontrer l'inégalité : $\forall x>0,\quad \ln x \leq x-1$. Pour quelle valeur de $x$ a-t-on égalité ?
\item Tracer soigneusement $(\mathcal{C}_f)$ dans le repère, en précisant la tangente au point d'abscisse $1$ et le comportement asymptotique en $0^+$.
\end{enumerate}

\textbf{Partie B : application}
Une entreprise de transformation de manioc à Bafoussam modélise son bénéfice mensuel, en centaines de milliers de francs CFA, par $B(x)=x-1-\ln x$, où $x$ est la quantité de manioc transformée, en tonnes ($x>0$).
\begin{enumerate}
\item Le bénéfice peut-il être négatif ? Justifier à l'aide de la Partie A.
\item Pour quelle quantité de manioc le bénéfice est-il minimal ? Quel est ce bénéfice minimal, en francs CFA ?
\end{enumerate}
\end{exercice}

\begin{exercice}[Datation au carbone 14 d'un site archéologique — Grassfield]
Les archéologues qui étudient les sites anciens du Grassfield camerounais utilisent la datation au carbone 14. La quantité de carbone 14 restant dans un échantillon de bois mort suit la loi :
\[ N(t)=N_0\,e^{-\lambda t}, \]
où $N_0$ est la quantité initiale, $t$ le temps écoulé en années depuis la mort du végétal, et $\lambda=\dfrac{\ln 2}{5730}$ (la période radioactive du carbone 14 est de \num{5730} ans).

\begin{enumerate}
\item Justifier que $\dfrac{N(t)}{N_0}=e^{-\lambda t}$, puis montrer que :
\[ t=-\dfrac{1}{\lambda}\ln\!\left(\dfrac{N(t)}{N_0}\right). \]
\item Un échantillon de bois prélevé sur un site ne contient plus que $40\,\%$ de son carbone 14 initial.
\begin{enumerate}
\item Écrire l'équation vérifiée par l'âge $t$ de l'échantillon.
\item Montrer que $t=\dfrac{5730\times\ln(2{,}5)}{\ln 2}$.
\item Donner une valeur approchée de $t$ à l'année près (on donne $\ln 2\approx 0,693$ et $\ln(2{,}5)\approx 0,916$).
\end{enumerate}
\item Un second échantillon est daté de \num{11460} ans. Quelle proportion de carbone 14 initial contient-il encore ? Justifier sans calculatrice.
\item Un laboratoire ne sait pas mesurer des proportions inférieures à $1\,\%$. Montrer que la datation au carbone 14 devient impossible pour des échantillons trop anciens, et déterminer l'âge limite correspondant (on donnera le résultat en fonction de $\ln 100$ et $\ln 2$, puis une valeur approchée sachant que $\ln 100\approx 4,605$).
\end{enumerate}
\end{exercice}

\begin{exercice}[La fonction $x\mapsto \dfrac{\ln x}{x}$ et discussion d'une équation]
Soit $f$ la fonction définie par $f(x)=\dfrac{\ln x}{x}$, et $(\mathcal{C})$ sa courbe représentative dans un repère orthonormé.

\textbf{Partie A : étude de $f$}
\begin{enumerate}
\item Déterminer l'ensemble de définition de $f$.
\item Calculer les limites de $f$ en $0^+$ et en $+\infty$. Interpréter graphiquement chacun des résultats.
\item Montrer que pour tout $x>0$ : $f'(x)=\dfrac{1-\ln x}{x^2}$.
\item Étudier le signe de $f'(x)$ et dresser le tableau de variations de $f$.
\item Montrer que $f$ admet un maximum égal à $\dfrac{1}{e}$, atteint en $x=e$.
\item Déterminer le signe de $f(x)$ sur $]0\,;\,+\infty[$.
\item Tracer l'allure de $(\mathcal{C})$ en plaçant les éléments remarquables (asymptotes, maximum, intersection avec l'axe des abscisses).
\end{enumerate}

\textbf{Partie B : discussion graphique}
On cherche à déterminer, suivant les valeurs du réel $k$, le nombre de solutions dans $]0\,;\,+\infty[$ de l'équation :
\[ (E_k):\quad \ln x = kx. \]
\begin{enumerate}
\item Montrer que résoudre $(E_k)$ revient à résoudre $f(x)=k$.
\item En utilisant le tableau de variations de la Partie A, déterminer le nombre de solutions de $(E_k)$ dans chacun des cas :
\[ k>\dfrac{1}{e}\quad ;\quad k=\dfrac{1}{e}\quad ;\quad 0<k<\dfrac{1}{e}\quad ;\quad k\leq 0. \]
\item Application : un opérateur télécoms modélise le nombre d'abonnés (en millions) d'une offre par $N(t)$, et la rentabilité impose de résoudre $\ln x = \dfrac{x}{4}$. À l'aide de la question précédente, montrer que cette équation admet exactement deux solutions. Encadrer chacune d'elles entre deux entiers consécutifs (on donne $f(1)=0$, $f(2)\approx 0,347$, $f(8)\approx 0,260$, $f(9)\approx 0,244$ et $\dfrac{1}{4}=0,25$).
\end{enumerate}
\end{exercice}

\begin{savaistu}[Le nombre $e$ et l'histoire des logarithmes]
John Napier (Neper), Écossais (1550--1617), a inventé les logarithmes pour simplifier les calculs astronomiques (multiplications $\to$ additions). Ses "logarithmes napiériens" n'étaient pas encore en base $e$. C'est Leonhard Euler (1707--1783) qui a identifié la base $e = \lim_{n\to\infty}(1+\frac{1}{n})^n \approx 2,71828$ et imposé la notation $\ln$ (logarithmus naturalis). Au Cameroun, les ingénieurs de CAMTEL ou de la SONEL utilisent quotidiennement les échelles logarithmiques (dB pour l'atténuation du signal sur la fibre optique, népers pour les lignes de transmission).
\end{savaistu}

\newpage

\coursec{Corrigés des exercices}

\begin{corrige}[Exercice 1 — Étude complète de fonction et inégalité fondamentale (manioc à Bafoussam)]

\textbf{Partie A : étude de la fonction}

\textbf{1. Ensemble de définition.} La fonction $f(x)=x-1-\ln x$ fait intervenir $\ln x$, qui n'existe que pour $x>0$. Les termes $x$ et $-1$ sont définis pour tout réel. Donc :
\[ D_f = ]0\,;\,+\infty[. \]

\textbf{2. Limites aux bornes.}

\emph{En $0^+$ :} on a $\lim\limits_{x\to 0^+} x = 0$, $\lim\limits_{x\to 0^+}(-1)=-1$ et $\lim\limits_{x\to 0^+}\ln x = -\infty$. Par somme de limites (aucune indétermination) :
\[ \lim_{x\to 0^+} f(x) = 0 - 1 - (-\infty) = +\infty. \]
Interprétation graphique : la droite d'équation $x=0$ (l'axe des ordonnées) est \textbf{asymptote verticale} à $(\mathcal{C}_f)$.

\emph{En $+\infty$ :} on a une forme indéterminée du type « $+\infty - \infty$ ». On lève l'indétermination en factorisant par $x$ (consigne de l'énoncé) :
\[ f(x) = x\left(1 - \frac{1}{x} - \frac{\ln x}{x}\right). \]
Or $\lim\limits_{x\to +\infty}\dfrac{1}{x}=0$ et, par le théorème des \textbf{croissances comparées}, $\lim\limits_{x\to +\infty}\dfrac{\ln x}{x}=0$. La parenthèse tend donc vers $1$, et par produit de limites :
\[ \lim_{x\to +\infty} f(x) = +\infty \times 1 = +\infty. \]

\textbf{3. Dérivée et signe.} $f$ est dérivable sur $]0\,;\,+\infty[$ comme somme de fonctions dérivables ($x\mapsto x-1$ est polynomiale, $\ln$ est dérivable sur $]0,+\infty[$), et :
\[ f'(x) = 1 - \frac{1}{x} = \frac{x-1}{x}. \]
Comme $x>0$ sur $D_f$, le dénominateur est strictement positif, donc $f'(x)$ a le signe de $x-1$ :
\begin{itemize}
\item $f'(x) < 0$ sur $]0\,;\,1[$ ;
\item $f'(1) = 0$ ;
\item $f'(x) > 0$ sur $]1\,;\,+\infty[$.
\end{itemize}

\textbf{4. Tableau de variations complet.}
\begin{center}
\begin{tabular}{|c|ccccc|}
\hline
$x$ & $0$ & & $1$ & & $+\infty$ \\
\hline
$f'(x)$ & & $-$ & $0$ & $+$ & \\
\hline
$f(x)$ & $+\infty$ & \searrow & $0$ & \nearrow & $+\infty$ \\
\hline
\end{tabular}
\end{center}
(On place les limites $+\infty$ aux deux bornes et la valeur $f(1) = 1-1-\ln 1 = 0$.)

\textbf{5. Minimum.} $f$ est strictement décroissante sur $]0,1]$ puis strictement croissante sur $[1,+\infty[$ : elle admet donc en $x=1$ un \textbf{minimum global} sur $D_f$, qui vaut :
\[ f(1) = 1 - 1 - \ln 1 = 0. \]

\textbf{6. Signe de $f$ et inégalité fondamentale.} Puisque le minimum de $f$ sur $D_f$ vaut $0$, on a pour tout $x>0$ : $f(x) \geq f(1) = 0$, c'est-à-dire :
\[ \forall x>0,\quad x-1-\ln x \geq 0 \iff \boxed{\ln x \leq x-1}. \]
L'égalité a lieu si et seulement si $f(x)=0$, c'est-à-dire au point où le minimum est atteint : $\boxed{x=1}$.

\textbf{7. Tracé de $(\mathcal{C}_f)$.} Éléments à placer :
\begin{itemize}
\item asymptote verticale $x=0$ (la courbe « monte » vers $+\infty$ en longeant l'axe des ordonnées) ;
\item point $(1\,;\,0)$, avec tangente \textbf{horizontale} car $f'(1)=0$ : la tangente en ce point a pour équation $y = f(1) + f'(1)(x-1) = 0$ ;
\item la courbe décroît de $+\infty$ vers $0$ sur $]0,1]$, touche l'axe des abscisses en $(1,0)$ sans le traverser (c'est un minimum), puis croît vers $+\infty$ ;
\item avec l'unité de \SI{2}{cm}, le point $(1,0)$ est à \SI{2}{cm} de l'origine sur l'axe des abscisses.
\end{itemize}

\textbf{Partie B : application}

\textbf{1.} D'après la question A.6, $\forall x>0,\; B(x) = x-1-\ln x = f(x) \geq 0$. Le bénéfice \textbf{ne peut donc jamais être négatif} : quelle que soit la quantité de manioc transformée, l'entreprise ne perd pas d'argent dans ce modèle.

\textbf{2.} Le bénéfice est minimal lorsque $f$ atteint son minimum, c'est-à-dire pour $x = 1$ tonne de manioc. Ce bénéfice minimal vaut $B(1) = 0$ centaine de milliers de francs CFA, soit $\boxed{0\ \text{FCFA}}$ : avec une tonne, l'entreprise couvre exactement ses frais (seuil de rentabilité).

\textbf{Barème indicatif (sur 10) :} $D_f$ : 0,5 ; limites : 2 (dont 1 pour la factorisation en $+\infty$) ; dérivée et signe : 1,5 ; tableau : 1 ; minimum : 0,5 ; inégalité et cas d'égalité : 1,5 ; tracé : 1,5 ; Partie B : 1,5.

\textbf{Points de vigilance :}
\begin{itemize}
\item En $+\infty$, la limite est une forme indéterminée : il est \textbf{obligatoire} de factoriser par $x$ et de citer la croissance comparée $\lim\limits_{x\to+\infty}\frac{\ln x}{x}=0$. Une réponse directe du type « $+\infty-\infty=+\infty$ » est sanctionnée.
\item Pour le signe de $f'(x)$, préciser que $x>0$ avant de conclure que $f'$ a le signe de $x-1$.
\item Dans la question 6, le raisonnement doit passer par le \textbf{minimum} : « $f(x)\geq$ minimum $=0$ ». Le cas d'égalité est explicitement demandé, ne pas l'oublier.
\item En B.2, l'unité du modèle est la \emph{centaine de milliers de FCFA} : bien convertir le résultat.
\end{itemize}
\end{corrige}

\begin{corrige}[Exercice 2 — Datation au carbone 14 d'un site du Grassfield]

\textbf{1.} On part de $N(t) = N_0 e^{-\lambda t}$. Comme $N_0 > 0$, on peut diviser les deux membres par $N_0$ :
\[ \frac{N(t)}{N_0} = e^{-\lambda t}. \]
Le membre de gauche est strictement positif (quotient de deux quantités strictement positives), on peut donc appliquer $\ln$ aux deux membres, l'égalité étant préservée car $\ln$ est une bijection de $]0,+\infty[$ sur $\mathbb{R}$ :
\[ \ln\!\left(\frac{N(t)}{N_0}\right) = \ln\!\left(e^{-\lambda t}\right) = -\lambda t, \]
d'où, en divisant par $-\lambda \neq 0$ (car $\lambda > 0$) :
\[ \boxed{t = -\frac{1}{\lambda}\ln\!\left(\frac{N(t)}{N_0}\right)}. \]

\textbf{2.(a)} L'échantillon ne contient plus que $40\,\%$ de son carbone 14 initial : $N(t) = 0{,}4\,N_0$, soit $\dfrac{N(t)}{N_0} = 0{,}4$. L'âge $t$ vérifie donc :
\[ e^{-\lambda t} = 0{,}4. \]

\textbf{2.(b)} D'après la question 1 avec $\lambda = \dfrac{\ln 2}{5730}$ :
\begin{align*}
t &= -\frac{1}{\lambda}\ln(0{,}4) = -\frac{5730}{\ln 2}\ln(0{,}4) \\
  &= \frac{5730}{\ln 2}\ln\!\left(\frac{1}{0{,}4}\right) \quad \left(\text{car } -\ln a = \ln\frac{1}{a}\right) \\
  &= \frac{5730}{\ln 2}\ln(2{,}5) \quad \left(\text{car } \frac{1}{0{,}4} = \frac{10}{4} = 2{,}5\right).
\end{align*}
Donc $\boxed{t = \dfrac{5730\,\ln(2{,}5)}{\ln 2}}$.

\textbf{2.(c)} Valeur approchée, avec $\ln(2{,}5)\approx 0{,}916$ et $\ln 2 \approx 0{,}693$ :
\[ t \approx 5730 \times \frac{0{,}916}{0{,}693} \approx 5730 \times 1{,}3218 \approx 7573{,}9. \]
L'échantillon a environ $\boxed{7574\ \text{ans}}$ (à l'année près).

\textbf{3.} On remarque que $11460 = 2 \times 5730$ : l'échantillon a vécu exactement \textbf{deux périodes} radioactives. Calculons sans calculatrice :
\begin{align*}
\frac{N(t)}{N_0} &= e^{-\lambda \times 11460} = e^{-\frac{\ln 2}{5730}\times 2\times 5730} = e^{-2\ln 2} \\
&= e^{\ln(2^{-2})} = 2^{-2} = \frac{1}{4},
\end{align*}
où l'on a utilisé $-2\ln 2 = \ln(2^{-2})$ (propriété $a\ln b = \ln(b^a)$) puis $e^{\ln u} = u$. L'échantillon contient encore $\boxed{\dfrac{1}{4} = 25\,\%}$ de son carbone 14 initial. (Cohérent : une période divise la quantité par 2, deux périodes la divisent par 4.)

\textbf{4.} La datation devient impossible lorsque la proportion résiduelle passe sous le seuil de mesure de $1\,\%$. L'âge limite correspond à $\dfrac{N(t)}{N_0} = 0{,}01$ :
\[ t_{\text{lim}} = -\frac{1}{\lambda}\ln(0{,}01) = \frac{5730}{\ln 2}\ln\!\left(\frac{1}{0{,}01}\right) = \boxed{\frac{5730\,\ln 100}{\ln 2}}. \]
Valeur approchée, avec $\ln 100 \approx 4{,}605$ :
\[ t_{\text{lim}} \approx 5730 \times \frac{4{,}605}{0{,}693} \approx 5730 \times 6{,}645 \approx 38\,076\ \text{ans}. \]
Comme $t \mapsto e^{-\lambda t}$ est strictement décroissante, pour $t > t_{\text{lim}}$ la proportion résiduelle est strictement inférieure à $1\,\%$ et le laboratoire ne peut plus rien mesurer : la datation au carbone 14 est donc limitée aux échantillons de moins d'environ $\boxed{38\,000\ \text{ans}}$.

\textbf{Barème indicatif (sur 10) :} question 1 : 2 ; question 2 : 3 (1 par sous-question) ; question 3 : 2 ; question 4 : 3 (expression littérale 1,5 ; valeur approchée 1 ; conclusion 0,5).

\textbf{Points de vigilance :}
\begin{itemize}
\item Justifier la positivité de $\frac{N(t)}{N_0}$ avant d'appliquer $\ln$ : c'est la condition d'application du logarithme, elle est attendue dans la rédaction.
\item Le signe moins : $-\ln(0{,}4) = \ln(2{,}5)$ et non $\ln(0{,}4)$. Une erreur de signe donne un âge négatif, absurde : pensez à vérifier la cohérence physique (un âge est toujours positif).
\item Question 3 : l'énoncé exige une justification \textbf{sans calculatrice} ; il faut utiliser $e^{-2\ln 2} = \frac{1}{4}$ via la relation $e^{a\ln b} = b^a$, et non une valeur approchée.
\item Question 4 : bien conclure par une phrase d'interprétation (seuil de détection), pas seulement par le calcul.
\end{itemize}
\end{corrige}

\begin{corrige}[Exercice 3 — La fonction $x\mapsto \dfrac{\ln x}{x}$ et discussion d'une équation]

\textbf{Partie A : étude de $f$}

\textbf{1.} Le numérateur $\ln x$ exige $x>0$ et le dénominateur $x$ exige $x\neq 0$. La condition la plus restrictive est $x>0$, donc :
\[ D_f = ]0\,;\,+\infty[. \]

\textbf{2. Limites et interprétations graphiques.}

\emph{En $0^+$ :} $\lim\limits_{x\to 0^+}\ln x = -\infty$ et $\lim\limits_{x\to 0^+} x = 0$ avec $x>0$. Le quotient est de la forme « $\frac{-\infty}{0^+}$ » : il n'y a \textbf{pas} d'indétermination, et par quotient de limites (numérateur négatif, dénominateur positif tendant vers $0$) :
\[ \lim_{x\to 0^+} f(x) = -\infty. \]
Interprétation : la droite $x=0$ (axe des ordonnées) est \textbf{asymptote verticale} à $(\mathcal{C})$, la courbe plonge vers $-\infty$ au voisinage de $0$.

\emph{En $+\infty$ :} forme indéterminée du type « $\frac{\infty}{\infty}$ », levée par le théorème des croissances comparées :
\[ \lim_{x\to +\infty} \frac{\ln x}{x} = 0. \]
De plus $f(x)>0$ pour $x>1$, donc la limite est $0^+$. Interprétation : la droite $y=0$ (axe des abscisses) est \textbf{asymptote horizontale} à $(\mathcal{C})$ en $+\infty$, approchée par valeurs supérieures.

\textbf{3. Dérivée.} $f$ est un quotient $\frac{u}{v}$ avec $u(x)=\ln x$, $u'(x)=\frac{1}{x}$, $v(x)=x$, $v'(x)=1$, fonctions dérivables sur $]0,+\infty[$ avec $v$ qui ne s'annule pas :
\[ f'(x) = \frac{u'v - uv'}{v^2} = \frac{\frac{1}{x}\cdot x - \ln x \cdot 1}{x^2} = \boxed{\frac{1-\ln x}{x^2}}. \]

\textbf{4. Signe de $f'$ et tableau de variations.} Comme $x^2 > 0$ sur $D_f$, $f'(x)$ a le signe de $1-\ln x$. Par stricte croissance de $\ln$ :
\[ f'(x) > 0 \iff \ln x < 1 \iff x < e \; ; \qquad f'(x) = 0 \iff x = e \; ; \qquad f'(x) < 0 \iff x > e. \]
Valeur au sommet : $f(e) = \dfrac{\ln e}{e} = \dfrac{1}{e}$.
\begin{center}
\begin{tabular}{|c|ccccc|}
\hline
$x$ & $0$ & & $e$ & & $+\infty$ \\
\hline
$f'(x)$ & & $+$ & $0$ & $-$ & \\
\hline
$f(x)$ & $-\infty$ & \nearrow & $\frac{1}{e}$ & \searrow & $0^+$ \\
\hline
\end{tabular}
\end{center}

\textbf{5. Maximum.} $f$ est strictement croissante sur $]0,e]$ puis strictement décroissante sur $[e,+\infty[$ : elle admet donc en $x=e$ un \textbf{maximum global} sur $]0,+\infty[$, égal à :
\[ f(e) = \frac{1}{e} \approx 0{,}368. \]

\textbf{6. Signe de $f$.} Pour $x>0$, le dénominateur $x$ est strictement positif, donc $f(x)$ a le signe de $\ln x$ :
\begin{itemize}
\item $f(x) < 0$ sur $]0\,;\,1[$ ;
\item $f(1) = \frac{\ln 1}{1} = 0$ ;
\item $f(x) > 0$ sur $]1\,;\,+\infty[$.
\end{itemize}

\textbf{7. Allure de $(\mathcal{C})$.} Éléments remarquables à placer :
\begin{itemize}
\item asymptote verticale $x=0$ (branche vers $-\infty$) ;
\item unique intersection avec l'axe des abscisses au point $(1\,;\,0)$ ;
\item maximum au point $\left(e\,;\,\frac{1}{e}\right) \approx (2{,}72\,;\,0{,}37)$, avec tangente horizontale en ce point ;
\item asymptote horizontale $y=0$ en $+\infty$, la courbe restant au-dessus de l'axe des abscisses pour $x>1$.
\end{itemize}

\textbf{Partie B : discussion graphique}

\textbf{1.} Pour $x>0$, on peut diviser les deux membres par $x$ (non nul) :
\[ \ln x = kx \iff \frac{\ln x}{x} = k \iff f(x) = k. \]
Résoudre $(E_k)$ revient donc à chercher les abscisses des points d'intersection de $(\mathcal{C})$ avec la droite horizontale d'équation $y=k$.

\textbf{2.} On détermine le nombre de solutions à l'aide du tableau de variations, en appliquant sur chaque intervalle de stricte monotonie le \textbf{corollaire du théorème des valeurs intermédiaires} (une fonction continue et strictement monotone sur un intervalle y réalise une bijection sur son intervalle image) :
\begin{itemize}
\item $\boxed{k > \frac{1}{e}}$ : la droite est au-dessus du maximum de $f$ : \textbf{0 solution}.
\item $\boxed{k = \frac{1}{e}}$ : la droite touche $(\mathcal{C})$ en son sommet : \textbf{1 solution}, $x=e$.
\item $\boxed{0 < k < \frac{1}{e}}$ : sur $]0,e]$, $f$ est continue, strictement croissante, et $k \in \left]0,\frac{1}{e}\right[ \subset \left]-\infty,\frac{1}{e}\right[$ : une unique solution dans $]0,e[$, qui est même dans $]1,e[$ car $f(x)\leq 0$ pour $x\leq 1$. Sur $[e,+\infty[$, $f$ est continue, strictement décroissante, et $k \in \left]0,\frac{1}{e}\right[$ : une unique solution dans $]e,+\infty[$. Au total : \textbf{2 solutions}, l'une dans $]1,e[$, l'autre dans $]e,+\infty[$.
\item $\boxed{k \leq 0}$ : sur $[e,+\infty[$, $f(x)>0 \geq k$ : pas de solution. Sur $]0,e]$, $f$ est continue, strictement croissante de $-\infty$ vers $\frac{1}{e}$, et $k \leq 0 < \frac{1}{e}$ : exactement une solution, située dans $]0,1]$ (pour $k=0$, c'est $x=1$). Au total : \textbf{1 solution}.
\end{itemize}

\textbf{3. Application.} L'équation $\ln x = \dfrac{x}{4}$ équivaut, pour $x>0$, à $f(x) = \dfrac{1}{4}$, c'est-à-dire au cas $k = \dfrac{1}{4} = 0{,}25$.
Or $0 < 0{,}25 < \dfrac{1}{e} \approx 0{,}368$ : d'après la question précédente, l'équation admet \textbf{exactement deux solutions}.

Encadrements, en utilisant la stricte monotonie de $f$ sur chaque branche et les valeurs fournies par l'énoncé :
\begin{itemize}
\item Sur $]1,e[$, $f$ est strictement croissante et $f(1) = 0 < 0{,}25 < f(2) \approx 0{,}347$, donc la première solution vérifie $\boxed{1 < x_1 < 2}$.
\item Sur $]e,+\infty[$, $f$ est strictement \emph{décroissante} et $f(8) \approx 0{,}260 > 0{,}25 > f(9) \approx 0{,}244$ (les inégalités se renversent quand on applique une fonction décroissante), donc la seconde solution vérifie $\boxed{8 < x_2 < 9}$.
\end{itemize}

\textbf{Barème indicatif (sur 12) :} Partie A : ensemble de définition 0,5 ; limites et interprétations 2,5 ; dérivée 1,5 ; signe et tableau 1,5 ; maximum 0,5 ; signe de $f$ 0,5 ; tracé 1. Partie B : équivalence 1 ; discussion (4 cas) 2,5 ; application et encadrements 1,5.

\textbf{Points de vigilance :}
\begin{itemize}
\item En $0^+$, il n'y a \textbf{pas} d'indétermination : écrire directement « $\frac{-\infty}{0^+} = -\infty$ » en précisant le signe du dénominateur. En $+\infty$, au contraire, citer explicitement la croissance comparée.
\item Chaque limite doit être suivie de son \textbf{interprétation graphique} (asymptote), c'est demandé par l'énoncé.
\item Dans la discussion, la justification rigoureuse repose sur le \textbf{corollaire du théorème des valeurs intermédiaires} appliqué sur chaque intervalle de stricte monotonie (existence \emph{et} unicité), pas seulement sur une lecture du dessin.
\item Pour l'encadrement de $x_2$, la fonction est \emph{décroissante} : les inégalités se renversent ($f(8) > 0{,}25 > f(9)$ donne $8 < x_2 < 9$). C'est l'erreur la plus fréquente.
\item Ne pas confondre $\frac{1}{e} \approx 0{,}368$ avec $e \approx 2{,}718$ dans la comparaison avec $k=\frac14$.
\end{itemize}
\end{corrige}

\newpage

\coursec{Fiche bilan}

\coursec{Définition et premières propriétés de la fonction $\ln$}

\begin{definition}[Fonction logarithme népérien]
On appelle \cle{fonction logarithme népérien}, notée $\ln$, l'unique fonction dérivable sur $]0\,;\,+\infty[$ qui vérifie :
\begin{itemize}
  \item $\forall x \in ]0\,;\,+\infty[,\quad \ln'(x) = \dfrac{1}{x}$ \quad (c'est une primitive de $x \mapsto \frac{1}{x}$) ;
  \item $\ln(1) = 0$ \quad (elle s'annule en $1$).
\end{itemize}
Son ensemble de définition est $\mathcal{D}_{\ln} = ]0\,;\,+\infty[$. On ne peut calculer $\ln x$ que pour $x > 0$.
\end{definition}

\begin{propriete}[Valeurs remarquables et nombre $e$]
\begin{itemize}
  \item $\ln 1 = 0$ (par définition).
  \item Il existe un unique réel $e > 0$ tel que $\ln e = 1$. Ce nombre, appelé \cle{nombre d'Euler}, vaut approximativement $e \approx 2{,}718\,281\,828$.
  \item Pour tout $x > 0$, $\ln x > 0 \Leftrightarrow x > 1$ et $\ln x < 0 \Leftrightarrow 0 < x < 1$.
\end{itemize}
\end{propriete}

\begin{experience}[Découverte : $\ln$ comme aire sous l'hyperbole]
On considère la fonction $f(x) = \frac{1}{x}$ sur $]0\,;\,+\infty[$. Pour $x > 0$, on note $\mathcal{A}(x)$ l'aire du domaine borné par l'axe des abscisses, la courbe $y = \frac{1}{t}$, les droites $t=1$ et $t=x$ (aire comptée positivement si $x>1$, négativement si $x<1$).
\begin{enumerate}
  \item Montrer que $\mathcal{A}'(x) = \frac{1}{x}$ et $\mathcal{A}(1)=0$.
  \item En déduire que $\mathcal{A}(x) = \ln x$.
\end{enumerate}
\emph{Cette approche intégrale justifie l'existence et l'unicité de $\ln$ sans faire appel à la fonction exponentielle.}
\end{experience}

\begin{exemplebox}[Calculs directs et signe]
\begin{enumerate}
  \item $\ln e^2 = 2\ln e = 2$ (on verra la propriété de puissance plus loin).
  \item $\ln(\frac{1}{e}) = -\ln e = -1$.
  \item Déterminer le signe de $\ln(0{,}5)$ : comme $0 < 0{,}5 < 1$, $\ln(0{,}5) < 0$.
  \item $\ln(\sqrt{e}) = \frac{1}{2}$.
\end{enumerate}
\end{exemplebox}

\begin{attention}[Piège classique : ensemble de définition]
Ne jamais écrire $\ln(-2)$, $\ln(0)$ ou $\ln(x^2-1)$ sans vérifier que l'argument est \textbf{strictement positif}.
\begin{itemize}
  \item $\ln(x^2)$ n'est défini que si $x^2 > 0 \Leftrightarrow x \neq 0$.
  \item $\ln(x^2-1)$ n'est défini que si $x^2-1 > 0 \Leftrightarrow x < -1$ ou $x > 1$.
\end{itemize}
Toujours commencer par : \textbf{« Posons les conditions d'existence : ... »}.
\end{attention}

\coursec{Propriétés algébriques de $\ln$}

\begin{propriete}[Formules fondamentales -- Démonstration exigible au Bac]
Pour tous réels strictement positifs $a$ et $b$, et pour tout entier relatif $n$ (et même tout réel $n$ en Terminale) :
\begin{align*}
  \ln(ab) &= \ln a + \ln b \quad &\text{(Produit)} \\
  \ln\!\left(\frac{a}{b}\right) &= \ln a - \ln b \quad &\text{(Quotient)} \\
  \ln\!\left(\frac{1}{a}\right) &= -\ln a \quad &\text{(Inverse)} \\
  \ln(a^n) &= n\ln a \quad &\text{(Puissance)} \\
  \ln(\sqrt{a}) &= \frac{1}{2}\ln a \quad &\text{(Racine carrée, cas particulier $n=\frac{1}{2}$)}
\end{align*}
\end{propriete}

\begin{experience}[Démonstration guidée de $\ln(ab) = \ln a + \ln b$]
Soient $a>0$, $b>0$ fixés. On considère la fonction $g(x) = \ln(ax) - \ln x$ sur $]0\,;\,+\infty[$.
\begin{enumerate}
  \item Calculer $g'(x)$.
  \item En déduire que $g$ est constante.
  \item Calculer $g(1)$ et conclure.
\end{enumerate}
\emph{Astuce : cette méthode (dériver pour montrer l'égalité) fonctionne pour toutes les propriétés.}
\end{experience}

\begin{methode}[Simplifier une expression contenant des $\ln$]
\begin{enumerate}
  \item Vérifier les conditions d'existence (arguments $>0$).
  \item Transformer les sommes/différences en produits/quotients : $\ln A + \ln B = \ln(AB)$, $k\ln A = \ln(A^k)$.
  \item Simplifier l'argument du $\ln$ final si possible.
  \item Remettre sous forme développée si l'énoncé le demande.
\end{enumerate}
\end{methode}

\begin{exoresolu}[Simplification]
Simplifier $E = 2\ln 3 - \ln 2 + \ln 6 - \ln 9$.
\tcblower
Conditions : arguments $3,2,6,9 > 0$ (OK).
$E = \ln(3^2) - \ln 2 + \ln 6 - \ln 9 = \ln 9 - \ln 2 + \ln 6 - \ln 9$
$E = \ln 6 - \ln 2 = \ln\left(\frac{6}{2}\right) = \ln 3$.
\end{exoresolu}

\begin{exercice}[Entraînement]
Simplifier :
\begin{enumerate}
  \item $\ln 50 - 2\ln 5 + \ln 2$
  \item $\frac{1}{2}\ln 8 - \ln 2 + \ln\sqrt{2}$
  \item $\ln(x^2-1) - \ln(x-1)$ (préciser le domaine)
\end{enumerate}
\end{exercice}

\coursec{Limites de la fonction $\ln$}

\begin{propriete}[Limites usuelles]
\begin{align*}
  \lim_{x \to 0^+} \ln x &= -\infty \\
  \lim_{x \to +\infty} \ln x &= +\infty
\end{align*}
La droite d'équation $x=0$ (axe des ordonnées) est une \cle{asymptote verticale} à la courbe $\mathcal{C}_{\ln}$.
\end{propriete}

\begin{propriete}[Croissances comparées -- Résultats à connaître]
\begin{align*}
  \lim_{x \to +\infty} \frac{\ln x}{x} &= 0 \quad \text{($\ln x$ croît plus lentement que toute fonction puissance $x^\alpha, \alpha>0$)} \\
  \lim_{x \to 0^+} x\ln x &= 0 \\
  \lim_{x \to 0} \frac{\ln(1+x)}{x} &= 1 \quad \text{(Équivalence $\ln(1+x) \sim x$ en $0$)}
\end{align*}
\end{propriete}

\begin{methode}[Calculer une limite faisant intervenir $\ln$]
\begin{enumerate}
  \item Identifier les formes indéterminées ($0 \times \infty$, $\frac{\infty}{\infty}$, $\frac{0}{0}$).
  \item Utiliser les croissances comparées : $\ln x = o(x)$ en $+\infty$, $x\ln x \to 0$ en $0^+$.
  \item Pour les composées $\ln(u(x))$ : calculer $\lim u(x) = \ell$, puis $\lim \ln(u(x)) = \ln \ell$ si $\ell > 0$, $-\infty$ si $\ell = 0^+$, $+\infty$ si $\ell = +\infty$.
\end{enumerate}
\end{methode}

\begin{exoresolu}[Limites composées]
Étudier les limites de $f(x) = \ln(x^2 - 3x + 2)$ aux bornes de son domaine.
\tcblower
Domaine : $x^2-3x+2 > 0 \Leftrightarrow (x-1)(x-2) > 0 \Leftrightarrow x \in ]-\infty;1[ \cup ]2;+\infty[$.
\begin{itemize}
  \item $x \to 1^-$ : $u(x) \to 0^+ \Rightarrow \ln(u(x)) \to -\infty$.
  \item $x \to 2^+$ : $u(x) \to 0^+ \Rightarrow \ln(u(x)) \to -\infty$.
  \item $x \to -\infty$ : $u(x) \sim x^2 \to +\infty \Rightarrow \ln(u(x)) \to +\infty$.
  \item $x \to +\infty$ : $u(x) \sim x^2 \to +\infty \Rightarrow \ln(u(x)) \to +\infty$.
\end{itemize}
\end{exoresolu}

\begin{savaistu}[Le pH et l'échelle logarithmique]
En chimie, le \cle{pH} (potentiel hydrogène) mesure l'acidité d'une solution : $\mathrm{pH} = -\log_{10}[\mathrm{H}^+] = -\frac{\ln[\mathrm{H}^+]}{\ln 10}$.
Une variation de 1 unité de pH correspond à une multiplication par 10 de la concentration en ions $\mathrm{H}^+$.
Au Cameroun, le contrôle de la qualité de l'eau (CAMWATER, forages ruraux) utilise cette échelle : une eau potable a un pH entre 6,5 et 8,5.
\end{savaistu}

\coursec{Dérivation, variations et courbe représentative}

\begin{propriete}[Dérivée de $\ln u$]
Soit $u$ une fonction dérivable et strictement positive sur un intervalle $I$.
La fonction $\ln \circ u$ est dérivable sur $I$ et :
\[ (\ln u)'(x) = \frac{u'(x)}{u(x)} \]
\textbf{Preuve :} Dérivée de la composée : $(\ln \circ u)' = \ln'(u) \times u' = \frac{1}{u} \times u' = \frac{u'}{u}$.
\end{propriete}

\begin{propriete}[Variations de $\ln$]
La fonction $\ln$ est \cle{strictement croissante} sur $]0\,;\,+\infty[$ car $\ln'(x) = \frac{1}{x} > 0$.
Elle est \cle{concave} car $\ln''(x) = -\frac{1}{x^2} < 0$.
\end{propriete}

\begin{popfigure}
\begin{tikzpicture}[scale=0.8, domain=0.05:5]
\draw[->] (-0.5,0) -- (5.5,0) node[right] {$x$};
\draw[->] (0,-3.5) -- (0,2.5) node[above] {$y$};
\draw[popcurve, thick, samples=200] plot (\x, {ln(\x)});
\draw[dashed, popmuted] (1,0) node[below left] {$1$} -- (1,-3);
\draw[dashed, popmuted] (2.718,0) node[below] {$e$} -- (2.718,1);
\draw[dashed, popmuted] (0,1) node[left] {$1$} -- (2.718,1);
\fill[popblue] (1,0) circle (2pt) node[below right] {$(1;0)$};
\fill[popblue] (2.718,1) circle (2pt) node[above left] {$(e;1)$};
\draw[poporange, thick, dashed] (0,0) -- (0,-3) node[left, pos=0.5] {Asymptote $x=0$};
\end{tikzpicture}
\legende{Courbe représentative $y = \ln x$ : passage par $(1;0)$ et $(e;1)$, asymptote verticale $x=0$, concavité vers le bas.}
\end{popfigure}

\begin{methode}[Étudier une fonction $f(x) = \ln(u(x))$]
\begin{enumerate}
  \item \textbf{Domaine :} Résoudre $u(x) > 0$.
  \item \textbf{Dérivée :} $f'(x) = \frac{u'(x)}{u(x)}$. Le signe de $f'$ est le signe de $u'$ (car $u>0$ sur le domaine).
  \item \textbf{Tableau de variations :} D'après le signe de $u'$.
  \item \textbf{Limites aux bornes :} Utiliser $\lim_{X\to 0^+}\ln X = -\infty$ et $\lim_{X\to +\infty}\ln X = +\infty$.
  \item \textbf{Courbe :} Tracer les asymptotes verticales (bornes où $u\to 0^+$), points remarquables ($u=1 \Rightarrow f=0$, $u=e \Rightarrow f=1$), concavité ($f'' = \frac{u''u - (u')^2}{u^2}$).
\end{enumerate}
\end{methode}

\begin{exoresolu}[Étude complète : $f(x) = \ln(x^2 - 2x + 2)$]
\tcblower
1. \textbf{Domaine :} $x^2-2x+2 = (x-1)^2+1 \ge 1 > 0$ pour tout $x \in \mathbb{R}$. $\mathcal{D}_f = \mathbb{R}$.
2. \textbf{Dérivée :} $u(x) = x^2-2x+2$, $u'(x) = 2x-2 = 2(x-1)$.
$f'(x) = \frac{2(x-1)}{x^2-2x+2}$. Dénominateur $>0$. Signe de $f'$ = signe de $x-1$.
3. \textbf{Variations :} Décroissante sur $]-\infty;1]$, croissante sur $[1;+\infty[$. Minimum en $1$ : $f(1) = \ln 1 = 0$.
4. \textbf{Limites :} $x \to \pm\infty \Rightarrow u(x) \sim x^2 \to +\infty \Rightarrow f(x) \to +\infty$.
5. \textbf{Courbe :} Symétrique par rapport à $x=1$ (car $u$ symétrique). Point bas $(1;0)$. Pas d'asymptote verticale. En $\pm\infty$, $f(x) \sim \ln(x^2) = 2\ln|x|$ : la courbe monte vers $+\infty$ très lentement (branche parabolique de direction $(Ox)$).
\end{exoresolu}

\begin{attention}[Signe de la dérivée]
Pour $f = \ln u$, $f' = \frac{u'}{u}$. Comme $u > 0$ sur le domaine, \textbf{le signe de $f'$ est exactement le signe de $u'$}. Inutile de faire un tableau de signes avec $u$ au dénominateur : il est positif !
\end{attention}

\coursec{Primitives de la forme $\frac{u'}{u}$}

\begin{propriete}[Primitive de $\frac{u'}{u}$]
Soit $u$ une fonction dérivable sur un intervalle $I$, ne s'annulant pas sur $I$.
Une primitive de $\frac{u'}{u}$ sur $I$ est $\ln|u|$.
\[ \int \frac{u'(x)}{u(x)}\,\mathrm{d}x = \ln|u(x)| + C \]
\textbf{Cas particulier :} $\int \frac{\mathrm{d}x}{x} = \ln|x| + C$ (sur $]-\infty;0[$ et sur $]0;+\infty[$).
\end{propriete}

\begin{attention}[Valeur absolue obligatoire !]
La primitive de $\frac{1}{x}$ sur $\mathbb{R}^*$ est $\ln|x| + C$, \textbf{pas} $\ln x + C$.
Exemple : $\int_{-2}^{-1} \frac{\mathrm{d}x}{x} = [\ln|x|]_{-2}^{-1} = \ln 1 - \ln 2 = -\ln 2$.
Si on utilisait $\ln x$, on aurait $\ln(-1)$ qui n'existe pas.
\end{attention}

\begin{methode}[Repérer et calculer $\int \frac{u'}{u}$]
\begin{enumerate}
  \item Identifier le dénominateur $u(x)$.
  \item Vérifier que le numérateur est bien $u'(x)$ (éventuellement à une constante près).
  \item Écrire $\ln|u(x)| + C$ (ou $k\ln|u(x)| + C$ si numérateur $= k u'$).
\end{enumerate}
\end{methode}

\begin{exoresolu}[Primitives]
Calculer les primitives suivantes :
\begin{enumerate}
  \item $\int \frac{2x}{x^2+1}\,\mathrm{d}x$
  \item $\int \tan x \,\mathrm{d}x$ sur $]-\frac{\pi}{2};\frac{\pi}{2}[$
  \item $\int \frac{e^x}{e^x+3}\,\mathrm{d}x$
\end{enumerate}
\tcblower
1. $u = x^2+1$, $u' = 2x$. Primitive $= \ln(x^2+1) + C$ (valeur absolue inutile car $x^2+1>0$).
2. $\tan x = \frac{\sin x}{\cos x} = \frac{-\cos' x}{\cos x}$. $u = \cos x$, $u' = -\sin x$. Primitive $= -\ln|\cos x| + C = \ln|\sec x| + C$.
3. $u = e^x+3$, $u' = e^x$. Primitive $= \ln(e^x+3) + C$.
\end{exoresolu}

\coursec{Équations, inéquations, systèmes et applications}

\begin{propriete}[Résolution d'équations et inéquations]
$\ln$ est strictement croissante sur $]0;+\infty[$, donc pour $A>0, B>0$ :
\begin{align*}
  \ln A = \ln B &\Leftrightarrow A = B \\
  \ln A < \ln B &\Leftrightarrow A < B \\
  \ln A \leq \ln B &\Leftrightarrow A \leq B
\end{align*}
\end{propriete}

\begin{methode}[Résoudre une équation $\ln(\dots) = \dots$]
\begin{enumerate}
  \item \textbf{Conditions d'existence (C.E.) :} Tous les arguments des $\ln$ doivent être $>0$. Noter l'ensemble $\mathcal{D}$.
  \item \textbf{Transformer :} Utiliser les propriétés algébriques pour n'avoir qu'un seul $\ln$ de chaque côté (ou isoler $\ln X$).
  \item \textbf{Résoudre :} Passer à l'équation sans $\ln$ (ex: $\ln A = \ln B \to A=B$ ; $\ln X = k \to X = e^k$).
  \item \textbf{Vérifier :} Ne garder que les solutions appartenant à $\mathcal{D}$.
\end{enumerate}
\end{methode}

\begin{exoresolu}[Équation]
Résoudre $\ln(x+3) + \ln(x-1) = \ln 5$.
\tcblower
1. \textbf{C.E.} : $x+3 > 0$ et $x-1 > 0 \Rightarrow x > 1$. $\mathcal{D} = ]1;+\infty[$.
2. $\ln[(x+3)(x-1)] = \ln 5$.
3. $(x+3)(x-1) = 5 \Leftrightarrow x^2+2x-3 = 5 \Leftrightarrow x^2+2x-8=0 \Leftrightarrow (x+4)(x-2)=0$.
Solutions candidates : $x = -4$ ou $x = 2$.
4. \textbf{Vérification :} $-4 \notin \mathcal{D}$, $2 \in \mathcal{D}$.
\textbf{Solution :} $S = \{2\}$.
\end{exoresolu}

\begin{methode}[Résoudre un système d'inconnues $\ln x, \ln y$]
Poser $X = \ln x$, $Y = \ln y$ (avec $x>0, y>0$). On obtient un système linéaire en $X, Y$.
Résoudre pour $X, Y$, puis retrouver $x = e^X, y = e^Y$.
\end{methode}

\begin{exoresolu}[Système]
Résoudre $\begin{cases} \ln x + 2\ln y = 3 \\ 2\ln x - \ln y = 0 \end{cases}$.
\tcblower
C.E. : $x>0, y>0$.
Posons $X=\ln x, Y=\ln y$.
$\begin{cases} X + 2Y = 3 \\ 2X - Y = 0 \end{cases} \Rightarrow Y = 2X \Rightarrow X + 4X = 3 \Rightarrow X = \frac{3}{5}, Y = \frac{6}{5}$.
$x = e^{3/5}, \quad y = e^{6/5}$.
Vérification : $x>0, y>0$ OK.
$S = \{(e^{0,6}\,;\, e^{1,2})\}$.
\end{exoresolu}

\begin{savaistu}[Croissance exponentielle et datation au Carbone 14]
La désintégration radioactive suit une loi $N(t) = N_0 e^{-\lambda t}$.
Le \cle{période} (demi-vie) $T$ vérifie $e^{-\lambda T} = \frac{1}{2} \Rightarrow \lambda = \frac{\ln 2}{T}$.
Pour le Carbone 14 ($T \approx 5730$ ans), on date un échantillon en mesurant $N/N_0$ :
$t = -\frac{1}{\lambda}\ln(N/N_0) = -T \frac{\ln(N/N_0)}{\ln 2}$.
Au Cameroun, cette méthode aide à dater les sites archéologiques (ex: gisements de Shum Laka dans le Nord-Ouest).
\end{savaistu}

\begin{exemplebox}[Application : Croissance démographique]
La population d'une ville suit $P(t) = P_0 e^{kt}$ ($t$ en années).
En 2020 ($t=0$), $P_0 = 500\,000$. En 2030, $P = 750\,000$.
Trouver $k$ et la date où la population atteindra $1\,000\,000$.
\tcblower
$750\,000 = 500\,000 e^{10k} \Rightarrow 1,5 = e^{10k} \Rightarrow \ln 1,5 = 10k \Rightarrow k = \frac{\ln 1,5}{10} \approx 0,0405$.
$1\,000\,000 = 500\,000 e^{kt} \Rightarrow 2 = e^{kt} \Rightarrow \ln 2 = kt \Rightarrow t = \frac{\ln 2}{k} = 10 \frac{\ln 2}{\ln 1,5} \approx 17,1$ ans.
Vers 2037.
\end{exemplebox}

\begin{exercice}[Synthèse Bac]
Soit $f(x) = \ln(x^2 - 4x + 5)$.
\begin{enumerate}
  \item Déterminer $\mathcal{D}_f$ et étudier la parité de $f$.
  \item Calculer $f'(x)$ et dresser le tableau de variations de $f$.
  \item Calculer les limites aux bornes et tracer la courbe $\mathcal{C}_f$.
  \item Résoudre $f(x) = \ln 2$.
  \item Calculer $\int_1^3 f'(x)\,\mathrm{d}x$ et interpréter graphiquement.
\end{enumerate}
\end{exercice}

\begin{corrige}[Exercice Synthèse]
1. $x^2-4x+5 = (x-2)^2+1 \ge 1 > 0$, $\mathcal{D}_f = \mathbb{R}$. $f(-x) = \ln(x^2+4x+5) \neq f(x)$ : ni paire ni impaire. Symétrie axiale $x=2$ car $u(2+x)=u(2-x)$.
2. $u(x) = (x-2)^2+1$, $u'(x) = 2(x-2)$. $f'(x) = \frac{2(x-2)}{(x-2)^2+1}$. Signe de $f'$ = signe de $x-2$.
Décroissante sur $]-\infty;2]$, croissante sur $[2;+\infty[$. Minimum $f(2) = \ln 1 = 0$.
3. $x\to\pm\infty \Rightarrow u\sim x^2 \to +\infty \Rightarrow f\to +\infty$. Branche infinie : $f(x) - 2\ln|x| \to \ln 1 = 0$. Courbe au-dessus de l'axe $x$ (minimum 0), symétrique axe $x=2$, point $(2;0)$.
4. $\ln(x^2-4x+5) = \ln 2 \Rightarrow x^2-4x+5 = 2 \Rightarrow x^2-4x+3=0 \Rightarrow (x-1)(x-3)=0$. $S=\{1;3\}$.
5. $\int_1^3 f'(x)\,\mathrm{d}x = f(3)-f(1) = \ln 2 - \ln 2 = 0$. Aire algébrique nulle (symétrie par rapport à $x=2$).
\end{corrige}

\courssub{Carte mentale de la leçon}

\begin{popfigure}
\begin{center}\tcbox[colback=popdark,colframe=popdark,coltext=white,arc=9pt,boxsep=4pt,left=6pt,right=6pt]{\bfseries\small Fonction $\ln$}\end{center}
\vspace{-2pt}
\begin{tcbraster}[raster columns=3,raster equal height=rows,raster column skip=6pt,raster row skip=6pt,enhanced,blankest]
\begin{tcolorbox}[enhanced,colback=white,colframe=popblue,boxrule=0.9pt,arc=3pt,title={Définition},fonttitle=\bfseries\scriptsize,coltitle=white,colbacktitle=popblue,left=4pt,right=4pt,top=2pt,bottom=2pt,boxsep=1pt,toptitle=1.5pt,bottomtitle=1.5pt,halign=left]
\begin{itemize}[nosep,leftmargin=*,itemsep=1pt,font=\scriptsize]
\item primitive de $\frac{1}{x}$
\item $\ln 1 = 0$
\item $\mathcal{D}=]0,+\infty[$
\end{itemize}
\end{tcolorbox}
\begin{tcolorbox}[enhanced,colback=white,colframe=poporange,boxrule=0.9pt,arc=3pt,title={Propriétés algébriques},fonttitle=\bfseries\scriptsize,coltitle=white,colbacktitle=poporange,left=4pt,right=4pt,top=2pt,bottom=2pt,boxsep=1pt,toptitle=1.5pt,bottomtitle=1.5pt,halign=left]
\begin{itemize}[nosep,leftmargin=*,itemsep=1pt,font=\scriptsize]
\item $\ln(ab)=\ln a+\ln b$
\item $\ln(a/b)=\ln a-\ln b$
\item $\ln(a^n)=n\ln a$
\item $\ln\sqrt{a}=\frac12\ln a$
\end{itemize}
\end{tcolorbox}
\begin{tcolorbox}[enhanced,colback=white,colframe=poppurple,boxrule=0.9pt,arc=3pt,title={Limites},fonttitle=\bfseries\scriptsize,coltitle=white,colbacktitle=poppurple,left=4pt,right=4pt,top=2pt,bottom=2pt,boxsep=1pt,toptitle=1.5pt,bottomtitle=1.5pt,halign=left]
\begin{itemize}[nosep,leftmargin=*,itemsep=1pt,font=\scriptsize]
\item $\ln x \to -\infty$ en $0^+$
\item $\ln x \to +\infty$ en $+\infty$
\item $\frac{\ln x}{x}\to 0$ en $+\infty$
\item $x\ln x\to 0$ en $0^+$
\item $\frac{\ln(1+x)}{x}\to 1$ en $0$
\end{itemize}
\end{tcolorbox}
\begin{tcolorbox}[enhanced,colback=white,colframe=popgreen,boxrule=0.9pt,arc=3pt,title={Dérivation},fonttitle=\bfseries\scriptsize,coltitle=white,colbacktitle=popgreen,left=4pt,right=4pt,top=2pt,bottom=2pt,boxsep=1pt,toptitle=1.5pt,bottomtitle=1.5pt,halign=left]
\begin{itemize}[nosep,leftmargin=*,itemsep=1pt,font=\scriptsize]
\item $(\ln u)'=\frac{u'}{u}$
\item $\ln$ strictement croissante
\item Concave
\end{itemize}
\end{tcolorbox}
\begin{tcolorbox}[enhanced,colback=white,colframe=poppink,boxrule=0.9pt,arc=3pt,title={Primitives},fonttitle=\bfseries\scriptsize,coltitle=white,colbacktitle=poppink,left=4pt,right=4pt,top=2pt,bottom=2pt,boxsep=1pt,toptitle=1.5pt,bottomtitle=1.5pt,halign=left]
\begin{itemize}[nosep,leftmargin=*,itemsep=1pt,font=\scriptsize]
\item $\frac{u'}{u}\to \ln|u|+C$
\item $\frac{1}{x}\to \ln|x|+C$
\end{itemize}
\end{tcolorbox}
\begin{tcolorbox}[enhanced,colback=white,colframe=popgold,boxrule=0.9pt,arc=3pt,title={Équations/Inéquations},fonttitle=\bfseries\scriptsize,coltitle=white,colbacktitle=popgold,left=4pt,right=4pt,top=2pt,bottom=2pt,boxsep=1pt,toptitle=1.5pt,bottomtitle=1.5pt,halign=left]
\begin{itemize}[nosep,leftmargin=*,itemsep=1pt,font=\scriptsize]
\item C.E. : arguments $>0$
\item $\ln A=\ln B \Leftrightarrow A=B$
\item Systèmes : $X=\ln x$
\end{itemize}
\end{tcolorbox}
\begin{tcolorbox}[enhanced,colback=white,colframe=popmuted,boxrule=0.9pt,arc=3pt,title={Applications},fonttitle=\bfseries\scriptsize,coltitle=white,colbacktitle=popmuted,left=4pt,right=4pt,top=2pt,bottom=2pt,boxsep=1pt,toptitle=1.5pt,bottomtitle=1.5pt,halign=left]
\begin{itemize}[nosep,leftmargin=*,itemsep=1pt,font=\scriptsize]
\item Croissance/Désintégration
\item pH, échelle Richter
\item Datation $^{14}$C
\end{itemize}
\end{tcolorbox}
\end{tcbraster}

\legende{Carte mentale de la leçon : la fonction logarithme népérien}
\end{popfigure}

\courssub{L'essentiel à retenir}

\begin{aretenir}[Fonction logarithme népérien -- Formulaire Bac]
\textbf{Définitions clés}
\begin{itemize}
  \item La fonction \cle{logarithme népérien}, notée $\ln$, est l'unique \cle{primitive} de $x \mapsto \dfrac{1}{x}$ sur $]0\,;\,+\infty[$ qui \cle{s'annule en 1} : $\ln 1 = 0$.
  \item Ensemble de définition : $\mathcal{D}_{\ln} = ]0\,;\,+\infty[$. On ne peut donc calculer $\ln x$ que si $x > 0$.
  \item Nombre d'Euler : $e$ est l'unique réel tel que $\ln e = 1$ ($e \approx 2{,}718$).
\end{itemize}

\textbf{Formules et théorèmes à connaître par cœur}
\begin{center}
\begin{tabularx}{\linewidth}{|l|X|}
\hline
\rowcolor{popblueL} \textbf{Domaine} & \textbf{Formule} \\
\hline
Produit & $\ln(ab) = \ln a + \ln b$ \quad ($a>0$, $b>0$) \\
\hline
Quotient & $\ln\!\left(\dfrac{a}{b}\right) = \ln a - \ln b$ ; \quad $\ln\!\left(\dfrac{1}{a}\right) = -\ln a$ \\
\hline
Puissance & $\ln(a^n) = n\ln a$ \quad ($n\in\mathbb{R}$, $a>0$) ; \quad $\ln(\sqrt{a}) = \dfrac{1}{2}\ln a$ \\
\hline
Dérivée & $(\ln x)' = \dfrac{1}{x}$ ; \quad $(\ln u)' = \dfrac{u'}{u}$ (avec $u>0$) \\
\hline
Primitive & $\displaystyle\int \frac{u'}{u}\,\mathrm{d}x = \ln|u| + C$ ; \quad $\displaystyle\int \frac{\mathrm{d}x}{x} = \ln|x| + C$ \\
\hline
Limites & $\displaystyle\lim_{x\to 0^+}\ln x = -\infty$ ; \quad $\displaystyle\lim_{x\to +\infty}\ln x = +\infty$ \\
\hline
Croissances comparées & $\displaystyle\lim_{x\to +\infty}\frac{\ln x}{x} = 0$ ; \quad $\displaystyle\lim_{x\to 0^+} x\ln x = 0$ ; \quad $\displaystyle\lim_{x\to 0}\frac{\ln(1+x)}{x} = 1$ \\
\hline
Équivalence & $\ln a = \ln b \Leftrightarrow a = b$ ; \quad $\ln a < \ln b \Leftrightarrow a < b$ ($\ln$ est \cle{strictement croissante}) \\
\hline
Signe & $\ln x < 0$ si $0 < x < 1$ ; \quad $\ln x > 0$ si $x > 1$ \\
\hline
\end{tabularx}
\end{center}

\textbf{Méthodes express}
\begin{itemize}
  \item \cle{Simplifier} une expression : regrouper avec $\ln a + \ln b = \ln(ab)$ et $n\ln a = \ln(a^n)$.
  \item \cle{Équation} avec $\ln$ : (1) poser les conditions d'existence, (2) se ramener à $\ln A = \ln B$ ou à $X = \ln x$, (3) résoudre, (4) vérifier que les solutions respectent les conditions.
  \item \cle{Inéquation} : isoler $\ln x$, puis utiliser la stricte croissance : $\ln x \geq \ln a \Leftrightarrow x \geq a$.
  \item \cle{Étude de fonction} $f = \ln u$ : domaine ($u>0$), dérivée $f' = \dfrac{u'}{u}$, signe de $f'$ = signe de $u'$, limites aux bornes.
  \item \cle{Primitive} de $\dfrac{u'}{u}$ : repérer que le numérateur est la dérivée du dénominateur, écrire $\ln|u| + C$.
\end{itemize}
\end{aretenir}

\begin{aretenir}[Checklist avant le Bac]
Avant l'épreuve, vérifie que tu sais :
\begin{itemize}
  \item[$\square$] Donner la définition de $\ln$ comme primitive de $x\mapsto \frac{1}{x}$ nulle en $1$, et son ensemble de définition $]0\,;\,+\infty[$.
  \item[$\square$] Appliquer sans erreur $\ln(ab)$, $\ln\!\left(\frac{a}{b}\right)$, $\ln(a^n)$ et $\ln(\sqrt{a})$.
  \item[$\square$] Citer les valeurs remarquables : $\ln 1 = 0$, $\ln e = 1$.
  \item[$\square$] Donner les limites de $\ln x$ en $0^+$ et en $+\infty$, et interpréter les asymptotes de la courbe.
  \item[$\square$] Utiliser les croissances comparées : $\lim\limits_{x\to+\infty}\frac{\ln x}{x}=0$ et $\lim\limits_{x\to 0^+} x\ln x = 0$.
  \item[$\square$] Dériver $\ln u$ en $\frac{u'}{u}$ et étudier le sens de variation d'une fonction contenant $\ln$.
  \item[$\square$] Déterminer une primitive de la forme $\frac{u'}{u}$ sans oublier la valeur absolue ni la constante.
  \item[$\square$] Résoudre une équation ou une inéquation avec $\ln$ \textbf{en commençant toujours par les conditions d'existence}.
  \item[$\square$] Résoudre un système d'inconnues $\ln x$ et $\ln y$ par substitution.
  \item[$\square$] Tracer l'allure de la courbe de $\ln$ (passage par $(1\,;\,0)$ et $(e\,;\,1)$, concavité, croissance lente).
  \item[$\square$] Exploiter $\ln$ dans un problème concret : croissance exponentielle, désintégration radioactive, pH, échelle logarithmique.
  \item[$\square$] Rédiger proprement : justifier la positivité des quantités avant d'écrire un logarithme.
\end{itemize}
\end{aretenir}

\findecours

\end{document}
